\documentclass[11pt,letterpaper]{article}
\usepackage[T1]{fontenc}
\usepackage[utf8]{inputenc}
\usepackage{lmodern}
\usepackage[margin=1in,headheight=15pt]{geometry}
\usepackage{amsmath,amssymb,amsthm,mathtools,mathrsfs,bm}
\usepackage{microtype,booktabs,longtable,array,enumitem,tabularx}
\usepackage[dvipsnames]{xcolor}
\usepackage{fancyhdr,xurl,needspace}
\usepackage{listings}
\lstset{basicstyle=\ttfamily\small,breaklines=true,columns=fullflexible,frame=single}
\usepackage[colorlinks=true,linkcolor=MidnightBlue,citecolor=MidnightBlue,urlcolor=MidnightBlue,breaklinks=true]{hyperref}
\usepackage[nameinlink,noabbrev]{cleveref}
\hypersetup{pdftitle={Vector and Tensor Fields over the Surreal Numbers},pdfauthor={Merged research report for the Surreal project},pdfsubject={Admissible calculus in every dimension, three-dimensional vector calculus, and Minkowski fields}}
\pagestyle{fancy}\fancyhf{}
\fancyhead[L]{\small Surreal vector and tensor fields}
\fancyhead[R]{\small\thepage}
\renewcommand{\headrulewidth}{0.3pt}
\setlength{\parskip}{0.3em}
\setlength{\parindent}{1.2em}
\setlength{\emergencystretch}{3em}
\setlist{leftmargin=2em,itemsep=0.2em,topsep=0.35em}
\clubpenalty=10000\widowpenalty=10000\displaywidowpenalty=10000
\setcounter{tocdepth}{1}\numberwithin{equation}{section}\allowdisplaybreaks[2]
\newtheorem{theorem}{Theorem}[section]
\newtheorem{proposition}[theorem]{Proposition}
\newtheorem{lemma}[theorem]{Lemma}
\newtheorem{corollary}[theorem]{Corollary}
\theoremstyle{definition}
\newtheorem{definition}[theorem]{Definition}
\newtheorem{example}[theorem]{Example}
\newtheorem{convention}[theorem]{Convention}
\theoremstyle{remark}
\newtheorem{remark}[theorem]{Remark}
\newtheorem{warning}[theorem]{Warning}
\crefname{proposition}{proposition}{propositions}
\crefname{lemma}{lemma}{lemmas}
\crefname{corollary}{corollary}{corollaries}
\crefname{definition}{definition}{definitions}
\crefname{example}{example}{examples}
\crefname{warning}{warning}{warnings}
\crefname{convention}{convention}{conventions}
\crefname{part}{Part}{Parts}
\Crefname{part}{Part}{Parts}
\newenvironment{principle}{\begin{quote}\small}{\end{quote}}
\newcommand{\No}{\mathbf{No}}
\newcommand{\R}{\mathbb R}\newcommand{\C}{\mathbb C}
\newcommand{\Q}{\mathbb Q}\newcommand{\N}{\mathbb N}\newcommand{\Z}{\mathbb Z}
\newcommand{\A}{\mathscr A}\newcommand{\OO}{\mathcal O}\newcommand{\mm}{\mathfrak m}
\newcommand{\supp}{\operatorname{supp}}\newcommand{\st}{\operatorname{st}}
\newcommand{\id}{\operatorname{id}}\newcommand{\tr}{\operatorname{tr}}
\newcommand{\Ric}{\operatorname{Ric}}\newcommand{\Scal}{\operatorname{Scal}}
\newcommand{\grad}{\operatorname{grad}}\newcommand{\Div}{\operatorname{div}}
\newcommand{\curl}{\operatorname{curl}}
\newcommand{\vol}{\operatorname{vol}}\newcommand{\GL}{\operatorname{GL}}
\newcommand{\SO}{\operatorname{SO}}\newcommand{\SL}{\operatorname{SL}}
\newcommand{\Hom}{\operatorname{Hom}}\newcommand{\diag}{\operatorname{diag}}
\newcommand{\Span}{\operatorname{span}}\newcommand{\sgn}{\operatorname{sgn}}
\newcommand{\im}{\operatorname{im}}
\newcommand{\End}{\operatorname{End}}\newcommand{\Hess}{\operatorname{Hess}}
\newcommand{\Lie}{\mathcal L}\newcommand{\HH}{\mathcal H}
\newcommand{\ev}{\operatorname{Ev}}\newcommand{\jet}{\operatorname{J}}
\newcommand{\halo}[1]{#1^{\mathrm h}}
\newcommand{\abs}[1]{\lvert#1\rvert}\newcommand{\norm}[1]{\lVert#1\rVert}
\newcommand{\vH}{v_{\!H}}
\newcommand{\dd}{\mathrm d}
\newcommand{\e}{\mathrm e}
\newcommand{\eps}{\varepsilon}
\newcommand{\Gret}{G_{\mathrm{ret}}}
\newcommand{\M}{\mathbb M}
\newcommand{\inner}[2]{\eta(#1,#2)}
\newcommand{\BM}{\partial_{\mathrm{BM}}}
\newcommand{\repoSHA}{d22a5b35d5b3040e870c3cfd6d1c8f7259e094b0}
\newcommand{\repo}{\texttt{VladimirReshetnikov/Surreal}}
\newcolumntype{P}[1]{>{\raggedright\arraybackslash}p{#1}}
\begin{document}
\hypersetup{pageanchor=false}
\begin{titlepage}
\centering
\vspace*{0.1in}
{\large\scshape A merged research report\par}
\vspace{0.35in}
{\Huge\bfseries Vector and Tensor Fields\par over the Surreal Numbers\par}
\vspace{0.25in}
{\Large Admissible calculus in every dimension,\par three-dimensional vector calculus, and Minkowski fields\par}
\vspace{0.25in}
{\large September 22, 2026\par}
\vspace{0.25in}
\begin{minipage}{0.94\textwidth}
\begin{abstract}
\noindent
Surreal coefficients extend finite-dimensional tensor algebra without changing
its algebraic identities, but they do not determine a calculus on their own.
This report, merged from three manuscripts that build the same layered theory,
separates three settings: finite algebra over a set-sized real-closed surreal
subfield; smooth real manifolds whose fields are Hahn series of ordinary smooth
tensor fields; and genuine surreal-coordinate neighbourhoods reached by
support-controlled Taylor prolongation and scale charts.

Part~I works in every finite dimension and every nondegenerate signature:
prolongation with intrinsic derivatives, tensors, forms, admissible metrics,
Levi--Civita connections, curvature, Hodge operators, coefficientwise
integration, Stokes and Poincar\'e theorems, de Rham base change, Hodge
decomposition for a real metric and for a regular Riemannian Hahn metric,
standard-part reduction, near-real flows, exact inversion of positive-order
operator perturbations, and nonlinear lifting from an invertible real
linearization. Part~II specializes to three dimensions: dot and cross products,
leading-order rank, vector-calculus identities, integral theorems, a monopole
flux that forbids a global vector potential, an explicit Helmholtz
decomposition, infinitesimally forced stationary incompressible flow on the
three-torus, and Maxwell, stress and elasticity equations. Part~III treats
Minkowski space: Lorentz transformations of arbitrary scale, a split reduction
of the bounded Lorentz group with an exponential description of its kernel, a
causal-stability criterion, the Lorentzian Hodge operator, Maxwell and scalar
fields, stress--energy with exact doubling of the leading field scale, Cauchy
and retarded lifting, and a transfinite cubic-wave deformation.

Every infinite operation carries an explicit support hypothesis. The
obstructions, among them pointwise-invertible fields without admissible
inverses, infinite-frequency oscillations, infinite boosts and infinite times,
are proved rather than assumed away. Nothing here is a physics claim, a
regularization of singularities, or a transfer of all real analytic existence
theorems.
\end{abstract}
\end{minipage}
\vfill
{\small Built from three manuscripts prepared in relation to the \repo{} project,\par
all pinned to commit \texttt{d22a5b3}, and merged after placement commit \texttt{5fe7f8d}.\par
An AI-assisted research draft with proofs and finite symbolic checks; not refereed;\par
not a claim of a completed Lean formalization.\par}
\end{titlepage}
\hypersetup{pageanchor=true}
\pagenumbering{roman}
\begingroup\small\setlength{\parskip}{0pt}\tableofcontents\endgroup
\clearpage\pagenumbering{arabic}

\part{The theory in every finite dimension}\label{vtf:part:general}

\section{Scope and the central choices}\label{vtf:sec:scope}

An $n$-dimensional vector with surreal components is an element of $\No^n$,
understood class-theoretically. A vector \emph{field} requires more: its base
space, allowed functions, derivative, coordinate transformations, and any
integration operation must be specified. The two natural questions are
\[
 X:U\subseteq\R^n\longrightarrow\No^n,
 \qquad
 X:W\subseteq\No^n\longrightarrow\No^n.
\]
They are different questions. This report treats both, without identifying
the usual topology of the real domain with a topology inherited from all
surreals.

Throughout, $n\ge1$ is an ordinary finite integer. Indices run from $1$ to
$n$, and repeated upper/lower indices are summed. Tensor rank is finite as
well. The default metric is Euclidean; arbitrary nondegenerate signatures are
introduced explicitly. No step of Part~I depends on three dimensions unless it
is labeled as such. Parts~II and~III specialize to three Euclidean dimensions and
to four-dimensional Minkowski space.

\subsection{Three complementary models}

\begin{center}\small
\begin{tabular}{P{0.17\textwidth}P{0.28\textwidth}P{0.43\textwidth}}
\toprule
Model & Base and coefficients & What it supplies\\
\midrule
Finite algebra & $K^n$, with $K\subset\No$ a set-sized real closed field &
Tensors, metrics at a point, contractions, determinants, orthogonal algebra;
no infinite analysis is implicit.\\
Smooth Hahn geometry & An ordinary smooth $n$-manifold $M$, with locally
Hahn-supported smooth coefficients & Coordinate-independent differential
calculus, real-chain integration, global topology and scale-by-scale
geometric equations.\\
Intrinsic coordinate geometry & Open subsets of $K^n$, with rational or
support-certified Taylor--Hahn functions & Actual derivatives with respect
to surreal coordinates, finite halos, arbitrary-center scale charts and
algebraic integration.\\
\bottomrule
\end{tabular}
\end{center}

The main bridge is \cref{vtf:thm:prolongation}: a smooth Hahn field has a canonical
Taylor--Hahn realization on a nearstandard open subset of $K^n$. The second
model therefore is not merely symbolic bookkeeping. Its differential
identities become genuine intrinsic ordered-field identities on that subset.
It is nevertheless only one specified class of intrinsic functions.

\subsection{Three manuscripts, one report}\label{vtf:subsec:merge}

This report is built from three manuscripts, written independently on
September 22, 2026 against the same repository pin (\cref{vtf:sec:audit}).
Each constructs the same layered theory: finite algebra over a set-sized
real-closed surreal subfield, Hahn series of ordinary smooth tensor fields, and
Taylor prolongation or scale charts to genuine surreal-coordinate
neighbourhoods. Their numbers are local to this directory.
\begin{description}
\item[Source 01] \emph{Vector and Tensor Fields over the Surreal Numbers:
admissible calculus, differential geometry, and multiscale field equations},
an $n$-dimensional framework. It is the base: its text is Part~I, and its
statements are the most general ones (every finite dimension, every
nondegenerate signature, intrinsic derivatives of the prolongation). Its files
carry the prefix \texttt{01-vector-tensor-fields-n-}.
\item[Source 02] \emph{Surreal-Valued Fields on Minkowski Space: Lorentz
geometry, tensor calculus, and causal Hahn dynamics}. Its Lorentzian material is
Part~III. Its dimension-free remarks, the cohomology of a fixed coherent
complex, the algebraic Poincar\'e lemma, the infinite-frequency and
infinite-time obstructions and its distributional coefficients are placed in
Part~I. Its files carry the prefix \texttt{02-minkowski-}; its delivered source
audit is \path{02-minkowski-SOURCE_AUDIT.md}.
\item[Source 03] \emph{Three-Dimensional Surreal Vector and Tensor Fields:
support-controlled differential geometry, integral theorems, and multiscale
continuum equations}. Its three-dimensional material is Part~II. Its
dimension-free results (convolution extension of real multilinear maps, exact
inversion of positive-order perturbations, cohomology of a Hahn-extended
cochain complex, the Hodge theorem for a regular Riemannian Hahn metric,
quadratic lifting, the scale-connection calculus and the finite-partition
obstruction) are placed in Part~I. Its files carry the prefix
\texttt{03-three-dimensional-fields-}.
\end{description}
When two or three sources prove the same statement it is printed once, in
source 01's formulation when 01 has one, and the other sources are named
there; \cref{vtf:subsec:printedonce} lists these places. Genuinely different
proofs are kept, for example the tree expansion for quadratic lifting
(\cref{vtf:3d:thm:quadratic}). Before a result from the three-dimensional or
the Lorentzian source was attached to a general statement, its proof was
re-read for uses of dimension three, of the Euclidean metric, or of the
Lorentzian signature. Each such attachment is marked where it is made. Every
label carries the prefix \texttt{vtf:}; material from source 02 carries
\texttt{vtf:mink:} and material from source 03 \texttt{vtf:3d:}, wherever it is
placed.

Part~III is mathematics of fields on Minkowski space. It is not an
assessment in the sense of the collection's report \emph{Surreal Scalars in
Theoretical Physics} (\path{docs/physics/surreal-scalars-and-spacetime}), and it
respects that report's claim-status firewall: no statement in this report is a
physics claim. Words such as energy, radiation, observer, charge and fluid name
mathematical objects in the fixed conventions of \cref{vtf:app:lorentzconv}.
None of them asserts a measurable prediction, a physical regime, or the
regularization of a singular field.

\subsection{Relation to the repository and the literature}

At the pin of all three manuscripts (\cref{vtf:sec:audit}) the repository's
documentation distinguished, as it still does, finite algebra,
set-sized Hahn workspaces, common-domain coefficient rings, strong summability,
fine differentiation, and the Berarducci--Mantova scalar derivation
\cite{Repo,RepoAnalysis,RepoFound}. We retain those distinctions. In particular,
the one-complex-variable coherent calculus is a precedent for our evaluation
construction; the present treatment extends its \emph{method} to several real
variables, smooth coefficients, tensor sheaves, and geometric equations.

Conway normal forms and real closedness are background results
\cite{Gonshor,vdDE,Poonen93}; the support lemma is classical Hahn--Neumann
algebra \cite{Neumann,Poonen26}. Differential calculus over non-real topological
fields already has a broad theory \cite{BGN}. The tensor identities themselves
are classical differential geometry \cite{Lee,Petersen,Warner}. We give direct
arguments for the support-sensitive extensions rather than claiming that
ordinary completeness or compactness transfers to $K$.

Our contribution is a unified construction with explicit sufficient hypotheses
and proofs of its compatibility results. No claim of priority for every
result, resolution of a named open problem, or exhaustive novelty search is
made. Repository observations are pinned in \cref{vtf:sec:audit}; no code in that
repository is modified. Sources 02 and 03 make the same disclaimers
(\cref{vtf:subsec:nonclaims}).

\subsection{Neighbouring reports and the Lean development}\label{vtf:subsec:neighbours}

Several results of this report restate, in the real tensor setting, material the
collection already contains, and several connect to newer reports.
\begin{itemize}
\item \path{docs/surcomplex/contours-and-stokes} constructs the Hahn de Rham
complex with complex coefficients on any smooth manifold, with coefficientwise
Stokes (\path{contours:thm:hahnstokes}), a support-preserving Poincar\'e lemma
(\path{contours:prop:poincare}), the cohomology of the globally supported
Hahn complex (\path{contours:thm:cohomology}), polynomial-chain Stokes
(\path{contours:thm:polstokes}), Stokes for Hahn-deformed chains
(\path{contours:thm:deformedstokes}) and the finite-partition obstruction
(\path{contours:prop:mesh}). \Cref{vtf:thm:stokes,vtf:thm:poincare,vtf:3d:thm:cohomology},
the polynomial Stokes theorem of \cref{vtf:sec:integration} and
\cref{vtf:3d:prop:darboux} are real-coefficient versions of those statements and
are not claimed as new. Source 03 said so at its pin; source 01 cited that report
only through the documentation catalogue.
\item \path{docs/surcomplex/analysis} develops coherent Hahn sections over
complex domains; the evaluation construction of \cref{vtf:sec:prolong} extends
its method to several real variables and smooth coefficients.
\item \path{docs/foundations-and-computation/foundations} records the size
conventions and the discreteness of small subsets and small-index nets in the
full fine topology (\path{found:thm:discrete}), which \cref{vtf:sec:prolong}
uses.
\item \path{docs/physics/surreal-scalars-and-spacetime} is the collection's
assessment of surreal scalars in physics. Its bounded-frame valuation invariance
(\path{phys:prop:frame}) is the statement of \cref{vtf:sec:algebra} and of
\cref{vtf:mink:prop:frame}. Its standard-part reduction over a Hahn coefficient
algebra for a Lorentzian $g_0$ (\path{phys:thm:reduction}) is the
four-dimensional Lorentzian case of \cref{vtf:thm:reduction}, together with its
Einstein-tensor consequence. Its conclusion that scalar extension does not
regularize a singular field is shared by
\cref{vtf:subsec:warped,vtf:3d:subsec:monopole,vtf:mink:sec:limits}.
\item \path{docs/surquaternions/surquaternions} treats Hamilton quaternions over
real-closed fields, including the spin description of $\SO(3,F)$; the cross
product of Part~II is the vector part of the product of two pure quaternions.
Its no-oscillation theorem (\path{squat:thm:nooscillation}) concerns the
scalar Berarducci--Mantova derivation. The infinite-frequency obstruction
\cref{vtf:mink:thm:highfrequency} concerns coefficientwise derivatives. The two
statements are different.
\item \path{docs/surreal/euclidean-three-space}, placed in the same batch as this
report, develops the coordinate geometry of $\No^3$, the group $\SO(3,\No)$,
Rodrigues rotations and finite-angle trigonometry. The finite algebra of $K^3$ in
\cref{vtf:3d:sec:algebra} overlaps its first layer and leaves angles to it.
\item \path{docs/surcomplex/trigonometry} treats the unit circle and its angle
group; the duality rotations of \cref{vtf:mink:sec:maxwell} use only the
rational parametrization of the circle.
\item \path{docs/surcomplex/spectral-theory} and
\path{docs/surcomplex/infinite-dimensional-hahn-spectral-theory} treat spectral
theory over real-closed and Hahn fields; Part~II uses only the pointwise
symmetric spectral theorem.
\item \path{docs/surreal/hahn-valued-measures-and-probability} develops
coefficientwise Hahn-valued measures. The positivity statement of
\cref{vtf:sec:integration} concerns smooth common-support densities only.
\item \path{docs/surcomplex/differential-equations} studies differential
equations for the scalar derivation; its phase obstruction is not the
coefficientwise obstruction of \cref{vtf:sec:boundaries}.
\end{itemize}
The repository's Lean Hahn layer proves the full Neumann support lemma,
including finiteness across all word lengths
(\path{Surreal/HahnSeries/Neumann.lean}, \path{NeumannWords.lean}), positive-order
evaluation and composition in finitely many variables
(\path{MvEvaluation.lean}, \path{MvComposition.lean}), binomial roots near one
and standard part on the nonnegative-order subring. These are the scalar-level
inputs of \cref{vtf:sec:scalars} and of the support half of
\cref{vtf:thm:prolongation}. No field-level statement of this report has a Lean
counterpart, and \path{docs/FORMALIZATION.md} maps no \texttt{vtf:} label.

\subsection{Conventions fixed for the whole report}\label{vtf:subsec:conventions}

The three sources use overlapping symbols with different meanings. The choices
below are made once; the source text has been renamed to match, and
\cref{vtf:app:ledger} repeats them.
\begin{itemize}
\item \emph{Monomials and scales.} $t^\gamma=\omega^{-\gamma}$ as in
\eqref{vtf:eq:normal}; when $1\in\Gamma$, $t=t^1=\omega^{-1}$. Source 02 wrote
this monomial as $\eps^\gamma$. Here $\eps$ always denotes a chosen positive
infinitesimal $t^\eta$ with $\eta>0$, which is source 01's usage. Ordinary time
is $\tau$, never a monomial.
\item \emph{Derivatives.} $D_i$ in Part~I and $\partial_i,\partial_\mu$ in
Parts~II and~III are the coefficientwise derivatives of
\cref{vtf:prop:diffalgebra}. They annihilate every constant of $K$ and agree with
intrinsic derivatives on halos (\cref{vtf:thm:prolongation}). The intrinsic
derivative of a realized function is written $\partial_{y^i}$. The
Berarducci--Mantova scalar derivation is written $\BM$ and the character
derivations of \cref{vtf:sec:BM} are written $\mathscr D_\chi$. Neither is a
spatial derivative.
\item \emph{Alternating symbols and the Minkowski metric.} Alternating symbols are
$\epsilon_{ijk}$ and $\epsilon_{\mu\nu\rho\sigma}$; source 03 wrote
$\varepsilon_{ijk}$. In Part~III, $\eta_{\mu\nu}$ and $\eta(u,w)$ denote the
Minkowski metric, its quadratic form is $q_\eta(u)=\eta(u,u)$, and no exponent
is called $\eta$ there; elsewhere $\eta$ is an exponent. The letter $q$ alone is
the number of negative signs of a signature, as in \eqref{vtf:eq:star2}. The group
$\SO^+(1,3;K)$ is written $\mathcal G$.
\item \emph{Valuations.} $v$ is the valuation of a scalar. $v_V$ is the least
component valuation of a finite tensor in a stated frame (source 02's $\nu$).
$\vH$ is the least exponent of a nonzero coefficient function of a field
(source 02's $\operatorname{ord}_\Gamma$ and source 03's $v(F)$ for fields). The
number $\vH(F)$ bounds the valuations of the values $F(x)$ from below but need
not equal them. The same symbol $\vH$ denotes the least exponent of an element of
a Hahn vector space $E((t^\Gamma))$ (\cref{vtf:subsec:hahnspaces}).
\item \emph{Tolerances.} In \eqref{vtf:eq:fineder} and the proof of
\cref{vtf:thm:prolongation}, $\epsilon$ without indices is a positive element of
$K$; with indices, $\epsilon$ is always an alternating symbol.
\item \emph{Classes of fields and metrics.} A \emph{coherent} field on $U$
(source 02's term) is an element of $\A_\Gamma(U)$, with one well-ordered
support on all of $U$. An \emph{admissible} metric has determinant a unit of the
coefficient algebra (\cref{vtf:sec:metric}). A \emph{regular} metric is $g_0+h$
with $g_0$ real and nondegenerate of any signature and $h$ of positive support
\eqref{vtf:eq:regularmetric}. A \emph{regular Riemannian} metric (source 03's
``regular Hahn metric'') is a regular metric with $g_0$ positive definite.
``Positive support'' always means support in $\Gamma_{>0}$; ``positive'' for a
scalar, a form or a density refers to the order of $K$.
\item \emph{Laplacians and codifferentials.} $\Delta_g=\Div_g\grad_g$ and the
Cartesian $\Delta=\sum_i\partial_i^2$ have the positive-second-derivative sign;
$\Delta_H=d\delta+\delta d$ is the Hodge Laplacian, with $\Delta_H=-\Delta_g$ on
functions; $\Box=\partial_\tau^2-\Delta$. The symbol $\delta$ with no index is a
codifferential; $\delta_{ij}$ is Kronecker's symbol and $\delta_0$ a Dirac
distribution.
\item \emph{Renamings in the added material.} Source 02's exponent set $\Delta$
and exponent $\delta$ (in the retarded-perturbation, cubic-wave and
infinite-frequency statements) are $S_Q$, $S_r$ and $\theta$; its angular
current $M$ is $\mathcal M$ and its fluid expansion $\theta$ is $\vartheta$.
Source 03's Maxwell permittivity and permeability are $\kappa_E,\kappa_B$, its
Lam\'e constants are $\lambda_{\mathrm L},\mu_{\mathrm L}$, and its metric
density $J=\sqrt{\det g}$ is $\mu_g$, as in source 01.
\item \emph{Curvature.} The three sources' conventions agree:
$R(D_i,D_j)D_k=R^\ell{}_{kij}D_\ell$ and $\Ric_{kj}=R^i{}_{kij}$
\eqref{vtf:eq:curvature}; source 03's formulas were checked to be the same after
renaming indices. Christoffel symbols $\Gamma^k{}_{ij}$ always carry indices and
are not the exponent group $\Gamma$.
\item \emph{Actions.} Source 01's action density $\tfrac12\abs{d\phi}_g^2+V(\phi)$
and source 02's Lagrangian $\tfrac12\partial_\mu\phi\,\partial^\mu\phi-V(\phi)$
use opposite signs for the potential; see \cref{vtf:rem:potentialsign}.
\end{itemize}

\section{Scalars, size, and admissible infinite sums}\label{vtf:sec:scalars}

\subsection{A fixed real closed Hahn workspace}

Write a surreal normal form in increasing-exponent notation as
\begin{equation}\label{vtf:eq:normal}
 a=\sum_{\gamma\in S}a_\gamma t^\gamma,
 \qquad a_\gamma\in\R,\qquad t^\gamma:=\omega^{-\gamma},
\end{equation}
where $S\subset\No$ is a well-ordered \emph{set}. Here $\omega^{-\gamma}$
denotes Conway's monomial map, not an unannounced identification with the
Gonshor exponential $\exp(-\gamma\log\omega)$. The relevant law is
$t^{\gamma+\eta}=t^\gamma t^\eta$.

Fix a nonzero, set-sized, divisible ordered additive subgroup
$\Gamma\subset\No$, and put
\begin{equation}\label{vtf:eq:workspace}
 K=K_\Gamma=\R((t^\Gamma)).
\end{equation}
The normal-form embedding identifies this field with a subfield of $\No$.
The Hahn-field algebraic-closure theorem applied to $\C((t^\Gamma))=K[i]$
implies that $K$ is real closed \cite{Poonen93}. Thus positive square roots,
finite orthogonal algebra and signature are available in $K$.

\begin{proposition}[Localization of set-sized scalar data]\label[proposition]{vtf:prop:localize}
Every set-sized collection of surreal numbers is contained in a workspace
of the form \eqref{vtf:eq:workspace}.
\end{proposition}
\begin{proof}
Take the union of the normal-form supports of the given numbers. It is a set.
Its rational additive span in $\No$, enlarged by a nonzero exponent if
necessary, is a set-sized divisible ordered subgroup $\Gamma$. Every original
normal form is a Hahn series over $\Gamma$.
\end{proof}

This proposition localizes scalar data, not arbitrary regularity claims about
functions. It does not assert that every set-valued field has smooth Hahn
coefficients. Nor does it assert that a selected workspace is closed under
every additional scalar derivation.

\begin{remark}[Abstract exponent groups and the class level]
An abstract divisible ordered group may serve as $\Gamma$ once a strictly
order-preserving additive embedding into $\No$ is fixed; its image replaces
$\Gamma$ in \eqref{vtf:eq:normal} (sources 02 and 03). No part of the
construction needs a positive exponent to be cofinal in $\Gamma$, and larger
groups accommodate several Archimedean classes of scales. At the level of the
full class, $\No^n$ and its finite tensor carriers may be treated as classes.
Every finite algebraic computation takes place in a set-sized workspace and
inherits the identities proved below. A field on the proper class $\No^n$ would
need a separately specified class-function notion; it is not an object of the
real-domain theory, and no proper-class sum is formed (source 02).
\end{remark}

\subsection{Valuation, order, and standard part}

For $a\ne0$ define $v(a)=\min\supp(a)$, and set $v(0)=+\infty$. Then
\[
 v(ab)=v(a)+v(b),\qquad v(a+b)\ge\min(v(a),v(b)).
\]
The order of $K$ is the sign of the first nonzero real coefficient. Define
\[
 \OO_K=\{a:v(a)\ge0\},\qquad \mm_K=\{a:v(a)>0\}.
\]
These are, respectively, the finite elements and the infinitesimals relative
to $\R$. Every $a\in\OO_K$ has a unique decomposition
\[
 a=\st(a)+h,\qquad \st(a)\in\R,\quad h\in\mm_K,
\]
and coefficient extraction at exponent zero gives a ring homomorphism
$\st:\OO_K\to\R$ with kernel $\mm_K$. Infinite elements have no standard
part under this definition.

Coefficient extraction at exponent zero is not multiplicative on all of $K$:
$[t^0](\eps\eps^{-1})=1$, although both separate extractions vanish (source 02).

\begin{definition}[Strong summability]
A family $(a_j)_{j\in J}$ of Hahn series is strongly summable when the union
of its supports is well ordered and every exponent occurs in the support of
only finitely many $a_j$. Its sum is obtained by finite coefficientwise
addition. The same definition applies to coefficients in any real algebra,
or in a real vector space when only addition is used.
\end{definition}

\begin{lemma}[Classical support lemma]\label[lemma]{vtf:lem:neumann}
If $S,T$ are well-ordered subsets of an ordered abelian group, then $S+T$ is
well ordered and a fixed group element has only finitely many representations
$s+t$, with $s\in S,t\in T$. If $T\subset\Gamma_{>0}$ is well ordered, the
additive monoid
\[
 \langle T\rangle=\{0\}\cup\{\tau_1+\cdots+\tau_r:r\ge1,\ \tau_i\in T\}
\]
is well ordered. Every exponent has only finitely many word representations
from $T$, counting all lengths. Consequently, $S+\langle T\rangle$ is well
ordered and all the corresponding coefficient contributions are finite.
\end{lemma}

This is the support theorem imported from Hahn--Neumann theory
\cite{Neumann,Poonen26}. Its last assertion follows by first applying the
finite-pair statement to $S$ and $\langle T\rangle$, then the finite-word
statement to the second summand. We will name these supports whenever a
Taylor substitution, inverse, or nonlinear recursion is used.

\begin{corollary}[Geometric inversion]\label[corollary]{vtf:cor:geometric}
If $H$ is a scalar or finite square matrix of Hahn series and the union of
its entry supports is strictly positive, then
\[
 (I-H)^{-1}=\sum_{r\ge0}H^r.
\]
The sum is strong, including all matrix-index choices. The result is an exact
Hahn identity, not a claim about convergence of ordinary partial sums.
\end{corollary}
\begin{proof}
There are finitely many matrix entries. Apply \cref{vtf:lem:neumann} to their
support union. At any fixed exponent, only finitely many lengths and index
words contribute. The finite identity
$(I-H)\sum_{r=0}^{N}H^r=I-H^{N+1}$ therefore gives the asserted identity at
each exponent once all words contributing there have been included.
\end{proof}

For $h\in\mm_K$ the families
\begin{equation}\label{vtf:mink:eq:formalfunctions}
 (1+h)^{-1}=\sum_{m\ge0}(-h)^m,\qquad
 e^h=\sum_{m\ge0}\frac{h^m}{m!},\qquad
 \log(1+h)=\sum_{m\ge1}\frac{(-1)^{m+1}}{m}h^m
\end{equation}
are strong sums by the same argument, and finite matrices with infinitesimal
entries obey the same support rules (source 02). Formal identities between these
series, such as $\exp\log(1+h)=1+h$, hold exactly, because only finitely many
words contribute at each exponent.

\subsection{Strong sums are not ordinary limits}

For example, $\sum_{m\ge1}t^{m/(m+1)}$ is a Hahn series over $\Q$.
Its tails have valuations below $1$, so its finite partial sums do not
converge to that series in the valuation topology. Likewise, if
$0<\eta\ll\lambda$ in $\Gamma$, meaning $m\eta<\lambda$ for every
positive integer $m$, then $\sum_{m\ge0}t^{m\eta}$ is strongly summable but
its partial sums are not cofinal in valuation precision. A support argument,
not the phrase ``the terms become small,'' is the justification for our sums.
Sources 02 and 03 give the same examples.

\section{Finite-dimensional vectors and tensors}\label{vtf:sec:algebra}

\subsection{Tensor spaces and transformation laws}

Let $V=K^n$, with basis $e_i$ and dual basis $e^i$. A tensor of type $(r,s)$ is
\[
 T\in T^r_s(V):=V^{\otimes r}\otimes_K(V^*)^{\otimes s},
 \qquad \dim_K T^r_s(V)=n^{r+s}.
\]
Equivalently, it is a multilinear map on $r$ covector inputs and $s$ vector
inputs. Its components are
\[
 T=T^{i_1\cdots i_r}{}_{j_1\cdots j_s}
 e_{i_1}\otimes\cdots\otimes e_{i_r}
 \otimes e^{j_1}\otimes\cdots\otimes e^{j_s}.
\]
Contraction, tensor product, symmetrization, antisymmetrization and permutation
of slots use only finite sums and products. They are therefore valid over
$K$ and, for any given finite data, over the full surreal field.

If $e'_a=e_iP^i{}_a$, with $P\in\GL_n(K)$, then
\begin{equation}\label{vtf:eq:transform}
 T'^{a_1\cdots a_r}{}_{b_1\cdots b_s}
 =(P^{-1})^{a_1}{}_{i_1}\cdots(P^{-1})^{a_r}{}_{i_r}
 P^{j_1}{}_{b_1}\cdots P^{j_s}{}_{b_s}
 T^{i_1\cdots i_r}{}_{j_1\cdots j_s}.
\end{equation}
Coordinate transformations $y=\Phi(x)$ instead give vector components
$X_y^a=(\partial y^a/\partial x^i)X_x^i$ and covector components
$\alpha_{y,a}=(\partial x^i/\partial y^a)\alpha_{x,i}$. Distinguishing a
basis change from a coordinate map prevents an inverse-Jacobian ambiguity.

\subsection{Metrics and norms at a point}

A metric at a point is a symmetric nondegenerate bilinear form
$g:V\times V\to K$. Sylvester's law of inertia holds over a real closed
field: in a suitable basis,
\[
 g(u,u)=u_1^2+\cdots+u_p^2-u_{p+1}^2-\cdots-u_n^2,
 \qquad p+q=n.
\]
This follows by symmetric Gaussian elimination and rescaling nonzero diagonal
entries by positive square roots; the numbers of positive and negative terms
are intrinsic by the ordinary finite-dimensional intersection argument.

For positive-definite $g$, put $\norm{u}_g=\sqrt{g(u,u)}\in K_{\ge0}$.
For $v\ne0$, positivity of
$g(u-av,u-av)$ at $a=g(u,v)/g(v,v)$ proves
\[
 g(u,v)^2\le g(u,u)g(v,v).
\]
Consequently the triangle inequality and Gram--Schmidt orthonormalization
hold. These are ordered-field inequalities. The norm is $K$-valued, not a
real-valued Banach norm. No completeness theorem is being asserted.

The metric identifies vectors and covectors by
$X^\flat_i=g_{ij}X^j$ and $\alpha^{\sharp i}=g^{ij}\alpha_j$.
A surreal infinitesimal is not nilpotent: for nonzero $\eps$, every finite
power $\eps^m$ is nonzero. Dropping higher powers therefore changes an exact
tensor into a truncation and must be indicated.

\subsection{Which component scales are invariant?}

The coordinate minimum $v_V(X)=\min_i v(X^i)$ depends on a chosen lattice
$\OO_K^n\subset K^n$. If $P\in\GL_n(\OO_K)$, where the inverse also has
entries in $\OO_K$, both $P$ and $P^{-1}$ cannot lower this minimum. Thus
\[
 v_V(PX)=v_V(X).
\]
The same argument applies to the minimum component valuation of any fixed
finite tensor type. Under arbitrary $\GL_n(K)$ changes this is false:
$e'_1=t^\eta e_1$ changes the first component of $e_1$ to $t^{-\eta}$.
Invariant scale claims should refer either to a specified integral lattice or
to scalar contractions such as $g(X,X)$, $\tr(T^2)$ and curvature invariants.

When $X$ has entries in $\OO_K$ as well, coefficient extraction commutes with
such a change of frame: $\st(PX)=(\st P)\,\st(X)$ for $P\in\GL_n(\OO_K)$,
because $\st$ is a ring homomorphism on $\OO_K$. Source 02 proves both
statements for the bounded Lorentz group (\cref{vtf:mink:prop:frame}) and shows
with an infinite boost that neither extends to all of $\GL_n(K)$
(\cref{vtf:mink:ex:boostshadow}); the invariance is also the collection's
\path{phys:prop:frame}. Source 03 records the case of real frames and adds
leading-order information in three dimensions (\cref{vtf:3d:prop:leading}).

\section{The sheaf of smooth Hahn fields}\label{vtf:sec:sheaf}

\subsection{A common-domain algebra}

For an ordinary open $U\subseteq\R^n$ define
\begin{equation}\label{vtf:eq:A}
 \A_\Gamma(U)=C^\infty(U)((t^\Gamma))
 =\left\{\sum_{\gamma\in S}f_\gamma t^\gamma:
 S\subset\Gamma\text{ well ordered},\ f_\gamma\in C^\infty(U)\right\}.
\end{equation}
Support means the exponents with coefficient function not identically zero.
Evaluation at any real $x\in U$ is a well-defined map into $K$.
The representation as a function $U\to K$ is injective: equality at every
$x$ forces equality of each real coefficient function.

\begin{proposition}[Differential algebra]\label[proposition]{vtf:prop:diffalgebra}
The algebra \eqref{vtf:eq:A} is a commutative $K$-algebra. The operators
\[
 D_i\left(\sum_\gamma f_\gamma t^\gamma\right)
   =\sum_\gamma(\partial_i f_\gamma)t^\gamma
\]
are commuting $K$-linear derivations. They preserve strong summability and
never enlarge support. Pullback along an ordinary smooth map is a strongly
linear algebra homomorphism and satisfies the ordinary chain rule.
\end{proposition}
\begin{proof}
Products have support in $S+T$. By \cref{vtf:lem:neumann}, each coefficient is a
finite sum of products of real smooth functions. Differentiating that finite
sum gives Leibniz's rule. Mixed derivatives commute coefficientwise, and
restriction, differentiation and real pullback can only delete exponents.
All the other assertions reduce to finite coefficient operations and the
real smooth chain rule.
\end{proof}

This definition does \emph{not} mean that $U\to K$ is continuous when $U$
has its usual topology and $K$ its intrinsic order topology. It is the chosen
smooth coefficient structure. The intrinsic realization is supplied later.

\subsection{Hahn vector spaces and positive-order perturbations}\label{vtf:subsec:hahnspaces}

The same coefficientwise constructions apply to any real vector space. This
subsection, from source 03, records the two facts used repeatedly below.

For a real vector space $E$, define
\[
 E((t^\Gamma))=\left\{\sum_{\gamma\in S}t^\gamma e_\gamma:
 e_\gamma\in E,\ S\text{ well ordered}\right\}.
\]
This is a $K$-vector space. A real bilinear map $B:E\times F\to G$ extends by convolution:
\begin{equation}\label{vtf:3d:eq:bilinear}
 B\left(\sum t^\alpha e_\alpha,\sum t^\beta f_\beta\right)
   =\sum_\gamma t^\gamma\sum_{\alpha+\beta=\gamma}
      B(e_\alpha,f_\beta).
\end{equation}
The inner sum is finite by \cref{vtf:lem:neumann}.

For a real smooth manifold $M$ and real vector bundle $E\to M$, source 03 writes
\begin{align}
 \A_\Gamma(M)&=C^\infty(M)((t^\Gamma)),\\
 \mathfrak X_\Gamma(M)&=C^\infty(M,TM)((t^\Gamma)),\\
 \mathcal T^r_{s,\Gamma}(M)&=
 C^\infty(M,T^r_sM)((t^\Gamma)),\\
 \Omega_\Gamma^k(M)&=\Omega^k(M)((t^\Gamma)).
\end{align}
These are global-support spaces; on a compact $M$, $\Omega^k_\Gamma(M)$ is the
space $\Omega^k_H(M)$ of \cref{vtf:sec:cohomology}. Evaluation at a real point is
coefficientwise and gives a tensor over $K$. Evaluation on all real points is injective: a coefficient field vanishing everywhere is zero.

\begin{theorem}[Algebraic closure of the field calculus; source 03]\label{vtf:3d:thm:closure}
Addition, tensor product, contraction, permutation, symmetric and alternating products, and every fixed real-linear differential operator preserve these spaces. A real-linear operator does not enlarge support. A bilinear operation has support contained in the sum of the input supports. All finite polynomial identities between such operations extend from real fields to Hahn-smooth fields.
\end{theorem}
\begin{proof}
Linear operations act on each coefficient. Bilinear operations use \eqref{vtf:3d:eq:bilinear}. For an identity, fix an exponent. Its coefficient is a finite sum of the corresponding real multilinear identities, after polarization when necessary. Thus it is zero. This proof is algebraic; continuity in a real function-space norm is not needed for the extension itself.
\end{proof}

\begin{theorem}[Support-controlled inverse; source 03]\label{vtf:3d:thm:inverse}
Let $E,F$ be real vector spaces and $L_0:E\to F$ a real-linear isomorphism. Let
\[
 P=\sum_{\sigma\in S}t^\sigma P_\sigma,
 \qquad S\subseteq\Gamma_{>0}\text{ well ordered},
 \qquad P_\sigma:E\to F\text{ real linear}.
\]
Then $L=L_0+P:E((t^\Gamma))\to F((t^\Gamma))$ is an isomorphism. Its inverse is the strongly defined operator expression
\begin{equation}\label{vtf:3d:eq:inverse}
 L^{-1}=\sum_{n\geq0}(-L_0^{-1}P)^nL_0^{-1}.
\end{equation}
For data $f$ with support $T$, the solution has support in
$T+\langle S\rangle$. If $P\ne0$ and $s=\min S$, then
\[
 \vH(L^{-1}f-L_0^{-1}f)\geq \vH(f)+s.
\]
The bound allows the left side to be $+\infty$.
\end{theorem}
\begin{proof}
Expanding the $n$th term produces compositions of $n$ real operators with exponents in words from $S$. By \cref{vtf:lem:neumann}, only finitely many such terms contribute to each output exponent, for all $n$ together. Thus \eqref{vtf:3d:eq:inverse} defines an actual map on Hahn vectors, and finite geometric-series cancellation proves both inverse identities coefficientwise. Alternatively, uniqueness follows by comparing the least exponent of a nonzero $u$: the leading coefficient of $L_0u$ is nonzero, and every coefficient of $Pu$ has strictly larger exponent. The support and valuation bounds follow directly from the displayed expansion.
\end{proof}

There is no smallness condition in an ordinary operator norm. The coefficients $P_\sigma$ may be differential operators unbounded in some Hilbert norm. What is required is that each one be defined on all of the specified real space $E$ with values in $F$, and that $L_0^{-1}$ exist on all of $F$. Domain and boundary conditions cannot be discarded.

\begin{remark}[More than a scalar extension]\label{vtf:3d:rem:nottensor}
Usually $E\otimes_\R K$ is a strict subspace of $E((t^\Gamma))$: its elements
involve only finitely many real coefficient vectors. For example
$F(x)=\sum_{m\ge0}t^mx^m\in C^\infty(\R)((t^\Q))$ has infinitely many linearly
independent coefficient functions, so it is not a finite sum
$\sum_{j=1}^ka_jf_j(x)$ with $a_j\in K$. This completion in the scale direction
is essential; it is not a completion in the spatial variable (sources 02 and
03).
\end{remark}

\begin{warning}[The field algebra is not a field; source 02]
$C^\infty(U)$ has zero divisors, for example two nonzero bump functions with
disjoint supports, and so does $\A_\Gamma(U)$. The least exponent
$\vH(f)=\min\{\gamma:f_\gamma\ne0\}$ satisfies $\vH(fg)\ge\vH(f)+\vH(g)$, not
equality in general. Statements about the valuation of numbers in $K$ must not be
copied verbatim to coefficient functions. Since $D_i$ acts on coefficients,
\begin{equation}\label{vtf:mink:eq:derivativeorder}
 \vH(D_if)\ge\vH(f).
\end{equation}
\end{warning}

Membership in the subring $\A_{\ge0}(U)$ of fields with nonnegative support is
equivalent to finiteness of the value at every real point: if a
negative-exponent coefficient function is nonzero at some point, the value there
has a negative leading exponent. Likewise, infinitesimal values at every real
point characterize the positive-support ideal $\A_{>0}(U)$. No uniform bound on
the sizes of the coefficient functions or their derivatives accompanies this
(source 02). For instance, $\sum_{m\ge0}t^m\sin(mx^1)$ belongs to $\A_\Q(\R^n)$
and has derivatives of every order; the growing real factors $m^k$ in its
$k$th derivative are harmless to the Hahn support. This does not assert
convergence after replacing $t$ by a positive real number, and it does not put
an infinite frequency $\sin(t^{-1}x^1)$ in the algebra (\cref{vtf:sec:boundaries}).

\subsection{Why local support is necessary}

The presheaf $U\mapsto\A_\Gamma(U)$ need not be a sheaf for arbitrary open
covers. On disjoint intervals $U_m$ one can prescribe the constant section
$t^{1/m}$. Each local section is admissible, but the union of the exponents
$\{1/m:m\ge1\}$ is not well ordered. The glued function has no single
representation \eqref{vtf:eq:A} on $\bigcup_mU_m$. Sources 02 and 03 give the
same obstruction with the sections $t^{-m}$, whose exponents are unbounded
below.

Define $\A^{\mathrm{loc}}_\Gamma$ to consist of functions that admit such a
representation near every real base point. This is a sheaf. Uniqueness of
coefficient extraction makes gluing unambiguous. All local algebra and
calculus statements apply to it. On a smooth manifold, use real coordinate
charts and the pullbacks of \cref{vtf:prop:diffalgebra}.

On a compact real manifold there is no global support defect. A finite
subcover of local presentations gives a finite union of well-ordered
supports, hence one well-ordered set. Coefficient functions, or coefficient
tensors, agree on overlaps and glue to global smooth objects. The analogous
statement holds for a section restricted to a neighborhood of a compact set
after a finite localization.

Global theorems below use explicitly specified common-domain sections. Global
support is not inferred from pointwise well ordering or from local
admissibility on a noncompact domain (sources 02 and 03).

\paragraph{A pointwise family that is not Hahn-smooth (source 03).} Suppose
$\Gamma=\R$ and put $F(x)=t^x$ for real $0<x<1$. Every value lies in $K$, but
$F\notin C^\infty((0,1))((t^\R))$: the coefficient of $t^\gamma$ would be the
indicator of $\{\gamma\}$, which is not smooth, and the supports have no common
well ordering. ``All components are surreal'' is weaker than the condition the
calculus needs.

\paragraph{Field valuation and point values (source 03).} A leading coefficient
field may vanish at some points. Thus $\vH(F)$ controls every value $F(x)$ from
below but need not equal its valuation. Estimates below use $\vH$; pointwise
statements evaluate leading coefficients explicitly.

\subsection{Not every nonzero field is a unit}\label{vtf:subsec:units}

A sufficient local unit condition is
\begin{equation}\label{vtf:eq:unit}
 f=t^\lambda(a+h),\qquad a\in C^\infty(U),\quad a(x)\ne0\text{ on }U,
 \quad \supp(h)\subset\Gamma_{>0}.
\end{equation}
Then
\[
 f^{-1}=t^{-\lambda}a^{-1}
       \sum_{m\ge0}(-a^{-1}h)^m\in\A_\Gamma(U).
\]
If $a>0$, divisibility of $\Gamma$ also gives
\[
 \sqrt f=t^{\lambda/2}\sqrt a
   \sum_{m\ge0}\binom{1/2}{m}(h/a)^m.
\]
All these supports are translates of positive generated monoids.

\begin{example}[A pointwise inverse outside the smooth Hahn algebra; sources 01 and 03]
Fix $\eps=t^\eta$, $\eta>0$, and set $f(x)=x^2+\eps^2$ on a real interval
around zero. Every $f(x)$ is positive and invertible in $K$. Nevertheless,
$f^{-1}$ is not in $\A^{\mathrm{loc}}_\Gamma$ at zero.

Indeed, at $x=0$ its coefficient at exponent $-2\eta$ is $1$. At every
nonzero real $x$ the exact expansion is
\[
 \frac1{x^2+\eps^2}=\sum_{m\ge0}(-1)^m x^{-2m-2}\eps^{2m},
\]
so that same coefficient is zero. It could not be a continuous real
coefficient function on a neighborhood of zero. This example will distinguish
pointwise nondegeneracy from an admissible metric inverse.
\end{example}

\section{Intrinsic derivatives and Taylor--Hahn prolongation}\label{vtf:sec:prolong}

\subsection{The topology that is actually used}

Give $K$ its \emph{intrinsic} order topology, whose radii lie in $K_{>0}$.
On $K^n$ use the equivalent coordinate maximum norm
$\norm{x}_\infty=\max_i|x^i|$. A derivative of $F:W\subseteq K^n\to K^m$
at $y$ is a $K$-linear map $L$ satisfying
\begin{equation}\label{vtf:eq:fineder}
 \forall\epsilon\in K_{>0}\ \exists\rho\in K_{>0}\quad
 0<\norm{k}_\infty<\rho\ \Longrightarrow\
 \norm{F(y+k)-F(y)-Lk}_\infty<\epsilon\norm{k}_\infty,
\end{equation}
with a ball about $y$ contained in $W$. The linear map is unique: test the
difference of two candidates on small multiples of each coordinate vector.
Polynomial and rational functions are differentiable on their denominator
nonzero domains by finite algebraic remainder identities.

The topology inherited by a set-sized $K\subset\No$ from the \emph{full}
surreal fine topology is different: it is discrete. To see the underlying
obstruction, let $S$ be any set of surreals and fix $x$. There is a positive
surreal smaller than every nonzero distance $|s-x|$, $s\in S$, by the
small-cut property. A sufficiently small ball isolates $x$ from $S\setminus
\{x\}$. Every set-sized subset is thus closed and discrete. In particular,
ordinary real-parameter curves into full $\No$ that are continuous in that
fine topology must be constant. These size-sensitive facts are also emphasized
in the repository \cite{RepoFound}. We never give $K$ that inherited topology
when using \eqref{vtf:eq:fineder}.

Sources 02 and 03 add three observations about the intrinsic topology. First, a
non-Archimedean ordered field $K\supset\R$ has $\mm_K$ as an open and closed
additive subgroup, and translates and rescalings of it separate points; so $K$
and $K^n$ are totally disconnected, and an intrinsically continuous map from a
connected real interval into $K^n$ is constant. Second, the map $\R\to K$
sending a real number to the same constant is not continuous from the ordinary
topology to the order topology: an infinitesimal ball about $0$ meets $\R$ only
in $0$. The derivative \eqref{vtf:eq:fineder} is therefore used on open subsets
of $K^n$, never along ordinary real curves, and the smooth structure of
$\A_\Gamma$ is specified coefficientwise. Third, a fixed support presentation
can be given coefficientwise smooth convergence; no argument below equates
strong summation, coefficientwise convergence and valuation-topological
convergence.

\begin{proposition}[Finite-partition obstruction; source 03]\label{vtf:3d:prop:darboux}
Let $K\supset\R$ be non-Archimedean. Ordinary finite-partition Darboux integration, with error required to be smaller than every positive element of $K$, fails for $f(x)=x$ on $[0,1]_K$.
\end{proposition}
\begin{proof}
For a partition with $N$ subintervals and lengths $\ell_j>0$, the difference between upper and lower sums is $\sum_j\ell_j^2$. Since $\sum_j\ell_j=1$, finite Cauchy--Schwarz gives $\sum_j\ell_j^2\geq1/N$. Every positive infinitesimal is smaller than $1/N$ for every ordinary finite $N$. Therefore no such partition achieves infinitesimal error. The argument permits partition points in $K$, not just in $\R$.
\end{proof}

The same obstruction, with an explicit bound, is \path{contours:prop:mesh}.
Integrals below are either coefficientwise over real chains or algebraic on
polynomial chains.

\subsection{Prolongation of arbitrary smooth coefficients}

For a real open $U\subset\R^n$ put
\[
 \halo U=\{a+h:a\in U,\ h\in\mm_K^n\}\subset K^n.
\]
It is intrinsically open. The decomposition $a=\st(y)$, $h=y-a$ is unique.
For $F=\sum_\gamma f_\gamma t^\gamma\in\A_\Gamma(U)$ define
\begin{equation}\label{vtf:eq:eval}
 (\ev F)(a+h)=
 \sum_{\gamma\in\supp(F)}\ \sum_{\alpha\in\N^n}
 \frac{\partial^\alpha f_\gamma(a)}{\alpha!}\,t^\gamma h^\alpha.
\end{equation}
This is a Taylor-\emph{jet} construction. It does not assert convergence of a
real Taylor series to a real smooth function on an ordinary neighborhood.

\begin{theorem}[Multivariable smooth prolongation]\label{vtf:thm:prolongation}
Formula \eqref{vtf:eq:eval} is strongly summable and defines an injective
$K$-algebra homomorphism into functions on $\halo U$. It extends evaluation
on real points. Its values are intrinsically $C^\infty$ in the sense of
iterated \eqref{vtf:eq:fineder}, and
\begin{equation}\label{vtf:eq:derbridge}
 \partial_{y^i}(\ev F)=\ev(D_iF).
\end{equation}
For every ordinary smooth map $\Phi:U\to V$,
\begin{equation}\label{vtf:eq:natural}
 \ev(\Phi^*F)=(\ev F)\circ(\ev\Phi),
\end{equation}
where $\ev\Phi$ is componentwise prolongation. These statements localize to
$\A^{\mathrm{loc}}_\Gamma$.
\end{theorem}
\begin{proof}
Let $S=\supp(F)$, and let $T$ be the union of the nonzero coordinate supports
of $h$. It is a finite union of well-ordered positive sets. Every summand of
\eqref{vtf:eq:eval} has support in $S+\langle T\rangle$. A fixed exponent has
only finitely many choices of $\gamma$, word length, monomials from the
coordinates of $h$, and multiindex: this is \cref{vtf:lem:neumann}, with the
finite coordinate labels attached to each word. Thus the displayed double
sum is strong. Evaluation at $h=0$ gives the asserted restriction and proves
injectivity by coefficient extraction on all real points.

For each $a$, taking the full Taylor jet
$\jet_a:C^\infty(U)\to\R[[Z_1,\ldots,Z_n]]$ is an algebra homomorphism.
Leibniz's formula proves this degree by degree. Substitution of positive-order
$h$ into these jets is a homomorphism by the same support lemma. This proves
additivity, multiplicativity and compatibility with $K$ constants.

We give the derivative estimate, since it is the important additional step.
Fix $y=a+h$. For $k\in\mm_K^n$, formal Taylor translation and strong
regrouping yield
\begin{equation}\label{vtf:eq:remainder}
 \ev F(y+k)=\ev F(y)+\sum_i\ev(D_iF)(y)k^i
             +\sum_{i,j}k^ik^j R_{ij}(h,k).
\end{equation}
The remainders are evaluations of formal real series with Hahn supports in
$S+\langle T_h\cup T_k\rangle$. If $F\ne0$ and $\beta=\min S$, every
$R_{ij}$ has valuation at least $\beta$, independently of the small $k$.
This follows by removing two of the $k$ factors from terms of $k$-degree at
least two; the remaining factors all have nonnegative valuation. In
particular,
\[
 v\left(\ev F(y+k)-\ev F(y)-\sum_i\ev(D_iF)(y)k^i\right)
 \ge\beta+2v(\norm{k}_\infty).
\]
Given $\epsilon>0$, choose $\rho=t^\theta$ with $\theta>0$ and
$\beta+\theta>v(\epsilon)$. For $0<\norm{k}_\infty<\rho$, the
remainder divided by $\norm{k}_\infty$ has valuation strictly greater than
$v(\epsilon)$ and hence absolute value less than $\epsilon$. The ball stays
in the same halo. This proves \eqref{vtf:eq:derbridge}. Applying the same argument
to every $D^\alpha F$ proves existence of all iterated derivatives; its
first-order difference estimate proves their continuity.

Finally, the ordinary smooth chain rule in every finite jet degree gives the
formal identity
\[
 \jet_a(f\circ\Phi)=
 (\jet_{\Phi(a)}f)\circ\bigl(\jet_a\Phi-\Phi(a)\bigr).
\]
The inner series have zero constant terms and evaluate to infinitesimals.
The common positive-support certificate justifies every regrouping, proving
\eqref{vtf:eq:natural}. Agreement of local representatives proves localization.
\end{proof}

\begin{remark}[Three proofs of the prolongation]
Sources 02 and 03 prove the algebraic half of \cref{vtf:thm:prolongation} by the
same support argument, with the same formal Taylor homomorphism: strong
summability on $\supp F+\langle T\rangle$, the homomorphism property,
compatibility with formal differentiation, and the chain rule. Source 03 also
proves the intrinsic derivative at each halo point, with a slightly different
bound: for any fixed $\rho\in\Gamma_{>0}$ the remainder is at most
$t^{\beta-\rho}r^2$, where $r=\max_i\abs{k^i}$ is infinitesimal. Source 01's
statement is the one printed, since it includes all iterated intrinsic
derivatives and the naturality \eqref{vtf:eq:natural}. Source 02 phrases the
prolongation on the set $\{z\in\OO_K^n:\st z\in U\}$ of finite points over $U$,
which is $\halo U$.
\end{remark}

\begin{remark}[Smooth is not analytic]
If $f(x)=e^{-1/x^2}$ for $x\ne0$ and $f(0)=0$, then $\ev f$ vanishes on
the whole halo $\mm_K$ of zero, because its jet there is zero. It still agrees
with $f$ at every nonzero real point. This is a well-defined smooth
prolongation, not an identity theorem for analytic functions.
\end{remark}

Source 03 sharpens the point with the one-sided function $f(x)=\e^{-1/x^2}$ for
$x>0$ and $f(x)=0$ for $x\le0$. Its prolongation vanishes at every positive
infinitesimal, whereas the surreal exponential of $-1/x^2$ does not. There is no
contradiction: $-1/x^2$ is not smooth at $0$, and the halo of $(0,\infty)$
contains no point with standard part $0$. Taylor prolongation is therefore not a
canonical extension of all smooth expressions to all surreal arguments, and it
does not address the extension problems of Costin and Ehrlich \cite{CE}. For
real-analytic coefficients the prolongation agrees with the admissible
infinitesimal Taylor expansion. At an unlimited argument neither a real base
point nor this prescription exists (source 02).

\subsection{Nonuniqueness outside the chosen class}\label{vtf:subsec:nonunique}

On $\halo\R=\OO_K$, the functions $y\mapsto y$ and $y\mapsto\st(y)$
agree at every real point. The first has derivative $1$; the second is locally
constant on each open halo and has derivative $0$. Thus restriction to real
points plus intrinsic smoothness does not determine an extension. Theorem
\ref{vtf:thm:prolongation} chooses the extension by all real Taylor jets.
The kernel of the derivative is constants on a connected \emph{real base}
inside our Hahn class; it is much larger among all intrinsic smooth functions.

\subsection{Surreal-valued coordinate changes}

Let $\Phi_0:U\to V$ be real smooth and let $H\in\A_\Gamma(U)^m$ have
strictly positive support. Write $\Phi=\Phi_0+H$. For
$F=\sum_\sigma f_\sigma t^\sigma\in\A_\Gamma(V)$ define
\begin{equation}\label{vtf:eq:nonrealpull}
 \Phi^*F=\sum_{\sigma,\alpha}
 \frac{(\partial^\alpha f_\sigma)\circ\Phi_0}{\alpha!}
 t^\sigma H^\alpha.
\end{equation}
The support is contained in $\supp(F)+\langle\supp(H)\rangle$.
The proof of \cref{vtf:thm:prolongation} applies verbatim to show that pullback is
an algebra map, obeys the chain rule, and realizes actual composition on halos.

\begin{proposition}[Lifting coordinate diffeomorphisms]\label[proposition]{vtf:prop:diffeo}
If $\Phi_0$ is a real diffeomorphism and $H$ has positive support, then
$\ev\Phi:\halo U\to\halo V$ is a bijection with an inverse of the same
Taylor--Hahn type. The correction to $\Phi_0^{-1}$ has support in
$\langle\supp(H)\rangle\setminus\{0\}$.
\end{proposition}
\begin{proof}
Write $\Psi=\Phi_0^{-1}+Q$. At exponent $\gamma>0$, the equation
$\Phi\circ\Psi=\id$ is a linear equation for $Q_\gamma$ with coefficient
matrix $D\Phi_0\circ\Phi_0^{-1}$. This matrix has a smooth real inverse.
All remaining terms depend on strictly smaller positive exponents. The
support lemma makes their number finite at each exponent. Transfinite
recursion over the generated monoid constructs a unique $Q$. Construct a
left inverse in the same way. A left inverse and a right inverse coincide by
associativity of the already certified compositions, proving both the formal
and realized statements.
\end{proof}

\section{Genuinely surreal-coordinate neighborhoods}\label{vtf:sec:charts}

\subsection{Arbitrary centers and scales}

The finite halo is not all of $K^n$. To work near arbitrary $a\in K^n$, choose
nonzero scales $r_1,\ldots,r_n\in K$ and coordinates
\[
 z^i=\frac{x^i-a^i}{r_i},\qquad z\in\mm_K^n.
\]
One may evaluate every element of
\begin{equation}\label{vtf:eq:formalhahn}
 \R[[Z_1,\ldots,Z_n]]((t^\Gamma))
\end{equation}
at such a $z$: its outer Hahn support is fixed, while the coefficient at
each exponent is an arbitrary real formal series. The support proof and the
remainder proof of \cref{vtf:thm:prolongation} still apply. In these coordinates
\[
 \partial_{x^i}=r_i^{-1}\partial_{z^i}.
\]
These scale factors matter: an infinitesimal spatial scale can turn a modest
coordinate gradient into an infinite physical-coordinate gradient.

Polynomials and rational functions over $K$ are available independently of
this presentation. Near a point where a rational denominator is nonzero,
shrink a $K$-radius so that its relative change has positive valuation, then
use geometric inversion. Such rational germs are covered by sufficiently
small scale charts even when they fail to be smooth Hahn fields over a large
real base, as in \cref{vtf:subsec:units}.

\paragraph{Linear scale charts (source 03).} More generally, for $c\in K^n$ and
$L\in\GL_n(K)$ put $x=c+Ly$ with $y\in\halo U$. This accommodates an infinite
centre, an infinitesimal cell, and different scales along non-coordinate
directions. For $F(x)=\ev f(L^{-1}(x-c))$ and a vector field $V(x)=L\,\ev v(y)$
the chain rule gives
\[
 \grad_xF=L^{-T}\grad_y\ev f,\qquad
 D_xV=L\,(D_y\ev v)\,L^{-1},\qquad
 \Div_xV=\Div_y\ev v,
\]
with Euclidean gradients and divergences in both charts. The pulled-back
Euclidean metric is $L^TL$ and the volume factor is $\abs{\det L}$. Source 03
states these formulas for $n=3$; the computation uses no property of dimension
three. Gluing two such charts requires a transition map admissible in both chart
calculi, exactly as in \cref{vtf:subsec:exclusion}. The three-dimensional cross
product transforms with the extra factor in \eqref{vtf:3d:eq:crossL}. Source 02
uses the isotropic case $x=a+t^bX$ in Minkowski space
(\cref{vtf:mink:subsec:inner}).

\subsection{An important exclusion}\label{vtf:subsec:exclusion}

An arbitrary element of $K[[Z]]$ is not promised to evaluate on the whole
infinitesimal halo. For example,
\[
 \sum_{m\ge0}t^{-m^2\eta}Z^m
 \quad\text{at}\quad Z=t^\eta
\]
has proposed exponents $(-m^2+m)\eta$, not a well-ordered support.
The uniform outer support in \eqref{vtf:eq:formalhahn} is a sufficient hypothesis
that excludes this example. Other, larger analytic classes are possible,
but they need their own support certificates or domains. A formal inverse
function theorem alone is not a convergence theorem on a specified domain.

An intrinsic $n$-dimensional atlas may now be built from such open subsets of
$K^n$, requiring every transition and its inverse to belong to an explicitly
admitted calculus. Real diffeomorphism prolongations, their positive-support
perturbations, affine scale changes, and rational local coordinate changes
provide concrete examples. Tensor and connection formulas below apply in
any such atlas once its chain rule and closure hypotheses have been checked.
Global integration and cohomology statements will continue to name the real
base or algebraic chains to which they apply.

\section{Vector fields, tensor fields, and exterior calculus}\label{vtf:sec:fields}

\subsection{Modules of geometric fields}

On a real coordinate domain write $\A=\A^{\mathrm{loc}}_\Gamma$ and define
\[
 \mathfrak X_\A(U)=\bigoplus_{i=1}^n\A(U)D_i,
 \qquad
 \Omega^1_\A(U)=\bigoplus_{i=1}^n\A(U)\dd x^i.
\]
These glue under the real Jacobian laws. Tensor fields are sections of the
corresponding finite tensor powers. For an ordinary bundle $E\to M$, its
Hahn extension has locally the same construction; globally a section has
locally well-ordered sums of real sections. We do not identify these
geometric fields with every possible algebraic derivation of a huge function
ring without a locality or strong-linearity condition.

A vector field acts by $X(F)=X^iD_iF$. Its Lie bracket is
\begin{equation}\label{vtf:eq:bracket}
 [X,Y]^i=X^jD_jY^i-Y^jD_jX^i.
\end{equation}
Because it is the commutator of derivations, the bracket is antisymmetric and
satisfies Jacobi. Coordinate invariance follows from the chain rule and
cancellation of the symmetric second derivatives of a coordinate change.
It is $K$-bilinear, but not $\A$-bilinear:
$[X,fY]=f[X,Y]+X(f)Y$.

\subsection{Forms, exterior derivative, and Cartan calculus}

Define
\[
 \Omega^p_\A=\bigwedge^p_\A\Omega^1_\A,
 \qquad \operatorname{rank}_\A\Omega^p_\A=\binom np.
\]
For increasing multiindices $I$, set
\[
 d\left(\sum_I\alpha_I\dd x^I\right)
   =\sum_{I,j}D_j\alpha_I\dd x^j\wedge\dd x^I.
\]
Alternation and commutation of the $D_i$ prove $d^2=0$. For homogeneous
$\alpha$ of degree $p$,
\[
 d(\alpha\wedge\beta)=d\alpha\wedge\beta+(-1)^p\alpha\wedge d\beta.
\]
Insertion is $(\iota_X\alpha)(X_2,\ldots,X_p)=\alpha(X,X_2,\ldots,X_p)$.
Define the Lie derivative on forms by Cartan's formula
\begin{equation}\label{vtf:eq:cartan}
 \Lie_X=d\iota_X+\iota_Xd.
\end{equation}
It follows algebraically that $[\Lie_X,d]=0$,
$[\Lie_X,\iota_Y]=\iota_{[X,Y]}$ and
$[\Lie_X,\Lie_Y]=\Lie_{[X,Y]}$. To verify these identities directly, note
that both sides are graded derivations and agree on functions and coordinate
one-forms, which generate the exterior algebra. No flow theorem is needed
for this definition.

For a tensor $T$ of type $(r,s)$ the component formula is
\begin{align}\label{vtf:eq:lietensor}
 (\Lie_XT)^{i_1\cdots i_r}{}_{j_1\cdots j_s}
 &=X^kD_kT^{i_1\cdots i_r}{}_{j_1\cdots j_s}\notag\\
 &\quad-\sum_{a=1}^r(D_kX^{i_a})
 T^{i_1\cdots k\cdots i_r}{}_{j_1\cdots j_s}
 +\sum_{b=1}^s(D_{j_b}X^k)
 T^{i_1\cdots i_r}{}_{j_1\cdots k\cdots j_s}.
\end{align}
Its signs express the contravariant and covariant transformation laws.
Theorem \ref{vtf:thm:prolongation} transports all these identities to halo
realizations, and the intrinsic atlas construction transports them to any
admitted coordinate domain.

\section{Admissible metrics and their Levi--Civita connections}\label{vtf:sec:metric}

\subsection{Pointwise versus differential nondegeneracy}

A symmetric matrix $g_{ij}\in\A(U)$ may be nondegenerate at every real
point without its inverse belonging to $\A(U)$. For a \emph{metric in this
calculus}, we require $\det g$ to be a unit of the chosen local coefficient
algebra. When a volume form is used, its square root is required in that
algebra as well. This is an admissibility condition, not an extra condition
on finite-dimensional metric algebra at one point.

A particularly useful class is
\begin{equation}\label{vtf:eq:regularmetric}
 g=g_0+h,\qquad g_0\text{ a real smooth nondegenerate metric},
 \quad\supp(h)\subset\Gamma_{>0}.
\end{equation}
Finite matrix inversion gives the exact formula
\begin{equation}\label{vtf:eq:metricinverse}
 g^{-1}=\sum_{m\ge0}(-g_0^{-1}h)^m g_0^{-1}.
\end{equation}
Its support is in $\langle\supp(h)\rangle$. The signature is that of $g_0$:
locally diagonalize the real metric smoothly, use nonzero real principal
minors there, and observe that infinitesimal changes cannot alter their signs.
The same argument holds at every point of the prolonged halo.

More generally, constant scale factors are allowed. If $B\in\GL_n(K)$ is
constant and $g=B^T(g_0+h)B$, its inverse remains admissible. This includes
strongly anisotropic metrics not having a nondegenerate real standard part.

\subsection{Construction and uniqueness}

An affine connection satisfies
\[
 \nabla_{fX}Y=f\nabla_XY,\qquad
 \nabla_X(fY)=X(f)Y+f\nabla_XY.
\]
Its coefficients are $\nabla_{D_i}D_j=\Gamma^k{}_{ij}D_k$.
The torsion is $T(X,Y)=\nabla_XY-\nabla_YX-[X,Y]$.

\begin{theorem}[Levi--Civita construction in the admissible algebra]
For every admissible metric there is a unique torsion-free, metric-compatible
connection. It is given by
\begin{equation}\label{vtf:eq:christoffel}
 \Gamma^k{}_{ij}=\frac12g^{k\ell}
 (D_i g_{j\ell}+D_j g_{i\ell}-D_\ell g_{ij}).
\end{equation}
\end{theorem}
\begin{proof}
The inverse metric and its derivatives lie in the chosen algebra, so the
formula is defined. It is symmetric in $i,j$, hence torsion-free. Multiplying
by $g_{mk}$ and adding the appropriate permutations gives
$D_i g_{jk}=g_{\ell k}\Gamma^\ell{}_{ij}
+g_{j\ell}\Gamma^\ell{}_{ik}$, which is metric compatibility.
Conversely those three permuted equations, together with symmetry in the
lower indices, force \eqref{vtf:eq:christoffel}. This proves uniqueness.
Coordinate independence can either be checked with the chain rule or read
from the coordinate-free Koszul identity
\begin{align*}
 2g(\nabla_XY,Z)&=Xg(Y,Z)+Yg(Z,X)-Zg(X,Y)\\
 &\quad-g(X,[Y,Z])+g(Y,[Z,X])+g(Z,[X,Y]).
\end{align*}
Every term is an operation already established in the differential algebra.
\end{proof}

Source 03 proves the same construction for regular Riemannian metrics, together
with their pointwise positive definiteness over $K$ at real and halo points:
after normalizing $u\ne0$ by its least component valuation, the leading
coefficient of $g(u,u)$ is $g_0(u_0,u_0)>0$. It also shows that the inverse,
volume density, Hodge star, connection and curvature are Hahn-smooth, with
nonzero supports in $\langle\supp h\rangle$ and the real objects at exponent
zero. These are the Riemannian cases of \eqref{vtf:eq:metricinverse}, of the
signature argument after it, and of \cref{vtf:prop:threshold}, so they are
printed once, here.

On general tensors the connection is determined by the product rule and
contraction compatibility:
\begin{align}\label{vtf:eq:covtensor}
 (\nabla_kT)^{i_1\cdots i_r}{}_{j_1\cdots j_s}
 &=D_kT^{i_1\cdots i_r}{}_{j_1\cdots j_s}\notag\\
 &\quad+\sum_{a=1}^r\Gamma^{i_a}{}_{k\ell}
 T^{i_1\cdots\ell\cdots i_r}{}_{j_1\cdots j_s}
 -\sum_{b=1}^s\Gamma^\ell{}_{kj_b}
 T^{i_1\cdots i_r}{}_{j_1\cdots\ell\cdots j_s}.
\end{align}
Metric compatibility makes raising and lowering indices commute with
$\nabla$. For a function,
$\Hess_g f{}_{ij}=D_iD_jf-\Gamma^k{}_{ij}D_kf$; the Hessian is symmetric
because the connection has no torsion.

\subsection{Parallel transport and geodesic equations}

Along an admitted curve $x(s)$, parallel transport and geodesics satisfy
\[
 \frac{dY^k}{ds}+\Gamma^k{}_{ij}(x(s))\dot x^iY^j=0,
 \qquad
 \ddot x^k+\Gamma^k{}_{ij}(x(s))\dot x^i\dot x^j=0.
\]
These are well-defined differential equations in an admitted calculus, not
yet universal existence assertions. Metric compatibility and the product
rule imply
\[
 \frac{d}{ds}g(Y,Z)=0,\qquad
 \frac{d}{ds}g(\dot x,\dot x)=0
\]
for parallel fields and geodesics, respectively. These quantities are constant
in the coefficientwise calculus on a connected real parameter interval, and
therefore on its Taylor--Hahn prolongation. In a bare intrinsic calculus,
zero derivative alone is not a global constancy theorem; the example in
\cref{vtf:subsec:nonunique} explains the distinction. Existence for positive-support
perturbations of real trajectories is proved in \cref{vtf:sec:flows}; exact
polynomial solutions may also be available outside that class.

\section{Curvature and geometric identities}\label{vtf:sec:curvature}

Fix the convention
\[
 R(X,Y)Z=\nabla_X\nabla_YZ-\nabla_Y\nabla_XZ-\nabla_{[X,Y]}Z,
 \qquad R(D_i,D_j)D_k=R^\ell{}_{kij}D_\ell.
\]
Then
\begin{equation}\label{vtf:eq:curvature}
 R^\ell{}_{kij}=D_i\Gamma^\ell{}_{jk}-D_j\Gamma^\ell{}_{ik}
 +\Gamma^\ell{}_{im}\Gamma^m{}_{jk}
 -\Gamma^\ell{}_{jm}\Gamma^m{}_{ik}.
\end{equation}
In particular $R$ is a tensor although $\Gamma$ is not. Define
\[
 \Ric_{kj}=R^i{}_{kij},\qquad \Scal=g^{kj}\Ric_{kj},\qquad
 G_{ij}=\Ric_{ij}-\tfrac12\Scal\,g_{ij}.
\]

\begin{proposition}[Bianchi identities and conservation]
For an admissible Levi--Civita connection, the first and second Bianchi
identities and the contracted identity hold:
\[
 R(X,Y)Z+R(Y,Z)X+R(Z,X)Y=0,
\]
\[
 (\nabla_XR)(Y,Z)+(\nabla_YR)(Z,X)+(\nabla_ZR)(X,Y)=0,
 \qquad \nabla^iG_{ij}=0.
\]
The curvature has the usual metric symmetries, and $\Ric$ is symmetric.
\end{proposition}
\begin{proof}
Insert \eqref{vtf:eq:curvature} into the cyclic sum. Derivative terms cancel
because $\Gamma^\ell{}_{ij}=\Gamma^\ell{}_{ji}$, and the product terms
cancel in pairs. For the second identity, write the connection matrix
$\Gamma=\Gamma_i\dd x^i$ and curvature matrix
$\Omega=d\Gamma+\Gamma\wedge\Gamma$. Direct expansion, using $d^2=0$,
gives
\[
 d\Omega+\Gamma\wedge\Omega-\Omega\wedge\Gamma=0.
\]
This is the covariant second Bianchi identity. Differentiating
$g(Z,W)$ twice and commuting derivatives gives antisymmetry in the two
metric-paired slots; combining it with the first identity gives pair
symmetry and symmetry of Ricci. Contracting the second identity twice and
using $\nabla g=0$ yields
$\nabla^i\Ric_{ij}=\tfrac12D_j\Scal$, hence the last equation.
All operations are finite tensor contractions in the admissible algebra.
\end{proof}

In dimension two $G$ vanishes identically. In dimensions at least three it
need not. Flatness means $R=0$, not that the coefficients happen to be small.
A nonzero infinitesimal curvature is still nonzero curvature.

\subsection{First-order metric response with an exact remainder}

Let $g=g_0+t^\eta h$, $\eta>0$, with $h$ a real symmetric tensor. The
coefficient at $t^\eta$ of the connection difference is the tensor
\begin{equation}\label{vtf:eq:variationGamma}
 C^k{}_{ij}=\tfrac12g_0^{k\ell}
 (\nabla^0_i h_{j\ell}+\nabla^0_j h_{i\ell}-\nabla^0_\ell h_{ij}).
\end{equation}
Subtracting the metric-compatibility equations for the two connections and
extracting the first coefficient proves this formula. Substitution into
\eqref{vtf:eq:curvature} gives
\begin{equation}\label{vtf:eq:variationR}
 R^\ell{}_{kij}(g)=R^\ell{}_{kij}(g_0)
 +t^\eta(\nabla^0_i C^\ell{}_{jk}-\nabla^0_j C^\ell{}_{ik})
 +E^\ell{}_{kij},\qquad \vH(E)\ge2\eta.
\end{equation}
Here $\vH$ is the minimum nonzero \emph{coefficient-function} exponent in
a fixed local real frame. The remainder is a genuine Hahn tensor. Formula
\eqref{vtf:eq:variationR} is an exact decomposition, rather than an equality in
a ring where $(t^\eta)^2$ is zero.

For $h=\Lie_Vg_0$, naturality under the prolonged diffeomorphisms implies
that the first curvature response is $\Lie_VR(g_0)$. This identifies the
usual infinitesimal coordinate freedom and explains why solving geometric
field equations requires a treatment of gauge, not just componentwise
inversion of a linear operator.

\section{Volume, Hodge duality, and vector calculus in every dimension}\label{vtf:sec:hodge}

\subsection{Volume and the Hodge star}

For an oriented admissible metric of signature $(p,q)$ set
\[
 \mu_g=\sqrt{|\det g|},\qquad
 \vol_g=\mu_g\dd x^1\wedge\cdots\wedge\dd x^n.
\]
For \eqref{vtf:eq:regularmetric}, the square root is obtained by the positive
binomial expansion after factoring $|\det g_0|$. The induced inner product
on forms is the determinant pairing. Define $*$ by
\begin{equation}\label{vtf:eq:star}
 \alpha\wedge *\beta=\langle\alpha,\beta\rangle_g\vol_g,
 \qquad *:\Omega^k_\A\to\Omega^{n-k}_\A.
\end{equation}
An oriented pseudo-orthonormal basis gives
\begin{equation}\label{vtf:eq:star2}
 *^2\alpha=(-1)^{k(n-k)+q}\alpha.
\end{equation}
Indeed the two permutations contribute $(-1)^{k(n-k)}$, and the product of
all metric signs is $(-1)^q$. This proves the formula without a positivity
assumption.

\subsection{Gradient, divergence, and Laplacians}

The definitions and their coordinate formulas are
\begin{align}
 \grad_g f&=(df)^\sharp=g^{ij}D_jf\,D_i,\label{vtf:eq:grad}\\
 \Lie_X\vol_g&=(\Div_gX)\vol_g,
 &\Div_gX&=\mu_g^{-1}D_i(\mu_gX^i)=\nabla_iX^i,\label{vtf:eq:div}\\
 \Delta_gf&=\Div_g\grad_gf
 =\mu_g^{-1}D_i(\mu_g g^{ij}D_jf).\label{vtf:eq:lap}
\end{align}
For \eqref{vtf:eq:div}, use Cartan's formula and expand
$d\iota_X\vol_g$. Equality with the connection trace follows from
$D_i\det g=(\det g)\tr(g^{-1}D_i g)$ and
$\Gamma^j{}_{ji}=D_i\log\mu_g$, with the latter notation meaning the
algebraic quotient $\mu_g^{-1}D_i\mu_g$, whether or not a logarithm has
been selected in the function ring.

On $k$-forms define the codifferential by
\begin{equation}\label{vtf:eq:codiff}
 \delta_g=(-1)^{n(k+1)+q+1}*d*,\qquad
 \Delta_H=d\delta_g+\delta_gd.
\end{equation}
The sign follows by integrating $d(\alpha\wedge*\beta)$ on real compact
chains, or by the equivalent local algebraic identity. On functions
$\Delta_H=-\Delta_g$. This distinguishes the nonnegative real Riemannian
Hodge-Laplacian convention from the positive-coordinate-second-derivative
convention in \eqref{vtf:eq:lap}.

\subsection{What replaces curl?}

The dimension-independent vorticity is the two-form
\[
 \mathcal W_X=d(X^\flat).
\]
Its dual $*\mathcal W_X$ has degree $n-2$. It is a scalar in dimension two,
a one-form (hence a vector after raising an index) in dimension three, and a
two-form in dimension four. In dimension one it vanishes. Thus a vector-valued
curl is not a dimension-independent operation.

Similarly, the metric and orientation give an $(n-1)$-ary cross product
\[
 X_1\times\cdots\times X_{n-1}
 =\bigl(*(X_1^\flat\wedge\cdots\wedge X_{n-1}^\flat)\bigr)^\sharp.
\]
It is the familiar binary cross product only when $n=3$. Extra exceptional
algebraic structures in other dimensions are not supplied by merely changing
the scalar field.

The identities $d^2=0$ and $\delta_g^2=0$ replace the usual ``curl of a
gradient'' and ``divergence of a curl'' formulas. In a flat constant metric,
\[
 \delta_gd(X^\flat)=d(\Div_gX)-\Delta_g(X^\flat),
\]
where the final Laplacian is taken componentwise. For a variable metric the
Hodge identity remains valid, while rewriting it using a rough covariant
Laplacian introduces the usual Ricci curvature term. The flat formula must
not be transplanted unchanged to curved coordinates or a curved metric.

\begin{remark}[Specializations checked for the merge]\label{vtf:rem:specializations}
For $n=3$ and Euclidean signature ($q=0$), \eqref{vtf:eq:codiff} gives
$\delta_g=(-1)^k*d*$ on $k$-forms, which is source 03's convention in Part~II.
For $n=4$ and signature $(1,3)$, so $q=3$, it gives $\delta=*d*$ in every degree,
and \eqref{vtf:eq:star2} gives $*^2=(-1)^{k(4-k)+3}$; these are source 02's
formulas in Part~III. For $\eta=\diag(1,-1,-1,-1)$ one has $\Delta_\eta=\Box$,
so $\Delta_H=-\Delta_g$ on functions is the degree-zero case of source 02's
$(d\delta+\delta d)\alpha=-\Box\alpha$. The vector-valued curl of Part~II is
$(*d(X^\flat))^\sharp$, and the four-dimensional substitutes for the cross
product are in \cref{vtf:mink:sec:tensors}.
\end{remark}

\section{Integration, Stokes, and local potentials}\label{vtf:sec:integration}

\subsection{Integration on real chains is coefficientwise}

For a smooth real oriented compact $p$-chain $C$ and a Hahn $p$-form defined
near it, use a finite localization and define
\begin{equation}\label{vtf:eq:integral}
 \int_C^{H}\sum_\gamma\alpha_\gamma t^\gamma
 =\sum_\gamma\left(\int_C\alpha_\gamma\right)t^\gamma.
\end{equation}
The sum is well defined because its support is a subset of a single
well-ordered support. Real integration happens \emph{inside each coefficient}
first. Formula \eqref{vtf:eq:integral} is not defined as a fine limit of
surreal-valued Riemann sums. It is independent of finite real chart choices
and is $K$-linear by the finite-pair support lemma.

\begin{theorem}[Hahn Stokes theorem]\label{vtf:thm:stokes}
For $C$ a smooth real compact oriented $p$-chain and
$\alpha\in\Omega^{p-1}_\A$ defined near it,
\[
 \int_C^{H}d\alpha=\int_{\partial C}^{H}\alpha.
\]
\end{theorem}
\begin{proof}
Every coefficient form satisfies real Stokes. Taking those equalities as the
coefficients of one Hahn series proves the identity. No exchange of an
uncertified infinite integral with an infinite sum is involved: the integral
was defined coefficientwise on a common support.
\end{proof}

For a real domain with smooth compact boundary this yields
\[
 \int_\Omega^{H}\Div_gX\,\vol_g
 =\int_{\partial\Omega}^{H}\iota_X\vol_g,
\]
provided the metric, inverse, density and field are admissible near the
closure. The flux form is the primary definition; a unit-normal expression
requires the corresponding induced metric and square roots.

For noncompact integration domains each coefficient must satisfy the relevant
ordinary integrability condition. Fubini, integration by parts and
differentiation under the integral are available coefficientwise under their
classical hypotheses; a common compact support or Schwartz coefficients are
convenient sufficient assumptions, not consequences of coherence (source 02).
When the integrand has nonnegative support, standard part commutes with both
sides of \cref{vtf:thm:stokes} (source 02; see also \cref{vtf:thm:reduction}).
Coefficientwise Stokes with complex coefficients on any smooth manifold with
corners is \path{contours:thm:hahnstokes}.

\begin{proposition}[Positivity for smooth common-support integrands]
Let $M$ be a nonempty compact real manifold without boundary, with a positive
real smooth density $\nu$. If $f\in\A_\Gamma(M)$ is pointwise nonnegative
at every real point and is not the zero field, then $\int_M^H f\nu>0$.
\end{proposition}
\begin{proof}
Let $\gamma_0$ be the least exponent whose smooth coefficient function is
not identically zero. Pointwise nonnegativity forces
$f_{\gamma_0}(x)\ge0$ everywhere: wherever it is nonzero it is the leading
coefficient of $f(x)$. It is positive on a nonempty real open set. Hence its
real integral is positive, and this is the leading coefficient of the Hahn
integral. The same argument works for compact domains with interior whenever
the nonzero smooth field is detected there.
\end{proof}

This result concerns smooth common-support densities on real domains. It is
not a statement about arbitrary Hahn-valued measures or a change from
coefficientwise to strong countable additivity.

\subsection{A support-preserving Poincar\'e lemma}

Suppose $U\subset\R^n$ is star-shaped about zero. For a real $p$-form,
$p\ge1$, define the radial homotopy
\[
 (H\alpha)_x(v_1,\ldots,v_{p-1})
 =\int_0^1 s^{p-1}\alpha_{sx}(x,v_1,\ldots,v_{p-1})\dd s.
\]
Extend it coefficientwise to Hahn forms, and put $H=0$ on functions.

\begin{theorem}[Support-preserving local exactness]\label{vtf:thm:poincare}
For common-support forms in $\Omega^p(U;\R)((t^\Gamma))$, the homotopy
never enlarges Hahn support and satisfies
\[
 dH+Hd=\id-e_0^*,
\]
where $e_0:U\to U$ is the constant map. Consequently every closed positive-
degree common-support Hahn form on $U$ is exact with a primitive having
no additional exponents. This proves local exactness for the sheaf as well. A Hahn function with all $D_i f=0$ is constant on every connected
real domain.
\end{theorem}
\begin{proof}
For the dilation $r_s(x)=sx$, differentiating $r_s^*\alpha$ in $s$ and
using Cartan's identity gives a total derivative whose integral from $0$ to
$1$ is $\alpha-e_0^*\alpha$. This is the displayed real homotopy identity;
each coefficient is smooth at $s=0$ because the factor $s^{p-1}$ is
nonnegative-degree. Applying the identity coefficientwise proves it for
Hahn forms. For a positive-degree form $e_0^*\alpha=0$. Finally, a real
coefficient function with zero gradient is constant on a connected real
open set, and the same conclusion holds at every Hahn exponent.
\end{proof}

In particular, a curl-free field has a scalar potential locally precisely
when $d(X^\flat)=0$; ``curl-free'' here means the two-form condition in every
dimension. A divergence-free field satisfies
$d(\iota_X\vol_g)=0$ and locally has an $(n-2)$-form potential when $n\ge2$.
These are differential-form statements, not a special three-dimensional
vector-potential ansatz.

There is also a purely algebraic version on $K^n$ for polynomial forms (source
02, stated on $K^4$; the argument uses no property of dimension four). Perform
the radial homotopy with the $s$-integral replaced by the formal rule
$s^m\mapsto1/(m+1)$. The identity $dH+Hd=\id-e_0^*$ is then a polynomial identity
with rational coefficients. This proves the polynomial Poincar\'e lemma in
characteristic zero without an integral over a surreal interval. The explicit
three-dimensional scalar and vector potentials are
\cref{vtf:3d:prop:potentials}.

\subsection{Integration on surreal-coordinate simplices}

There is a second, algebraic integration theory that does not require a real
base. For polynomial forms over $K$ on an oriented affine $p$-simplex with
vertices in $K^n$, pull back to barycentric coordinates
$\lambda_0+\cdots+\lambda_p=1$ and define its moments by
\begin{equation}\label{vtf:eq:simplex}
 \int_{\Delta^p}^{\mathrm{alg}}
 \lambda_0^{a_0}\cdots\lambda_p^{a_p}
 \dd\lambda_1\cdots\dd\lambda_p
 =\frac{a_0!\cdots a_p!}{(p+a_0+\cdots+a_p)!}.
\end{equation}
Extend by $K$-linearity and the oriented affine pullback. Polynomial Stokes
holds. One proof applies the polynomial fundamental theorem successively in
simplex coordinates. Another observes that, at each fixed degree, both sides
are polynomial expressions in the coefficients and vertices with rational
coefficients; their equality for all real data is a polynomial identity and
hence remains true over $K$.

A support-certified formal form on an infinitesimal scaled simplex can also
be integrated termwise by these moments. The same positive-support proof
justifies its sums. This supplies exact small-scale integrals, but neither
construction defines a countably additive Lebesgue theory on every subset of
$K^n$. Classical global integral theorems require the chain category stated
in their hypotheses.

A smooth real parameter domain may also describe a surreal-valued chain through
an admissible Taylor or scale-chart map. Its integral is defined by the
pulled-back form, provided that pullback has a common well-ordered support and
commutes with $d$. \Cref{vtf:thm:prolongation} and the polynomial constructions
supply these properties for the stated chart classes, and Stokes then follows on
the ordinary parameter domain (source 03; compare
\path{contours:thm:deformedstokes}). A general pointwise surreal map has
neither property, and the parameter map is usually not intrinsically
continuous. The parametric sphere integrals of \cref{vtf:subsec:radial} and
\cref{vtf:3d:subsec:monopole} are to be read in this way.

\section{Global cohomology and a precise Hodge theorem}\label{vtf:sec:cohomology}

Let $M$ be a compact smooth real manifold without boundary. Denote by
$\Omega_H^p(M)$ the global sections of the Hahn form sheaf. Compactness of
the \emph{real} base gives
\[
 \Omega_H^p(M)=\Omega^p(M;\R)((t^\Gamma)).
\]
This is a statement about supports of real coefficient forms, not compactness
of the much larger halo in its intrinsic topology.

\begin{theorem}[Hahn extension of a real cochain complex; sources 02 and 03]\label{vtf:3d:thm:cohomology}
Let $(E^\bullet,d)$ be a cochain complex of real vector spaces. There is a
natural isomorphism
\[
 H^k\bigl(E^\bullet((t^\Gamma)),d\bigr)\ \cong\ H^k(E^\bullet,d)((t^\Gamma)).
\]
If $H^k(E^\bullet,d)$ is finite dimensional, the right side is naturally
$H^k(E^\bullet,d)\otimes_\R K$.
\end{theorem}

\begin{proof}
A Hahn series is closed exactly when each coefficient is closed. Send it to the series of coefficient cohomology classes. This map is natural and has support contained in the original support. It is surjective: choose a representative for each nonzero coefficient class, retaining the same support. Its kernel consists of series all of whose coefficients are exact. Choose a primitive for each such coefficient; their series again has the same well-ordered support and differentiates to the given series. Thus the kernel is exactly the exact Hahn series. These existence choices are ordinary set-sized choices, not a proper-class selection. In finite dimension, expansion in a finite basis identifies the Hahn coefficient space with the ordinary tensor product.
\end{proof}

In infinite dimension an arbitrary Hahn family of cohomology classes need not
lie in a finite-dimensional real span, so writing the right side as an algebraic
tensor product with $K$ would lose information (source 02). Applied to the de
Rham complex of a real manifold with globally supported coefficients (source 02
on a fixed real domain, source 03 on a manifold), the theorem gives
\[
 H^p\bigl(\Omega^\bullet(M;\R)((t^\Gamma)),d\bigr)\cong
 H^p_{\mathrm{dR}}(M;\R)((t^\Gamma)).
\]
This is the cohomology of the specified coefficientwise complex, not singular
cohomology of the halo or of $K^n$. Betti dimensions are unchanged after scalar
extension, while periods can be arbitrary elements of $K$. The complex-coefficient
version is \path{contours:thm:cohomology}. For compact $M$ the next theorem,
proved independently by source 01, adds the graded-algebra structure.

\begin{theorem}[Compact de Rham base change]\label{vtf:thm:derham}
There is a canonical isomorphism of graded $K$-algebras
\begin{equation}\label{vtf:eq:cohom}
 H^\bullet(\Omega_H^\bullet(M),d)
 \cong H^\bullet_{\mathrm{dR}}(M;\R)\otimes_\R K.
\end{equation}
In particular, the dimension in degree $p$ is the ordinary real Betti number
$b_p(M)$.
\end{theorem}
\begin{proof}
The real cohomology in each degree is finite-dimensional. Choose real closed
forms $\theta_1,\ldots,\theta_b$ representing a basis. Choose real linear
coordinate maps on closed forms and a linear choice of primitive for exact
forms; such a choice exists by a vector-space splitting. Let $P$ apply that
primitive choice after subtracting the selected cohomology representative.
Thus every closed
real form has an identity
\[
 \alpha=\sum_{j=1}^b c_j(\alpha)\theta_j+dP\alpha,
\]
where each $c_j$ and $P$ is real linear. Apply these fixed operators to every
coefficient of a closed Hahn form. They cannot add exponents. We obtain
\[
 \alpha=\sum_{j=1}^b a_j\theta_j+d\beta,
 \qquad a_j\in K,\quad\beta\in\Omega_H^{p-1}(M).
\]
This proves surjectivity of the canonical map from the right side of
\eqref{vtf:eq:cohom}. If a finite $K$-linear combination of the $\theta_j$ is
exact, coefficient extraction gives a real exact combination at every
exponent. Real independence forces all its coefficients to vanish, proving
injectivity. The map itself sends scalar-extended real cohomology classes to
the same representative forms, so it is independent of the auxiliary
splittings. Products agree because wedge products have finite coefficient
convolutions.
\end{proof}

For an oriented compact $M$ choose an ordinary real Riemannian metric $g_0$.
The classical real Hodge theorem \cite{Warner} supplies harmonic projection
$P_0$ and Green operator $G_0$ on smooth forms, with
\[
 \id=P_0+d\delta_0G_0+\delta_0dG_0.
\]
Every operator in this identity extends coefficientwise without enlarging
support. We immediately obtain the following exact, but carefully scoped,
Hodge theorem.

\begin{corollary}[Coefficientwise Hodge decomposition]\label[corollary]{vtf:cor:hodge}
For the real metric $g_0$,
\[
 \Omega_H^p(M)=d\Omega_H^{p-1}(M)
 \oplus\delta_0\Omega_H^{p+1}(M)
 \oplus\bigl(\HH_{g_0}^p(M;\R)\otimes_\R K\bigr).
\]
The summands are orthogonal for the $K$-valued pairing
$\int_M^H\alpha\wedge *_0\beta$. This pairing is positive definite.
Every Hahn cohomology class has a unique harmonic representative.
\end{corollary}
\begin{proof}
Extend the real operator identities and orthogonality coefficientwise.
Uniqueness follows by applying the real projections at each exponent.
For positivity, the leading coefficient of the pointwise squared norm of a
nonzero form is the real squared norm of its leading nonzero coefficient
form; its integral is positive. Finite dimensionality of the harmonic space
identifies its Hahn extension with tensoring by $K$.
\end{proof}

Source 03 proves the same decomposition for a real background metric, and notes
that the three projections $P_0$, $d\delta_0G_0$ and $\delta_0dG_0$ preserve
support and that the harmonic space has $K$-dimension $b_k(M)$.

This does not claim an unrestricted Hodge theorem for every surreal metric,
every intrinsic $K$-manifold, or a completed infinite-dimensional normed
space. The fixed real elliptic theory is the input. For a regular Riemannian
Hahn metric, source 03 obtains the $g$-orthogonal decomposition without
constructing a new Green operator (\cref{vtf:3d:thm:regularhodge}); for metrics
without a nondegenerate real leading part no Hodge theorem is proved here.

\subsection{Hodge decomposition for a regular Riemannian Hahn metric}
\Cref{vtf:cor:hodge} used a real metric. A positive-order change of metric can also be handled without postulating a new surreal Hilbert-space theory.

Let $g=g_0+h$ be a regular Riemannian metric on a compact oriented manifold
without boundary. For each degree $k$, let $R_k=I+Q_k$ be the pointwise bundle operator defined by
\[
 (\alpha,\beta)_g=(R_k\alpha,\beta)_{g_0}.
\]
It is obtained from the inverse metrics and volume densities, is self-adjoint relative to $g_0$, and $Q_k$ has positive support. It has an admissible inverse by \cref{vtf:3d:thm:inverse}. Comparing formal adjoints gives
\begin{equation}\label{vtf:3d:eq:adjointR}
 \delta_g=R_{k-1}^{-1}\delta_0R_k
 \quad\text{on degree }k.
\end{equation}

\begin{theorem}[Regular-metric Hahn--Hodge theorem; source 03]\label{vtf:3d:thm:regularhodge}
For a compact oriented boundaryless real manifold and a regular Riemannian
Hahn metric,
\[
 \Omega_\Gamma^k(M)=\dd\Omega_\Gamma^{k-1}(M)
 \ \mathbin{\perp_g\oplus}\ \mathcal H_g^k
 \ \mathbin{\perp_g\oplus}\ \delta_g\Omega_\Gamma^{k+1}(M),
\]
where $\mathcal H_g^k=\ker\dd\cap\ker\delta_g=\ker\Delta_{H,g}$.
Every cohomology class has a unique $g$-harmonic representative, and
$\dim_{K}\mathcal H_g^k=b_k(M)$. If $S$ is the positive metric support and $T$ the input support, all three components have support in $T+\langle S\rangle$.
\end{theorem}
\begin{proof}
Let $Z=\ker\dd$ and $E=\im\dd$ in the ordinary real degree-$k$ space. Real Hodge theory supplies support-preserving orthogonal projections $P_Z^0$ and $P_E^0$. Closed and exact Hahn forms are respectively $Z((t^\Gamma))$ and $E((t^\Gamma))$ by \cref{vtf:3d:thm:cohomology} and its proof.

For $W=Z$ or $E$, the operator
\[
 P_W^0R_k\big|_{W((t^\Gamma))}:W((t^\Gamma))\to W((t^\Gamma))
\]
is $I$ plus a positive-support operator. It is invertible by \cref{vtf:3d:thm:inverse}. Hence
\[
 P_W^g\alpha=
 \left(P_W^0R_k\big|_{W((t^\Gamma))}\right)^{-1}P_W^0R_k\alpha
\]
is well defined. Its defining equation says that $R_k(\alpha-P_W^g\alpha)$ is $g_0$-orthogonal to $W$, equivalently that $\alpha-P_W^g\alpha$ is $g$-orthogonal to $W$. Thus it is the desired $g$-orthogonal projection.

Write
\[
 e=P_E^g\alpha,\qquad
 h=P_Z^g\alpha-P_E^g\alpha,\qquad
 c=\alpha-P_Z^g\alpha.
\]
Because $E\subseteq Z$, $P_E^gP_Z^g=P_E^g$, and these three fields are mutually orthogonal. The field $h$ is closed and orthogonal to exact fields; integration by parts implies $\delta_gh=0$. The pairing is positive definite by the leading-coefficient argument used in \cref{vtf:cor:hodge}, now with leading metric $g_0$.

To show $c$ is coexact, observe that $R_kc$ is $g_0$-orthogonal to all closed forms. Real Hodge theory coefficientwise gives $R_kc=\delta_0\beta$. By \eqref{vtf:3d:eq:adjointR},
$c=\delta_g(R_{k+1}^{-1}\beta)$. Conversely every coexact form is orthogonal to closed forms. This proves the decomposition and its uniqueness. Closed forms therefore split uniquely into exact and harmonic parts, giving the cohomology isomorphism and dimension statement. Finally,
$(\Delta_{H,g}\alpha,\alpha)_g=\norm{\dd\alpha}_g^2+\norm{\delta_g\alpha}_g^2$ proves the asserted description of the Laplacian kernel. Every inverse used above is a positive-support Neumann inverse, so the support bound follows.
\end{proof}
The theorem does not cover arbitrary singular standard-part metrics such as $\operatorname{diag}(1,t^2,t^4)$ without additional analysis. A metric being pointwise positive definite is not a substitute for the uniform expansion hypotheses.

\begin{remark}[Dimension and signature in \cref{vtf:3d:thm:regularhodge}]
Source 03 proves this theorem inside its three-dimensional development. Its
statement and proof use only a compact oriented boundaryless real manifold, real
Hodge theory for $g_0$, \cref{vtf:3d:thm:cohomology} and
\cref{vtf:3d:thm:inverse}; no property of dimension three is used, so it is
printed here in every dimension. Positive definiteness of $g_0$ is used, for the
positivity of the pairing, and remains a hypothesis: regular metrics of
indefinite signature are not covered.
\end{remark}

For example, on the real $n$-torus,
$\dim_K H_H^p=\binom np$. On a circle, $a\dd\theta$ has period $2\pi a$
and is nonexact for every $a\in K\setminus\{0\}$, even when $a$ is
infinitesimal. Scales enrich the coefficients of topological classes; in this
model they do not create extra Betti numbers. This cohomology is not asserted
to be singular cohomology of the halo with its fine, highly disconnected
topology.

\section{Standard-part geometry and scale stability}\label{vtf:sec:reduction}

Let $\A_{\ge0}$ be the subalgebra with no negative exponents, and
$\A_{>0}$ its ideal of positive-support sections. Coefficient extraction
extends standard part to
\[
 \st:\A_{\ge0}\to C^\infty(M),\qquad
 \ker\st=\A_{>0}.
\]
The maps $D_i$ commute with it because they act only on coefficient functions.

\begin{theorem}[Regular reduction of tensor geometry]\label{vtf:thm:reduction}
Suppose $g\in\A_{\ge0}$ is symmetric and $g_0=\st(g)$ is a real smooth
nondegenerate metric. For tensor fields with nonnegative support,
standard part commutes with tensor products, contractions, exterior
products, $d$, Lie brackets and Lie derivatives. Moreover,
\[
 \st(g^{-1})=g_0^{-1},\quad
 \st(\Gamma(g))=\Gamma(g_0),\quad
 \st(R(g))=R(g_0),\quad
 \st(\Ric(g))=\Ric(g_0),
\]
and similarly for scalar curvature, volume density, Hodge star,
gradient and divergence. Coefficientwise real-chain integrals commute with
standard part whenever the integrand has nonnegative support.
\end{theorem}
\begin{proof}
All finite polynomial expressions commute with the ring homomorphism $\st$,
and it commutes with $D_i$. Formula \eqref{vtf:eq:metricinverse} shows that
$g^{-1}$ has no negative exponents and has constant coefficient $g_0^{-1}$.
The binomial density formula has constant coefficient
$\sqrt{|\det g_0|}$. Substituting into the explicit formulas for connection,
curvature, star and divergence proves every assertion. Integration commutes
by its definition. At a halo point $a+h$, positive powers of $h$ disappear
under standard part, so these reductions take place at the real base point
$a$.
\end{proof}

\begin{proposition}[A quantitative support threshold]\label[proposition]{vtf:prop:threshold}
If $g=g_0+h$ and $\supp(h)\subseteq[\eta,\infty)$ for some $\eta>0$,
then the differences from the corresponding real quantities for $g^{-1}$,
$\Gamma$, $R$, $\Ric$, $\Scal$ and the volume density all have support in
$[\eta,\infty)$, locally in real coordinates.
\end{proposition}
\begin{proof}
Differentiation does not lower coefficient exponents. Multiplication by real
smooth coefficients does not change them. In every difference expansion,
including the inverse and binomial expansions, each term contains at least
one factor of $h$ or one of its real derivatives. It therefore has exponent
at least $\eta$. The expansions are strong by the positive-support lemma.
\end{proof}

The hypothesis is a support bound on a field and its chosen smooth
coefficient representation, not merely smallness of its values at a few
points. It also uses a nondegenerate real leading metric. If
$g=\operatorname{diag}(t^{2\eta},1,\ldots,1)$, the metric is perfectly
nondegenerate over $K$, but its standard part is degenerate and its inverse
contains $t^{-2\eta}$. The theorem intentionally excludes this situation.

\begin{remark}[Additions from sources 02 and 03]
For a regular Riemannian metric, source 03 adds to \cref{vtf:thm:reduction} the
$g$-musical maps, the covariant derivative, $\st(\nabla^gT)=\nabla^{g_0}\st T$,
and the Hodge projections of \cref{vtf:3d:thm:regularhodge}. These reduce to the
$g_0$-projections because every Neumann inverse in their construction has the
real inverse as constant term. Consequently a regular, nonnegative-support
solution of a polynomial field equation reduces to a real solution of its
leading equation; this is a consistency statement, not a converse existence
theorem. Source 02 proves the scalar part on real spacetime,
\begin{equation}\label{vtf:mink:eq:stdifferential}
 \st(fg)=\st(f)\st(g),\qquad \st(D_if)=D_i\st(f)
 \qquad(f,g\in\A_{\ge0}(U)),
\end{equation}
an exact differential-algebra quotient rather than an interchange of a limit and
a derivative. It also notes that standard part preserves rational expressions
only when their denominators have nonvanishing real body on the chosen domain:
$\eps$ is finite with an unlimited inverse, and reduction cannot turn
$\eps\eps^{-1}=1$ into a product of real shadows. The collection's
\path{phys:thm:reduction} is the four-dimensional Lorentzian case with the
Einstein tensor.
\end{remark}

\subsection{Information that standard part destroys}\label{vtf:3d:subsec:destroys}
Standard part annihilates every positive-order observable (source 03). A monopole flux $4\pi t^\eta q_0$ becomes zero although it obstructs a global vector potential. The area of the vectors $(1,0,0)$ and $(1,t^\eta,0)$ becomes zero although they are exactly independent. Curvature of order $t$ in \eqref{vtf:3d:eq:conformal} disappears from the ordinary metric.

To retain this information use the leading normalized observable
\[
 \operatorname{in}(F)=t^{\vH(F)}F_{\vH(F)},
\]
or record several successive coefficients. This associated-scale data is covariant under real frame changes. A surreal frame change may shift the hierarchy and must be recorded with the observable.

\section{Examples at finite, infinitesimal, and infinite scales}\label{vtf:sec:examples}

\subsection{Anisotropic flat geometry}

For constant $\lambda_i\in\Gamma$ let
\[
 g=\sum_{i=1}^n t^{2\lambda_i}(\dd x^i)^2.
\]
It has
\[
 g^{-1}=\operatorname{diag}(t^{-2\lambda_i}),\quad
 \mu_g=t^{\lambda_1+\cdots+\lambda_n},\quad
 \Gamma=0,\quad R=0,
\]
and
\[
 \grad_g f=\sum_i t^{-2\lambda_i}D_i f\,D_i,
 \qquad \Delta_g f=\sum_i t^{-2\lambda_i}D_i^2f.
\]
The coordinate change $y^i=t^{\lambda_i}x^i$ makes the metric Euclidean
over $K$. It is generally not an ordinary real coordinate change. Different
component scales can be an exact change of units or frame, without creating
curvature. Source 03's three-dimensional metric $\diag(1,t^2,t^4)$
(\cref{vtf:3d:subsec:anisotropic}) shows in addition that raising an index and
taking standard part do not commute when the standard part is degenerate.

\subsection{A smooth intrinsic metric with infinite curvature}\label{vtf:subsec:warped}

For $n\ge2$ and $\eps=t^\eta>0$ consider, on genuine $K$ coordinates,
\begin{equation}\label{vtf:eq:warped}
 g_\eps=(\dd x)^2+(x^2+\eps^2)(\dd y)^2
          +\sum_{j=3}^n(\dd z^j)^2.
\end{equation}
It is positive definite at every point. Its coefficients and inverse are
intrinsically smooth rational functions. The only nonzero Christoffel
coefficients in the two-dimensional factor are
\[
 \Gamma^x{}_{yy}=-x,\qquad
 \Gamma^y{}_{xy}=\Gamma^y{}_{yx}=\frac{x}{x^2+\eps^2}.
\]
Our curvature convention gives
\begin{equation}\label{vtf:eq:warpedcurv}
 \mathcal K_\eps=-\frac{\eps^2}{(x^2+\eps^2)^2},\qquad
 \Scal(g_\eps)=-\frac{2\eps^2}{(x^2+\eps^2)^2}.
\end{equation}
The extra flat factors do not change the scalar curvature. These formulas
follow directly from \eqref{vtf:eq:curvature}; the supplied symbolic program
checks the calculation independently.

At $x=0$, the sectional curvature is $-\eps^{-2}$, an infinite surreal.
At every fixed nonzero real $x$, it has leading term $-\eps^2/x^4$.
The inverse coefficient $(x^2+\eps^2)^{-1}$ has no smooth Hahn expansion
on a real neighborhood crossing zero. Thus this is an intrinsic metric, not
an admissible metric over that whole real base in the algebra
\eqref{vtf:eq:A}. There is no contradiction between those conclusions.

The inner coordinate $x=\eps X$ gives the two-dimensional factor
\[
 \eps^2\bigl((\dd X)^2+(1+X^2)(\dd y)^2\bigr),
 \qquad \mathcal K_\eps=-\eps^{-2}(1+X^2)^{-2}.
\]
A scale chart resolves the transition and exposes its curvature scale.
Replacing a small real parameter by a nonzero infinitesimal removes a zero
of this particular metric coefficient, but does not make the curvature
finite. Nondegeneracy, regular standard part and bounded invariants are
three different requirements.

\subsection{Radial fields and scale-independent flux}\label{vtf:subsec:radial}

In Euclidean $K^n$ let $r=(\sum_i(x^i)^2)^{1/2}$ and, away from $r=0$, set
\[
 X^i=c\,x^i r^{-n},\qquad c\in K.
\]
The local algebraic derivative of the positive square root gives
$D_i r=x^i/r$. Consequently
\[
 \Div X=c\left(nr^{-n}-nr^{-n-2}\sum_i(x^i)^2\right)=0.
\]
For $n\ne2$ this is the gradient of $c\,r^{2-n}/(2-n)$.
For $n=2$ a local logarithmic primitive may be chosen on an admitted positive
radius chart; no unrestricted global trigonometric or logarithmic function
is needed for the divergence calculation.

On the sphere parametrized by $x=R\theta$, $R\in K_{>0}$ and
$\theta\in S^{n-1}(\R)$, the pulled-back flux form is exactly
$c$ times the real unit-sphere area form: the factor $R^{1-n}$ in the field
cancels $R^{n-1}$ in the area. Define this parametric integral by its smooth
Hahn coefficients. Its value is
\[
 c\,\operatorname{area}(S^{n-1}),
\]
for infinitesimal, finite or infinite $R$. This is an explicit admissible
parametric calculation, not a claim that $\theta\mapsto R\theta$ is a
nonconstant fine-continuous map from a usual real sphere into full $\No^n$.
The central singularity at $r=0$ has not been assigned a value. For $n=3$ this
is the monopole of \cref{vtf:3d:subsec:monopole}, where Stokes on the sphere
also rules out a global vector potential; the Coulomb field of
\cref{vtf:mink:sec:limits} is the same field in electrostatic language.

\section{Flows and parallel transport by positive-support recursion}\label{vtf:sec:flows}

Let $X_0$ be a real smooth vector field on an open $U\subset\R^n$. Suppose
its real flow $\phi_0(s,a)$ is defined smoothly for $(s,a)$ in a fixed real
flow domain containing $s=0$, and stays in $U$. Consider
\[
 X=X_0+\sum_{\gamma\in S}t^\gamma X_\gamma,
 \qquad S\subset\Gamma_{>0}\text{ well ordered},
\]
with all coefficient fields smooth on $U$.

\begin{theorem}[Near-real Hahn flows]\label{vtf:thm:flow}
There is a solution of
\[
 \partial_s\phi=(\ev X)(\phi),\qquad \phi(0,a)=a,
\]
of the form
\[
 \phi(s,a)=\phi_0(s,a)+\sum_{\gamma\in\langle S\rangle\setminus\{0\}}
 t^\gamma u_\gamma(s,a).
\]
Each coefficient is real smooth on the same real flow domain, the support is
well ordered, and the equation holds exactly in the Taylor--Hahn calculus.
The solution is unique among positive-support Hahn corrections with the
specified initial condition. The realized flow maps satisfy the local flow
composition law wherever all terms are defined.
\end{theorem}
\begin{proof}
Expand each coefficient vector field by its Taylor jet at $\phi_0$.
At exponent $\gamma$, the unknown $u_\gamma$ occurs only in the linear term
from $X_0$. Every other occurrence contains a strictly positive lower
exponent, either from $X-X_0$ or another correction. Thus
\[
 \partial_s u_\gamma
 =DX_0(\phi_0)u_\gamma+F_\gamma,
 \qquad u_\gamma(0,a)=0,
\]
where $F_\gamma$ is a finite expression in already constructed coefficients
and real derivatives. Finiteness even across all Taylor orders is
\cref{vtf:lem:neumann}. The ordinary inhomogeneous linear variational equation
has a unique smooth solution on the given real flow domain, by its
fundamental-matrix formula. This defines the coefficients by transfinite
recursion over $\langle S\rangle$.

The family is Hahn-supported by construction. Substitution back into the
Taylor--Hahn equation is strong and yields equality at every exponent.
For uniqueness against a second solution, take the union of the two positive
supports and the forcing support and apply the same least-differing-exponent
linear equation. The coefficient at that exponent must agree, a contradiction.
Finally, composing two flow maps is admissible by \eqref{vtf:eq:nonrealpull};
the usual autonomous equation and matching initial data, followed by this
uniqueness, prove the local flow law. The identities can be prolonged in the
initial coordinates by \cref{vtf:thm:prolongation}.
\end{proof}

The real flow domain is deliberately part of the hypothesis. This result
asserts neither completeness of arbitrary real fields nor existence of an
all-real-time flow for a field with infinite leading coefficients.
For example $X=t^{-\eta}D_1$ has the exact intrinsic solution
$x^1(s)=a^1+s t^{-\eta}$, which leaves the finite halo for nonzero real $s$.
The rescaled time $s=t^\eta\tau$ makes this motion finite.

For $n\ge2$, the field
\[
 X=t^\eta(x^1D_1-x^2D_2)
\]
has the exact flow
\[
 x^1(s)=a^1\sum_{m\ge0}\frac{(s t^\eta)^m}{m!},\qquad
 x^2(s)=a^2\sum_{m\ge0}\frac{(-s t^\eta)^m}{m!}.
\]
Its Jacobian determinant is $1$, by the strong exponential addition identity;
the other coordinates are fixed. This illustrates an exact infinitesimal
strain with exact volume conservation.

The same recursion proves parallel transport for a connection that is a
positive-support perturbation of a real connection along a near-real curve.
Geodesics reduce to a first-order system on the position--velocity space,
so \cref{vtf:thm:flow} applies near any real geodesic on its specified real
existence interval. Metric conservation then holds coefficientwise and on
the intrinsic realizations.

\section{Nonlinear field equations and a support-controlled lifting theorem}\label{vtf:sec:PDE}

The preceding ODE construction generalizes to a large class of equations for
scalar, vector or tensor fields. The point is not to assume that every real
PDE is solvable. It is to identify the exact real solvability input and then
prove the Hahn extension.

Let $E,F$ be real vector spaces of smooth sections with specified boundary,
normalization or gauge conditions. Suppose a finite-degree polynomial
differential operator has an expansion at a real field $u_0$,
\begin{equation}\label{vtf:eq:nonlinear}
 \mathcal P(u_0+w)=\mathcal P(u_0)+Lw+
       \sum_{r=2}^d B_r(w,\ldots,w),
\end{equation}
where $L:E\to F$ is real linear and each $B_r:E^r\to F$ is real multilinear.
The finite differential order is absorbed into these maps.

\begin{theorem}[Regular nonlinear Hahn lifting]\label{vtf:thm:lifting}
Assume $L:E\to F$ is bijective, with a fixed inverse taking the prescribed
real data to smooth sections with the prescribed conditions. Given a Hahn
forcing
\[
 f=\mathcal P(u_0)+\sum_{\gamma\in S}f_\gamma t^\gamma,
 \qquad S\subset\Gamma_{>0}\text{ well ordered},
\]
there is a unique solution $u=u_0+w$, with $w$ of positive Hahn support,
to the coefficientwise equation $\mathcal P(u)=f$. The correction $w$ has
support contained in $\langle S\rangle\setminus\{0\}$.
\end{theorem}
\begin{proof}
At $\gamma\in\langle S\rangle\setminus\{0\}$ define
\begin{equation}\label{vtf:eq:recursion}
 w_\gamma=L^{-1}\left(f_\gamma-
 \sum_{r=2}^d\ \sum_{\gamma_1+\cdots+\gamma_r=\gamma}
 B_r(w_{\gamma_1},\ldots,w_{\gamma_r})\right),
\end{equation}
where all $\gamma_i>0$ and missing forcing coefficients are zero. Every
$\gamma_i<\gamma$, so the recursion is well founded. Each inner sum is
finite by the support lemma. Since $L^{-1}$ maps real coefficient data to
real coefficient solutions, it cannot introduce Hahn exponents.
This constructs a Hahn family on the stated support; substitution and
finite coefficient comparison prove the equation.

For uniqueness, the union of two proposed positive supports is well ordered.
At the least exponent where the solutions differ, expand each nonlinear
difference as a telescoping sum of multilinear terms. Every such term
contains a difference of coefficients at a strictly smaller exponent and
therefore vanishes. The difference equation is $Lz=0$. Bijectivity gives
$z=0$, a contradiction.
\end{proof}

The quadratic case around $u_0=0$ has a second proof and a sharper conclusion
(source 03).

\begin{theorem}[Exact infinitesimal quadratic solution; source 03]\label{vtf:3d:thm:quadratic}
Let $E,F$ be real vector spaces, $L_0:E\to F$ a real-linear isomorphism, and $B:E\times E\to F$ a real-bilinear map. If $f\in F((t^\Gamma))$ has positive support, the equation
\begin{equation}\label{vtf:3d:eq:quadratic}
 L_0u+B(u,u)=f
\end{equation}
has exactly one positive-support solution $u\in E((t^\Gamma))$. Its support lies in
$\langle\supp f\rangle\setminus\{0\}$. For $f\ne0$,
\[
 \vH(u)=\vH(f),\qquad
 \vH(u-L_0^{-1}f)\geq2\vH(f).
\]
For $f=0$ the unique positive-support solution is zero.
\end{theorem}
\begin{proof}
Set $a=L_0^{-1}f$ and $C=L_0^{-1}B$. Then $u=a-C(u,u)$. Expand this identity over finite rooted planar full binary trees. A leaf carries a coefficient of $a$, and each internal vertex applies $-C$ to its two incoming values. A tree with $m$ leaves has exponent equal to the sum of its $m$ positive leaf exponents. For a fixed total exponent the support lemma allows only finitely many leaf words, and for each word only finitely many planar binary trees. Therefore the tree sum is strongly defined and has the stated support. Splitting trees into the single-leaf tree and trees with a root vertex proves $u=a-C(u,u)$ exactly.

For uniqueness, suppose $u$ and $v$ are positive-support solutions and $w=u-v\ne0$. Then
\[
 w=-C(u,w)-C(w,v).
\]
If both solutions are nonzero, the right side has valuation at least
$\vH(w)+\min\{\vH(u),\vH(v)\}>\vH(w)$, impossible. If one is zero, omit its term and obtain the same contradiction. This also treats zero forcing. For nonzero forcing, $C(u,u)$ has valuation at least $2\vH(u)>\vH(u)$, so comparing leading exponents in $u+C(u,u)=a$ gives $\vH(u)=\vH(a)=\vH(f)$. The remainder estimate follows.
\end{proof}

The proof works even when the sequence $n\eta$ is not cofinal in $\Gamma$. It constructs coefficients by finite combinatorics, not by assuming that a Picard sequence converges in a topology. The theorem is a support-controlled formal implicit-function principle; it is not a claim that a real nonlinear problem has solutions for arbitrary real forcing.

With $\mathcal P(u)=L_0u+B(u,u)$, $u_0=0$ and $d=2$, existence, uniqueness and
the support statement of \cref{vtf:3d:thm:quadratic} are those of
\cref{vtf:thm:lifting}. The tree expansion is a different construction, and the
valuation estimates are additional.

\paragraph{A scalar recurrence (source 03).} The scalar model $u+u^2=t$
illustrates the tree expansion:
\[
 u=t-t^2+2t^3-5t^4+14t^5-42t^6+\cdots
   =\frac{\sqrt{1+4t}-1}{2}.
\]
Its coefficients count binary trees with alternating signs. The equality is exact in the Hahn field. The program \path{code/03-three-dimensional-fields-verify_examples.py} checks
the first twelve coefficients against the quadratic equation; the infinite existence and uniqueness statement rests on \cref{vtf:3d:thm:quadratic}, not on this finite test.

If $L$ has only a specified right inverse, the same recursion constructs a
solution with each correction in that right inverse's image. Uniqueness then
holds only with that normalization. A nontrivial cokernel creates compatibility
conditions. Gauge kernels and conservation identities in geometric PDEs
must be handled before applying the theorem.

\subsection{An exact nonlinear example in dimension \texorpdfstring{$n$}{n}}

On the ordinary flat torus $\mathbb T^n=(\R/2\pi\Z)^n$, take
\[
 \mathcal P(u)=(1-\Delta)u+u^2,
 \qquad f=\eps\cos x^1,\qquad \eps=t^\eta.
\]
The real operator $L=1-\Delta$ is bijective on smooth periodic functions:
its Fourier multiplier on mode $k\in\Z^n$ is $1+|k|^2$, and division by
it preserves rapid Fourier decay. The theorem gives an exact Hahn solution
$u=\sum_{m\ge1}u_m\eps^m$, with
\[
 u_1=L^{-1}\cos x^1,\qquad
 u_m=-L^{-1}\sum_{a+b=m}u_au_b\quad(m\ge2).
\]
The first terms are
\begin{align}\label{vtf:eq:examplePDE}
 u={}&\frac{\eps}{2}\cos x^1
 -\eps^2\left(\frac18+\frac1{40}\cos 2x^1\right)\notag\\
 &+\eps^3\left(\frac{11}{160}\cos x^1+
                   \frac1{800}\cos 3x^1\right)
 +\sum_{m\ge4}u_m\eps^m.
\end{align}
The supplied program computes six coefficients and verifies the residual
through that order with exact rational Fourier arithmetic. The all-orders
existence and uniqueness are furnished by the theorem, not by that finite
check. The solution is exact even in value groups where the sequence
$m\eta$ is not cofinal and ordinary partial-sum convergence fails.

\subsection{Tensor and constrained systems}

For tensor equations the real spaces $E,F$ may be smooth tensor sections and
$B_r$ can include contractions and covariant derivatives of a fixed real
background. For equations containing $g^{-1}$, the positive-support inverse
expansion gives an infinite Taylor nonlinearity rather than a polynomial.
The same coefficient recursion works when those multilinear terms have the
Neumann finite-word certificate. This additional certificate must be checked,
as must the invertibility of the gauge-fixed real linearization. The theorem
does not, by itself, establish global Einstein evolution, Navier--Stokes
regularity, or a singularity-removal result.

Two further instances, with their own hypotheses, are proved in Parts~II
and~III: infinitesimally forced stationary incompressible flow on the
three-torus (\cref{vtf:3d:cor:ns}) and the positive-order deformation of a real
cubic-wave solution (\cref{vtf:mink:thm:cubic}), whose real linearization is a
Cauchy problem rather than a bijection of section spaces. The same Neumann
mechanism as in \cref{vtf:3d:thm:inverse} gives the retarded inverse of a
positive-order perturbation of a hyperbolic operator
(\cref{vtf:mink:thm:retardedperturb}).

\section{Variational calculus and gauge-valued fields}\label{vtf:sec:variation}

\subsection{A scalar action and its Euler--Lagrange equation}

Let a real compact domain have an admissible metric and density, and let
$V$ be a polynomial potential with coefficients in $K$. For an admissible
scalar field set
\[
 \mathcal S[\phi]=\int^H
 \left(\tfrac12g^{ij}D_i\phi D_j\phi+V(\phi)\right)\vol_g.
\]
For a compactly supported variation $\psi$, take the coefficient of the
ordinary formal variation parameter $s$ in $\mathcal S[\phi+s\psi]$.
Leibniz's rule and Hahn Stokes give
\[
 \delta\mathcal S[\phi](\psi)
 =\int^H\bigl(-\Delta_g\phi+V'(\phi)\bigr)\psi\vol_g.
\]
The fundamental lemma remains valid in this setting: if this integral is zero
for every real smooth compactly supported $\psi$, then each coefficient of
the admissible residual density has zero integral against every real test
function, so each is zero by the real fundamental lemma. Hence stationarity
is equivalent to
\[
 -\Delta_g\phi+V'(\phi)=0.
\]
This is a precise formal-coefficient variational problem; its stationary
points and their existence are separate questions. Minimizers are not inferred
from order completeness, which $K$ lacks, and no measure on all of $K^n$ is used
(source 02).

Define the stress tensor
\[
 T_{ij}=D_i\phi D_j\phi-
 g_{ij}\left(\tfrac12g^{ab}D_a\phi D_b\phi+V(\phi)\right).
\]
Metric compatibility, Hessian symmetry and the product rule yield
\[
 \nabla^iT_{ij}=(\Delta_g\phi-V'(\phi))D_j\phi.
\]
Thus the field equation implies covariant stress conservation exactly,
including all infinitesimal orders. This identity makes no positivity claim
when the metric has indefinite signature.

\begin{remark}[Sign of the potential]\label{vtf:rem:potentialsign}
For $\eta=\diag(1,-1,-1,-1)$ one has $\Delta_\eta=\Box$, and the
Euler--Lagrange equation above reads $-\Box\phi+V'(\phi)=0$. Source 02
(Part~III) uses the Lagrangian $\tfrac12\partial_\mu\phi\,\partial^\mu\phi-V(\phi)$,
so its $V$ is $-V$ here. Its equation $\Box\phi+V'(\phi)=J$ and its stress
$\partial^\mu\phi\,\partial^\nu\phi-\eta^{\mu\nu}L_0$ are the equations above
after $V\mapsto-V$ and the addition of a source. Neither sign is a positivity
statement; the energy positivity of Part~III uses its own hypotheses
$m^2,\lambda\ge0$.
\end{remark}

\subsection{Finite-rank gauge fields}

Let $A$ be a matrix-valued admissible one-form acting on a finite-rank
$K$-bundle. Its curvature and covariant exterior derivative are
\[
 F_A=dA+A\wedge A,\qquad
 d_A\beta=d\beta+A\wedge\beta-(-1)^{\deg\beta}\beta\wedge A.
\]
Direct expansion gives $d_AF_A=0$ and $d_A^2\beta=[F_A,\beta]$ in the
graded sense. Under an admissible frame transformation $P$ with admissible
inverse,
\[
 A'=P^{-1}AP+P^{-1}dP,\qquad F_{A'}=P^{-1}F_AP.
\]
The derivative of $P^{-1}P=I$ proves the transformation law. Positive-support
transformations $P=P_0(I+H)$ with real invertible $P_0$ are a concrete
admissible class.

In any dimension, the Yang--Mills-type equation $d_A(*F_A)=0$ is therefore
well defined whenever the metric star and products are admissible. As with
the scalar action, finite-rank algebra and variational identities do not
supply an automatic global existence theorem. Surcomplex coefficients can
be introduced by $K[i]$, with conjugation and Hermitian pairings specified
separately; they do not alter the real-coordinate support arguments.

\section{Scalar derivations and scale connections}\label{vtf:sec:BM}

\subsection{Spatial and scalar derivatives are different}

Spatial differentiation in this report fixes every $K$ constant. In
particular,
\[
 D_i\omega=0.
\]
The Berarducci--Mantova derivation $\BM$ on surreal scalars instead has
$\BM\omega=1$ and is strongly additive on summable normal forms \cite{BM}. It
describes differential-field structure, not the variation of a constant scalar
field across space. All three sources insist that the two derivatives must not
be substituted for one another. The infinite-argument function and composition
calculus of omega-series \cite{BMGerms,Mantova26} is also additional structure,
not a universal replacement for coordinate calculus.

There is a compatible mixed construction when it is useful (source 01). For a
given Hahn field with smooth real coefficient functions, define
\[
 \BM F=\sum_\gamma f_\gamma\,\BM(t^\gamma).
\]
Strong additivity of the scalar derivation provides a well-ordered support
union and finite contributions. Coefficients are finite real linear
combinations of the $f_\gamma$, so they are smooth. The image may require a
larger workspace containing the derivative supports. On such admissible
expressions,
\[
 [D_i,\BM]F=0,
\]
because the $D_i$ act only on the real coefficient functions, while $\BM$ acts
only on the monomials.

An additional connection direction can be written $\nabla_{\mathrm{BM}}=\BM+B$,
with the semilinear Leibniz rule
$\nabla_{\mathrm{BM}}(aY)=\BM(a)Y+a\nabla_{\mathrm{BM}}Y$. Together with
$\nabla_i=D_i+\Gamma_i$ its mixed curvature is
\[
 [\nabla_{\mathrm{BM}},\nabla_i]
 =\BM\Gamma_i-D_iB+[B,\Gamma_i].
\]
This is a differential-module extension of the theory. Interpreting it as an
extra physical time or scale direction would be a further modelling choice; it
is not forced by the definition of a surreal-valued tensor field.

\subsection{Explicit scalar derivations and their support}

For $\Gamma\subseteq\R$ containing $1$ and the normalization $\BM\omega=1$, so
that $\BM t=-t^2$, one has
\[
 \BM(t^\gamma)=-\gamma t^{\gamma+1}\qquad(\gamma\in\Gamma\subseteq\R)
\]
(sources 02 and 03). The formula must not be used unchanged for arbitrary
surreal exponents, which can themselves have nonzero scalar derivative, and
closure of the workspace under the shift by $1$, together with strong summation
of the result, must be ensured. On general surreal exponent workspaces a
strongly acting scalar derivation needs more structure than this formula
\cite{BM,MM}.

A simpler support-preserving derivation exists at every rank (source 03). For
an additive character $\chi:\Gamma\to\R$ put
\[
 \mathscr D_\chi\Bigl(\sum_\gamma t^\gamma F_\gamma\Bigr)
 =\sum_\gamma t^\gamma\chi(\gamma)F_\gamma.
\]
It is a derivation because $\chi$ is additive. It commutes with every $D_i$ and
is $\R$-linear but not $K$-linear. For $\Gamma\subseteq\R$ and $\chi=\id$ it is
the Euler scale derivation $t\,\partial_t$, with $\mathscr D_{\id}t=t$; this is
source 02's derivation $\delta$ on $\R((t^\Q))$.

\subsection{Scale connections}

Let $a=a_i\,dx^i$ be a real smooth one-form and put
\[
 \mathcal D_i=D_i+a_i\mathscr D_\chi
\]
(source 03; source 02's coupled scale connection $\partial_\mu+a_\mu\mathscr
D_{\id}$ is the case $\chi=\id$, $n=4$). These are real-linear derivations of
$\A_\Gamma(U)$ modelling a chosen comparison rule for scale-dependent
quantities. They are not the coordinate derivatives of the scalar-constant
geometry unless the connection is trivialized.

\begin{proposition}[Scale curvature; sources 02 and 03]\label{vtf:3d:prop:scale}
Put $b_{ij}=D_ia_j-D_ja_i$. Then
\[
 [\mathcal D_i,\mathcal D_j]=b_{ij}\,\mathscr D_\chi .
\]
If $\chi\ne0$, the $\mathcal D_i$ commute exactly when $da=0$. If the $a_i$ are
allowed to be Hahn-valued, the curvature coefficient becomes
$D_ia_j-D_ja_i+a_i\mathscr D_\chi a_j-a_j\mathscr D_\chi a_i$.
\end{proposition}
\begin{proof}
Expand the commutator. The terms $[D_i,a_j\mathscr D_\chi]=(D_ia_j)\mathscr
D_\chi$ use that $\mathscr D_\chi$ commutes with $D_i$. The term
$[a_i\mathscr D_\chi,a_j\mathscr D_\chi]$ equals
$(a_i\mathscr D_\chi a_j-a_j\mathscr D_\chi a_i)\mathscr D_\chi$, because the
second-order terms $a_ia_j\mathscr D_\chi^2$ cancel; it vanishes for real $a_i$,
which $\mathscr D_\chi$ annihilates. If $\chi(\gamma)\ne0$, evaluating a
commutator on $t^\gamma$ proves the equivalence with $da=0$.
\end{proof}

Source 03 states the first two assertions for $n=3$, with the vector identities
of \cref{vtf:3d:cor:scalevector}; the commutator computation is the same in every
dimension. Source 02 states the Hahn-valued formula for $n=4$ and $\chi=\id$.
When $a=d\psi$ with real smooth $\psi$, the support-preserving algebra
automorphism
\[
 \mathcal E_\psi\Bigl(\sum_\gamma t^\gamma F_\gamma\Bigr)
 =\sum_\gamma t^\gamma\e^{\chi(\gamma)\psi}F_\gamma
\]
satisfies $\mathcal D_i=\mathcal E_\psi^{-1}D_i\mathcal E_\psi$. A locally exact
scale connection is thus a change of scale convention, and nonzero $da$ is an
integrability obstruction (source 03). Each of these is a possible additional
model, not the flat derivative of the theory already constructed: changing the
derivative rules changes the equations (sources 01 and 02).

\section{Boundaries of the coherent calculus}\label{vtf:sec:boundaries}

Two explicit obstructions show why the support hypotheses cannot be dropped, and
a third subsection records how far distributions go. All three concern the chosen
coefficient algebra $\A_\Gamma$, not every conceivable function class.

\subsection{No coherent oscillation at a negative valuation}

Let $\A_\Gamma(I;\C)=C^\infty(I,\C)((t^\Gamma))$ for a nonempty real open
interval $I$. By \eqref{vtf:mink:eq:derivativeorder}, differentiation cannot
lower the least exponent of a coherent series.

\begin{theorem}[No coherent pure oscillation at a negative valuation; sources 02 and 03]\label{vtf:mink:thm:highfrequency}
Let $\theta>0$ in $\Gamma$. The equation
\begin{equation}\label{vtf:mink:eq:highfreq}
 u'=it^{-\theta}u
\end{equation}
has no nonzero solution in $\A_\Gamma(I;\C)$. Likewise,
\begin{equation}
 u''+t^{-2\theta}u=0
\end{equation}
has no nonzero coherent solution. More generally, if a real finite-order differential operator $P_0$ preserves exponent support and $c\in K[i]$ has $v(c)<0$, then $P_0u+cu=0$ has no nonzero coherent scalar solution.
\end{theorem}
\begin{proof}
If $\alpha=\vH(u)$, then $\vH(u')\geq\alpha$, while
$\vH(it^{-\theta}u)=\alpha-\theta$. The nonzero leading coefficient at that smaller exponent cannot be matched. The second-order statement is identical. For the general statement, multiplication by a nonzero constant has exact order $\vH(cu)=v(c)+\alpha$, even though multiplication by an arbitrary smooth coefficient need not have exact order. This order is strictly smaller than any nonzero exponent in $P_0u$.
\end{proof}

The theorem is a limitation of the \emph{chosen coherent amplitude algebra}, not a proof that an oscillatory extension is impossible. One can adjoin phase symbols, work with oscillatory distributions, or construct a larger transseries-type field. Each approach must state its multiplication, derivative, summation, and real-observable rules. The usual exponential symbol by itself does not provide those rules.

Source 03 proves the special cases $Y''+t^{-2}Y=0$ and $Y'+t^{-1}Y=0$ (for
$1\in\Gamma$) by the same least-exponent argument. It adds that the formal
series $\sum_{m\ge0}(-1)^m\tau^{2m}t^{-2m}/(2m)!$ solves the second equation in
$K[[\tau]]$, but its exponents are unbounded below and it does not define a
Hahn-smooth field of ordinary time; in the fast-time chart $\tau'=\tau/t$ the
ordinary function $Y=\cos\tau'$ solves the rescaled equation. Fast dynamics
therefore need a fast-time domain, an oscillatory enrichment, or another
function class; the same issue arises in thin boundary layers and rapidly
varying material laws. The proof of \cref{vtf:mink:thm:highfrequency} uses only
\eqref{vtf:mink:eq:derivativeorder} and the exact order of multiplication by a
nonzero constant, so it applies verbatim on any real open domain in place of $I$;
source 02 states it on an interval.

\subsection{Every real time is not every surreal time}\label{vtf:mink:subsec:realtime}
The initial value problem
\begin{equation}
 y'=y^2,\qquad y(0)=t
\end{equation}
has the coherent solution
\begin{equation}\label{vtf:mink:eq:blowupexample}
 y(\tau)=\sum_{n\geq0}t^{n+1}\tau^n
         =\frac{t}{1-t\tau}
 \qquad(\tau\in\R).
\end{equation}
The formal series is well defined at every real time and differentiates coefficientwise to $y^2$. But the rational expression has no value at the surreal time $\tau=t^{-1}$, where its denominator is zero. At that argument the displayed series would contribute infinitely many copies of the same exponent and is not strongly summable either.

Thus a common-domain coefficient solution on $\R$ is not a theorem of existence at every $K$-valued time. In this example the obstruction is completely explicit. Naming the blow-up time as a surreal number does not supply a continuation through it.

\subsection{Distributional coefficients}\label{vtf:subsec:distributions}

Let $\mathcal D'(M,E)$ be the real space of distributional sections of a real
bundle. Its Hahn extension $\mathcal D'(M,E)((t^\Gamma))$ has coefficientwise
derivatives and pairings with fixed real test sections, and fixed real-linear
distributional operators extend to it without enlarging support. Multiplication
by a Hahn-smooth function is defined by convolution. Linear equations, weak
divergences and smooth-coefficient differential operators therefore have
distributional versions in the same workspace (sources 02 and 03). A general
product of two distributions is not thereby defined: the ambiguity of
$\delta_0^2$ is not removed by changing the scalar field, and nonlinear stress,
energy or convective terms formed from unrelated distributional fields need a
further product theory or regularity hypothesis.

For an infinitesimal displacement $h\in\mm_K^n$ the family
\[
 \delta_h^{\mathrm{jet}}
 =\sum_{\alpha\in\N^n}\frac{(-h)^\alpha}{\alpha!}\,\partial^\alpha\delta_0
\]
is strongly summable, and its value on a real test function $\varphi$ is
$\sum_\alpha h^\alpha\partial^\alpha\varphi(0)/\alpha!=(\ev\varphi)(h)$. It is a
precisely defined infinitesimally shifted jet distribution, not an ordinary
distribution supported at a nonreal point, and it inherits the flat-function
limitation of \cref{vtf:thm:prolongation}. Source 03 states it for $n=3$; its
support argument is that of \cref{vtf:thm:prolongation} in any dimension. The
three-dimensional identity $\Div(q\,x/r^3)=4\pi q\,\delta_0$ is in
\cref{vtf:3d:subsec:weak}.

\section{Computation, proof boundaries, and implementation}\label{vtf:sec:implementation}

\subsection{Representations and exactness}

A practical implementation should keep four objects separate: finite tensors
over an exact scalar field; sparse finite Hahn truncations; support-certified
infinite Hahn expressions; and coefficient operators on smooth or formal
real functions. A truncation is not an exact infinite series unless its tail
is separately represented or certified.

For a grid generated by positive exponents $\eta_1,\ldots,\eta_m$, a useful
normal form records a finite exponent vector, a real symbolic coefficient,
and a support certificate for allowed sums. Negative initial shifts are
stored separately. At arbitrary ordered-group rank, a cutoff $\gamma<B$
need not leave finitely many terms: for $0<\eta\ll B$, all $m\eta$ lie
below $B$. Algorithms must therefore use an explicit finite set of target
coefficients, a degree bound, or a proved left-finiteness condition. This
restriction is consistent with the repository's distinction between exact
denotation and effective support access \cite{Repo}.

Source 03 lists what a field object must carry: base domain, tensor type,
coefficient representation, exponent group, support certificate and derivative
convention. A metric additionally needs an inverse and a volume-density
certificate, and a chain needs an orientation, a parameter domain, a pullback
rule and an integral convention. For infinite fields a coefficient oracle alone
is insufficient; useful restricted formats include bounded-denominator Puiseux
series, nonnegative integer power series, and positive-monoid expressions with a
proved summability construction. A safe implementation never treats a failed
support check as zero. It reports an unsupported operation, as for inversion
across a degenerating leading coefficient, evaluation of a formal series at an
argument of nonpositive valuation, or a threshold truncation that is not finite.
Already in $\Gamma=\Q$ the support $\{1-1/m:m\ge2\}$ is well ordered with
infinitely many exponents below $1$, so ``retain all terms below a threshold''
need not give a finite representation.

\subsection{A minimal implementation architecture}

A suitable dependency order is: finite tensor algebra over an abstract real
closed field; Hahn support and strong evaluation; smooth coefficient
operators; the prolongation bridge; admissible inverses and metrics; tensor
calculus and curvature; coefficientwise chain integration; and finally
support-controlled equation solvers. A theorem about a metric inverse should
consume a unit certificate. A theorem about a substitution should consume a
joint support certificate. These are mathematical data, not comments that a
simplifier may silently discard.

Source 03 expresses the separation of responsibilities as follows; the words
``require'' and ``certified'' indicate obligations to be discharged, not checks
available for arbitrary infinite objects.

\begin{lstlisting}
Field(domain, tensor_type, exponent_group,
      coefficient_family, support_certificate)

spatial_derivative(F, i):
    return map_coefficients(F, real_partial_derivative(i))

tensor_product(F, G):
    require compatible_domains_and_exponent_groups(F, G)
    certificate = well_ordered_sum(F.support, G.support)
    return finite_at_each_exponent_convolution(F, G, certificate)

inverse_regular_operator(L0_inverse, positive_P, forcing):
    require positive_well_ordered_support(positive_P)
    return certified_neumann_expression(L0_inverse, positive_P, forcing)

integrate_form(alpha, chain):
    pulled_back = certified_pullback(alpha, chain)
    return map_coefficients(pulled_back, ordinary_chain_integral)
\end{lstlisting}

For the Lorentzian part, source 02 proposes the following order of layers.

\begin{center}\small
\begin{tabularx}{\textwidth}{@{}P{0.2\textwidth}>{\raggedright\arraybackslash}XP{0.25\textwidth}@{}}
\toprule
Layer & Objects and proofs & Main dependency\\
\midrule
Finite Lorentz algebra & Metric, cones, explicit boosts, inverse formulas, tensor actions, Hodge square & Ordered real-closed field and finite matrices\\
Electromagnetic algebra & Invariants, stress tensor, Rainich identity, dominant energy, null decomposition & Finite Lorentz algebra and sums of squares\\
Hahn tensor modules & Finite-component supports, minimum order, bounded-frame invariance & Hahn summable families and coefficient maps\\
Reduction and halo jets & Split bounded group, matrix exp/log, Taylor prolongation, chain rule & Positive-support substitution and real smooth jets\\
Differential coefficient algebra & Commuting derivations, forms, Stokes lifting, Euler--Lagrange identities & Coefficientwise maps and ordinary calculus\\
Hyperbolic lifting & Cauchy operator, retarded inverse, nonlinear recursion & Precisely packaged real existence and uniqueness theorems\\
\bottomrule
\end{tabularx}
\end{center}

For Lean, the finite part can be formulated over an abstract ordered field,
with additional square-root or real-closed hypotheses where needed. The Hahn
part should be proved separately and then instantiated through a verified
normal-form bridge. A global class of all surreals should not be forced into
a smaller universe by asserting that all its cuts or all its coefficients
form a set. Our proofs describe this layering; no new Lean implementation or
kernel verification is claimed here.

Finite Lorentz statements should be proved over an abstract real-closed field and
only then instantiated with a surreal workspace. Source 03 proposes the same
staging for three dimensions: finite vector identities over an ordered field with
the needed square roots, then tensor type and convolution extension, then an
abstract differential algebra with commuting derivations, then coefficientwise
smooth fields and the support-preserving operator lemmas. The nontrivial real
analytic inputs, namely the real Hodge theorem, the torus inverses and the Leray
projection, and the real Cauchy and Green theorems of Part~III, should be
explicit interfaces with documented hypotheses; proving a lift is not the same
as formalizing the underlying real theory. No Lean source is delivered with any
of the three manuscripts, the repository was not modified by them, and no Lean
build was performed for them. Since their pin, the repository's Lean Hahn layer
has grown (\cref{vtf:subsec:neighbours}), but it contains no counterpart of the
field-level theorems here.

\subsection{Included finite verification programs}\label{vtf:subsec:verification}

Three programs are shipped under \path{code/}, byte-identical to the delivered
ones; their recorded outputs are under \path{data/}.
\begin{description}
\item[\path{code/01-vector-tensor-fields-n-verify.py}] checks the warped-metric
curvature, the first-order metric inverse, the Lie-bracket Jacobi identity for
sample polynomial fields, $d^2=0$ for polynomial forms, Hodge-star-square signs
in dimensions $1$ through $6$ for every diagonal signature, a radial homotopy
identity, and six coefficients of the nonlinear periodic example: seven check
groups and $5{,}533$ exact assertions, recorded in
\path{data/01-vector-tensor-fields-n-verification.txt}.
\item[\path{code/02-minkowski-check_identities.py}] performs 42 named exact
checks: boosts (metric preservation, determinant, null eigenvectors, the
infinite-boost example); electromagnetic tensors (both invariants, Hodge signs,
the symmetric traceless stress and its blocks, the dominant-energy identity, the
Rainich square, null radiation stress); the Hodge square on all 16 basis wedges
in dimension four; current conservation and the four stress-divergence
components for a general potential; the null nilpotent; and the cubic-oscillator
truncation. Records: \path{data/02-minkowski-verification_log.txt} and
\path{data/02-minkowski-verification_results.json}.
\item[\path{code/03-three-dimensional-fields-verify_examples.py}] runs 11 check
groups: finite vector identities and the anisotropic cross-product rule,
grad/div/curl identities, the polynomial box, radial potentials, conformal
curvature from Christoffel symbols, scale-connection defects, elastic energy,
the Beltrami field, the quadratic recursion through degree $12$, the softened
monopole density, and the worked three-scale field. Record:
\path{data/03-three-dimensional-fields-verification_results.txt}.
\end{description}
For this merge all three were rerun on a copy of the directory with
Python~3.14.4 and SymPy~1.14.0. Every check passed, and the outputs agree with the
records apart from line endings and the Python version printed by programs 01
and 03 (the records used Python~3.13.5). Program 02 writes
\path{verification_results.json} next to itself, so it should be run on a copy.

These are exact finite symbolic checks. They do not prove real closedness, the
support lemma, arbitrary-rank summability, transfinite recursion, PDE existence,
the prolongation theorem, the operator theorems, the imported real Hodge and
Cauchy theorems, or global cohomology; those are proved or explicitly imported
in the text. Successful typesetting has no bearing on mathematical correctness.

\part{Three-dimensional fields}\label{vtf:part:3d}

\section{Scope of Part II}\label{vtf:3d:sec:scope}

This part is source 03's three-dimensional development; its dimension-free
results have been moved to Part~I and are cited from there. The analytic model
is
\begin{equation}\label{vtf:3d:eq:central}
 V(x)=\sum_{\gamma\in S}t^\gamma V_\gamma(x),\qquad
 x\in M,\quad V_\gamma\in C^\infty(M,TM),
\end{equation}
where $M$ is a real smooth three-manifold or a domain of $\R^3$ and $S$ is well
ordered. The value $V(x)$ is an actual surreal vector, not a numerical
approximation, and spatial differentiation acts on the $V_\gamma$, not on
$t^\gamma$. A nonlinear equation for $V$ becomes a support-controlled family of
real equations, and the resulting series solves it as an exact algebraic
identity. Not every pointwise surreal-valued assignment admits a common
well-ordered support of smooth coefficient fields (\cref{vtf:sec:sheaf}).
Everything in Part~I applies with $n=3$, with the Euclidean metric unless
another metric is named, and with $\partial_i=D_i$. Source 03 organizes the
layers as in \cref{vtf:sec:scope}, with polynomial and rational fields on $K^3$
as a separate layer carrying coordinate derivatives, algebraic forms and
polynomial-cell integration (\cref{vtf:sec:charts,vtf:sec:integration}).

\section{Vectors and tensors in three dimensions}\label{vtf:3d:sec:algebra}
\subsection{Euclidean vector algebra over surreal scalars}
Let $K$ be any real-closed ordered field. On $V=K^3$, define
\[
 u\cdot v=\sum_{i=1}^3u_iv_i,\qquad
 \norm{u}=\sqrt{u\cdot u},\qquad
 (u\times v)_i=\sum_{j,k=1}^3\epsilon_{ijk}u_jv_k,
\]
where $\epsilon_{ijk}$ is the alternating symbol with $\epsilon_{123}=1$.
The positive square root exists in $K$. No Archimedean assumption appears in these definitions.

\begin{theorem}[Finite vector identities]\label{vtf:3d:thm:vector}
For $u,v,w\in K^3$,
\begin{align}
 u\cdot(v\times w)&=\det[u\ v\ w],\label{vtf:3d:eq:triple}\\
 u\times(v\times w)&=v(u\cdot w)-w(u\cdot v),\\
 \norm{u\times v}^2&=\norm{u}^2\norm{v}^2-(u\cdot v)^2,\label{vtf:3d:eq:lagrange}\\
 |u\cdot v|&\leq\norm{u}\norm{v},\qquad
 \norm{u+v}\leq\norm{u}+\norm{v}.
\end{align}
If $u\ne0$, then $\norm{u}>0$, even when its norm is infinitesimal. Orthogonal projection onto $Ku$ is
$P_u(v)=u(u\cdot v)/(u\cdot u)$.
\end{theorem}
\begin{proof}
The first three identities follow by finite expansion of the alternating symbol. The right side of \eqref{vtf:3d:eq:lagrange} is a sum of three squares, hence nonnegative. Taking nonnegative square roots gives Cauchy--Schwarz. Expanding $\norm{u+v}^2$ and using that inequality gives the triangle inequality. A sum of squares in an ordered field is zero only when every term is zero; this proves positive definiteness and justifies the projection denominator.
\end{proof}

The groups $O(3,K)$ and $\SO(3,K)$ consist of matrices satisfying $Q^TQ=I$ and, in the second case, $\det Q=1$. Cross product is equivariant under $\SO(3,K)$; an orientation-reversing orthogonal map changes its sign. These are algebraic rotation groups. Parametrizing all their elements by an arbitrarily large surreal angle would
require a separate trigonometric convention. The collection's report
\path{docs/surreal/euclidean-three-space}, placed with this one, treats rotations
of $\No^3$ with finite angles, Rodrigues rotations and quaternions, and the spin
description of $\SO(3,F)$ over a real-closed field is in
\path{docs/surquaternions/surquaternions}. For pure quaternions $u,v$ one has
$uv=-u\cdot v+u\times v$, so the cross product is the vector part of their
product (source 03).

Tensors of type $(r,s)$ on $K^3$ have $3^{r+s}$ components, and all the finite
operations of \cref{vtf:sec:algebra} apply.

Repeated coordinate indices run from $1$ to $3$ and are summed. A $(1,1)$ tensor is an endomorphism; a $(0,2)$ tensor is a bilinear form. Their matrices can coincide in an orthonormal frame but have different transformation laws.

\begin{proposition}[Symmetric spectral theorem]
Every real symmetric $3\times3$ matrix over a real-closed field has an orthonormal eigenbasis and real eigenvalues in that field.
\end{proposition}
\begin{proof}
The algebraic closure is $K(i)$. For an eigenpair $Az=\lambda z$ there, $z^*z>0$ and $z^*Az$ is fixed by conjugation; hence $\lambda=(z^*Az)/(z^*z)\in K$. The nullspace of $A-\lambda I$ has a nonzero vector over $K$. Normalize it. Its orthogonal complement is $A$-invariant, so induction on dimension completes the proof.
\end{proof}
This is a pointwise theorem, not a guarantee of globally smooth eigenvector fields across repeated eigenvalues. The repository's spectral report supplies a broader treatment \cite{RepoSpectral}.

The proof is by induction on the dimension and works for symmetric $n\times n$
matrices; source 03 states it for $n=3$. In Lorentzian signature the analogous
statement fails (\eqref{vtf:mink:eq:nullnilpotent} in Part~III).

\subsection{Leading directions and multiscale rank}
\begin{proposition}[Leading geometry]\label{vtf:3d:prop:leading}
For nonzero $u\in K^3$, $v(\norm{u})=v_V(u)$. If $a=v_V(u)$ and $b=v_V(v)$, then
\[
 v_V(u\times v)\geq a+b,
\]
with equality exactly when the leading real vectors
$\st(t^{-a}u)$ and $\st(t^{-b}v)$ are linearly independent. The analogous statement for a determinant says that its leading valuation is the sum of the three column valuations exactly when their normalized leading columns are independent.
\end{proposition}
\begin{proof}
The coefficient of $t^{2a}$ in $u\cdot u$ is the positive real sum of squares of the normalized leading coordinates. Thus it cannot cancel. For the cross product, the coefficient of $t^{a+b}$ is precisely the cross product of the two leading real vectors. It vanishes exactly when they are dependent. The determinant argument is identical with three columns.
\end{proof}
For example, $u=(1,0,0)$ and $v=(1,t^\eta,0)$ are independent for $\eta>0$, but their standard parts coincide. Their area is $t^\eta$, not zero. Infinitesimal linear independence is genuine independence that ordinary reduction loses.

For a matrix $A$, valuations of its minors encode a hierarchy of rank. The ordinary rank of $\st A$ records only scale zero. A higher-order nonzero minor may restore the full exact rank. This is often the relevant distinction in nearly incompressible materials, weakly coupled modes, and anisotropic metrics.

\section{Vector calculus in three dimensions}\label{vtf:3d:sec:calculus}

On a real coordinate domain $U\subseteq\R^3$ the coefficientwise derivatives
$\partial_i=D_i$ of \cref{vtf:prop:diffalgebra} are $K$-linear commuting
derivations annihilating every constant of $K$; under a change of ordinary real
coordinates the real Jacobian factors do not enlarge scale support, so the
construction is intrinsic on real manifolds although local components are
surreal.

\subsection{Gradient, divergence, and curl in a Cartesian frame}
With the Euclidean metric,
\begin{align}
 (\grad F)^i&=\partial_iF, &
 \Div V&=\partial_iV^i, &
 (\curl V)^i&=\epsilon_{ijk}\partial_jV^k,\label{vtf:3d:eq:operators}\\
 \Delta F&=\sum_i\partial_i^2F, &
 ((V\cdot\nabla)W)^i&=V^j\partial_jW^i.
\end{align}
Here $\Delta$ is the usual spatial Laplacian, with nonpositive Fourier eigenvalues on a flat torus. Later, the Hodge Laplacian $\Delta_H$ has the opposite sign on functions.

\begin{theorem}[Vector differential identities]\label{vtf:3d:thm:identities}
For Hahn-smooth scalar fields $F,G$ and vector fields $V,W$,
\begin{align}
 \curl\grad F&=0, &\Div\curl V&=0,\\
 \curl\curl V&=\grad\Div V-\Delta V,\label{vtf:3d:eq:curlcurl}\\
 \Div(FV)&=\grad F\cdot V+F\Div V,\\
 \curl(FV)&=\grad F\times V+F\curl V,\\
 \Div(V\times W)&=W\cdot\curl V-V\cdot\curl W,\label{vtf:3d:eq:crossdiv}\\
 \curl(V\times W)&=V\Div W-W\Div V
 +(W\cdot\nabla)V-(V\cdot\nabla)W,\\
 (V\cdot\nabla)V&=\grad\bigl(\norm V^2/2\bigr)-V\times\curl V.
\end{align}
\end{theorem}
\begin{proof}
The first two identities contract a symmetric second derivative with an alternating symbol. For \eqref{vtf:3d:eq:curlcurl}, use
$\epsilon_{ijk}\epsilon_{klm}=\delta_{il}\delta_{jm}-\delta_{im}\delta_{jl}$ and commute derivatives. The product identities follow from Leibniz's rule and the same finite contractions. Every coefficient in a nonlinear expression is a finite sum, so \cref{vtf:3d:thm:closure} justifies the calculation at arbitrary rank.
\end{proof}

\subsection{Forms explain the three vector operators}

Lie derivatives, forms and covariant derivatives of arbitrary tensor type are
those of \cref{vtf:sec:fields,vtf:sec:metric} with $n=3$. The Lie derivative is
a formal infinitesimal transport operator; its existence does not prove that an
arbitrary surreal vector field has a globally defined flow.

For the Euclidean Hodge star $*$ and musical identifications,
\[
 (\grad F)^\flat=\dd F,\qquad
 (\curl V)^\flat=*\dd(V^\flat),\qquad
 \dd(\iota_V\vol)=(\Div V)\vol.
\]
Thus $\curl\grad=0$ and $\Div\curl=0$ are the degree-one and degree-two instances of $\dd^2=0$.

\subsection{Metric vector calculus}
Let $\mu_g=\sqrt{\det(g_{ij})}$ in an oriented chart. Then
\begin{align}
 (\grad_gF)^i&=g^{ij}\partial_jF,\\
 \Div_gV&=\mu_g^{-1}\partial_i(\mu_gV^i),\label{vtf:3d:eq:metricdiv}\\
 (\curl_gV)^i&=\mu_g^{-1}\epsilon^{ijk}\partial_j(g_{k\ell}V^\ell),\label{vtf:3d:eq:metriccurl}\\
 \Delta_gF&=\mu_g^{-1}\partial_i(\mu_gg^{ij}\partial_jF).
\end{align}
The raised alternating symbol in \eqref{vtf:3d:eq:metriccurl} denotes the numerical symbol, not the result of raising three indices on a tensor without its density factor. The invariant definition is
$(\curl_gV)^\flat=*_{g}\dd(V^\flat)$.

For $k$-forms in dimension three,
\[
 \delta_g=(-1)^k*_{g}\dd*_{g},\qquad
 \Delta_{H,g}=\dd\delta_g+\delta_g\dd.
\]
In particular $\delta_g(V^\flat)=-\Div_gV$ and $\Delta_{H,g}F=-\Delta_gF$.
The curved vector identity is
\begin{equation}\label{vtf:3d:eq:curvedcurl}
 \curl_g\curl_gV
 =\grad_g\Div_gV-\nabla^j\nabla_jV+\Ric^\sharp(V).
\end{equation}
It follows by commuting the two covariant derivatives in the alternating-symbol calculation. Omitting the Ricci term would be incorrect on a curved metric, surreal or real.

\subsection{Scale charts in three dimensions}

The linear scale charts $x=c+Ly$ of \cref{vtf:sec:charts} transform gradients
and divergences as stated there. Curl must be transformed with the metric and
orientation, not as three unrelated scalar derivatives; the finite identity
\begin{equation}\label{vtf:3d:eq:crossL}
 (Lu)\times(Lv)=(\det L)\,L^{-T}(u\times v)\qquad(L\in\GL_3(K))
\end{equation}
encodes the required area transformation.

Let $L=\operatorname{diag}(t^{-1},t,t^2)$ and $f(y)=y_1^2+y_2^2+y_3^2$. Then
\[
 F(x)=t^2x_1^2+t^{-2}x_2^2+t^{-4}x_3^2,
 \qquad \grad_xF=(2t^2x_1,2t^{-2}x_2,2t^{-4}x_3).
\]
A function of order one in scaled coordinates has very large derivatives in thin directions and a small derivative in the long direction. This is an exact identity, not a limit as a real parameter tends to zero.

\section{Two metric examples}\label{vtf:3d:sec:metricexamples}

\subsection{Conformal infinitesimal curvature}
Take $g=\e^{2t\psi}\delta$ on a real domain, with $\psi$ smooth. The exponential is the strong series in the positive-order field $t\psi$. In three dimensions,
\begin{align}
 \Gamma^k_{ij}&=t(\delta^k_i\partial_j\psi
              +\delta^k_j\partial_i\psi-\delta_{ij}\partial_k\psi),\\
 \Ric_{ij}&=-t\partial_i\partial_j\psi
 +t^2\partial_i\psi\partial_j\psi
 -\bigl(t\Delta\psi+t^2|\grad\psi|^2\bigr)\delta_{ij},\\
 \Scal_g&=\e^{-2t\psi}
 \bigl(-4t\Delta\psi-2t^2|\grad\psi|^2\bigr).\label{vtf:3d:eq:conformal}
\end{align}
For $\psi=|x|^2/2$ this becomes
\[
 \Scal_g=\e^{-t|x|^2}(-12t-2t^2|x|^2).
\]
The standard metric is flat, but exact curvature is nonzero at infinitesimal order. The included symbolic checks verify this example directly from \eqref{vtf:eq:christoffel}.

\subsection{Anisotropic nonregular metrics}\label{vtf:3d:subsec:anisotropic}
The constant metric
\[
 g=\operatorname{diag}(1,t^2,t^4)
\]
is positive definite, flat, and has admissible inverse
$\operatorname{diag}(1,t^{-2},t^{-4})$ and density $\mu_g=t^3$. It is not regular in the sense of \eqref{vtf:eq:regularmetric} because its standard part is degenerate. Its scalar Laplacian is
\[
 \Delta_g=\partial_1^2+t^{-2}\partial_2^2+t^{-4}\partial_3^2.
\]
For $F=t^4x_3$, the scalar field has zero standard part while
$\grad_gF=(0,0,1)$. Thus raising indices and taking standard part do not commute for this metric. The exact geometry is valid; its naive real reduction is not.

\section{Integral theorems in three dimensions}\label{vtf:3d:sec:integrals}

The integrals are those of \cref{vtf:sec:integration}: coefficientwise over
compact real chains, algebraic over polynomial chains, and by pullback over
admissibly deformed chains.

By \cref{vtf:thm:stokes}, for a compact oriented real surface $S$ and region $D$,
\begin{align}
 \int_S\curl V\cdot n\,\dd S&=\int_{\partial S}V\cdot\dd\ell,\\
 \int_D\Div V\,\dd x&=\int_{\partial D}V\cdot n\,\dd S.
\end{align}
For regular Riemannian metrics use the flux form $\iota_V\vol_g$. The normal and area density are defined from that metric with their admissible square roots, so the same statement becomes
$\int_D\Div_gV\,\vol_g=\int_{\partial D}\iota_V\vol_g$.

\subsection{Polynomial boxes at arbitrary scale}

With the algebraic moments \eqref{vtf:eq:simplex}, polynomial cells can be
integrated at any scale.

For a box with side lengths $a,b,c\in K_{>0}$,
\[
 \int_0^a\int_0^b\int_0^c
 x^py^qz^r\,\dd z\,\dd y\,\dd x
 =\frac{a^{p+1}b^{q+1}c^{r+1}}{(p+1)(q+1)(r+1)}.
\]
For $V=(x^2,y^2,z^2)$, its divergence integral and outward flux are both
\[
 abc(a+b+c).
\]
The equality remains exact when $a=t^{-1}$, $b=t$, and $c=t^2$. The coordinate volume then equals $t^2$ although one side is infinite.

\subsection{A monopole whose flux survives every scale}\label{vtf:3d:subsec:monopole}
On $K^3\setminus\{0\}$ consider the algebraic-radical field
\[
 B(x)=q\frac{x}{r^3},\qquad r=\sqrt{x\cdot x},\quad q\in K.
\]
Direct differentiation gives $\Div B=0$ away from zero. For an oriented sphere of any positive radius $R\in K$, use the scaled ordinary parametrization $x=Rn$, $n\in S^2(\R)$. Then
\[
 B\cdot n=\frac q{R^2},\qquad \dd S=R^2\dd\Omega,
 \qquad \int_{S_R}B\cdot n\,\dd S=4\pi q.
\]
Thus an infinitesimal sphere carries the same exact flux as a unit sphere; this
is the case $n=3$ of \cref{vtf:subsec:radial}. When $q$ is a nonzero infinitesimal, the flux has zero standard part but is not zero.

On the real punctured space with Hahn coefficients, no global admissible vector potential $A$ can satisfy $\curl A=B$ for $q\ne0$: Stokes on the closed sphere would force its flux to vanish. This is a topological obstruction, not an analytic failure at a small scale. The scalar potential $q/r$ still satisfies $B=-\grad(q/r)$; scalar and vector potentials solve different problems.

\subsection{Radial scalar and vector potentials}

Specializing the radial homotopy of \cref{vtf:thm:poincare} to $n=3$ gives
explicit formulas (source 03).

\begin{proposition}[Radial scalar and vector potentials]\label{vtf:3d:prop:potentials}
If $\curl V=0$, then
\[
 \phi(x)=\int_0^1x\cdot V(sx)\,\dd s
 \quad\text{satisfies}\quad \grad\phi=V.
\]
If $\Div B=0$, then
\[
 A(x)=-x\times\int_0^1sB(sx)\,\dd s
 \quad\text{satisfies}\quad\curl A=B.
\]
Both potentials have support contained in the original field's support.
\end{proposition}
\begin{proof}
For the first formula,
$\partial_i\phi=\int_0^1[V_i(sx)+s x_j\partial_iV_j(sx)]\dd s$.
The curl-free condition replaces $\partial_iV_j$ by $\partial_jV_i$, and the integrand is $\frac{\dd}{\dd s}(sV_i(sx))$. For the second formula apply $H$ to the closed two-form $\iota_B\vol$; its one-form components are exactly the displayed $A$. Equivalently a direct cross-product calculation differentiates $s^2B(sx)$. Coefficient integration does not enlarge support.
\end{proof}
These are specializations of the support-preserving Poincar\'e construction already present in \cite{RepoContours}, included to make the vector theory usable without consulting a different dimensional notation.

\subsection{The weak form of the monopole}\label{vtf:3d:subsec:weak}

On ordinary $\R^3$, the locally integrable real vector field $x/r^3$ has distributional divergence $4\pi\delta_0$. Testing against a compactly supported smooth function and applying the divergence theorem outside a small real sphere proves this: the boundary term tends to $4\pi$ times the test function's value at zero, while the ordinary divergence vanishes away from zero. Multiplication by $q\in K$ gives the exact Hahn-distribution identity
\[
 \Div\left(q\frac{x}{r^3}\right)=4\pi q\delta_0.
\]
This is the weak counterpart of the flux computation, and distinguishes an infinitesimal source from no source.

Distributional coefficients in general, and the limits of products of
distributions, are discussed in \cref{vtf:subsec:distributions}.

\section{Exact inversion on the three-torus}\label{vtf:3d:sec:torus}

\subsection{An explicit Helmholtz decomposition on the three-torus}
Let $\mathbb T^3=(\R/2\pi\Z)^3$ with flat metric. On ordinary mean-zero smooth functions, $\Delta^{-1}$ is defined by Fourier multipliers $-1/|k|^2$ for $k\in\Z^3\setminus\{0\}$. Smoothness follows because smooth Fourier coefficients decay faster than any power, and division by $|k|^2$ preserves that property. Extend this real operator coefficientwise.

For $V\in\mathfrak X_\Gamma(\mathbb T^3)$, let $h$ be its coefficientwise spatial mean and set
\[
 \phi=\Delta^{-1}\Div V,\qquad
 A=-\Delta^{-1}\curl V.
\]
Then
\begin{equation}\label{vtf:3d:eq:helmholtz}
 V=h+\grad\phi+\curl A,\qquad
 \Div A=0,\qquad \int_{\mathbb T^3}\phi=0,\quad
 \int_{\mathbb T^3}A=0.
\end{equation}
Indeed, \eqref{vtf:3d:eq:curlcurl} gives
$\curl A=-\Delta^{-1}(\grad\Div V-\Delta V)=V-h-\grad\phi$.
These conditions uniquely fix both potentials: a mean-zero harmonic scalar is zero, and a mean-zero vector with zero divergence and curl is zero. All supports stay inside $\supp V$.

If $V=t^{-2}\grad f+t\curl A_0+t^3h_0$ with real mean-zero $f$, real divergence-free mean-zero $A_0$, and real constant $h_0$, the decomposition separates an infinite irrotational part, an infinitesimal solenoidal part, and a still smaller topological harmonic part. No relative numerical smallness threshold has to be selected.

\subsection{An elliptic example with an explicit reference inverse}
On $\mathbb T^3$ take $L_0=1-\Delta:C^\infty\to C^\infty$. Its inverse has Fourier multiplier $(1+|k|^2)^{-1}$. Let $A=I+H(x,t)$ be a symmetric matrix with positive-support $H$, and set
\[
 Lu=u-\Div(A\grad u).
\]
Then $P=-\Div(H\grad\cdot)$ has positive support, so by \cref{vtf:3d:thm:inverse} every $f\in\A_\Gamma(\mathbb T^3)$ has one
exact Hahn-smooth solution. The metric or material perturbation can contain infinitely many positive exponents in an arbitrary-rank group. Each coefficient of the solution is obtained by a finite number of ordinary smooth operations and applications of $(1-\Delta)^{-1}$.

For instance, if $H=tH_1$ and $f=f_0+tf_1$, then
\begin{align*}
 u_0&=L_0^{-1}f_0,\\
 u_1&=L_0^{-1}\bigl(f_1+\Div(H_1\grad u_0)\bigr),\\
 u_n&=L_0^{-1}\Div(H_1\grad u_{n-1})\quad(n\geq2).
\end{align*}
The output is exact in $C^\infty((t^\Gamma))$, even if its coefficients grow too rapidly for convergence at any nonzero real parameter.

\subsection{Stationary incompressible flow on \texorpdfstring{$\mathbb T^3$}{the three-torus}}
Let
\[
 E=\{u\in C^\infty(\mathbb T^3,\R^3):\Div u=0,\ \langle u\rangle=0\}.
\]
Let $\Pi$ be the real Leray projection onto this space. In Fourier variables,
\[
 \widehat{\Pi f}(k)=
 \begin{cases}
 \left(I-kk^T/|k|^2\right)\widehat f(k),&k\ne0,\\
 0,&k=0.
 \end{cases}
\]
It preserves smoothness. Fix a \emph{real} viscosity $\nu>0$. The map
$L_0=-\nu\Delta:E\to E$ is an isomorphism with multiplier $1/(\nu|k|^2)$ at nonzero modes. Let
\[
 B(u,v)=\Pi((u\cdot\nabla)v).
\]
It is a bilinear map $E\times E\to E$. The convective field has zero spatial mean because periodic integration by parts and $\Div u=0$ give
$\int u_j\partial_jv_i=0$.

\begin{corollary}[Infinitesimally forced steady flow]\label{vtf:3d:cor:ns}
For every positive-support, coefficientwise mean-zero forcing
$f\in C^\infty(\mathbb T^3,\R^3)((t^\Gamma))$, there is a unique positive-support velocity $u\in E((t^\Gamma))$ and a unique positive-support mean-zero pressure $p$ solving
\begin{equation}\label{vtf:3d:eq:NS}
 -\nu\Delta u+(u\cdot\nabla)u+\grad p=f,
 \qquad\Div u=0.
\end{equation}
Both supports are contained in the positive monoid generated by $\supp f$.
\end{corollary}
\begin{proof}
Apply \cref{vtf:3d:thm:quadratic} to forcing $\Pi f$, omitting its vanishing coefficients. Then define
\[
 p=\Delta^{-1}\Div\bigl(f-(u\cdot\nabla)u\bigr)
\]
with zero mean. Helmholtz decomposition reconstructs \eqref{vtf:3d:eq:NS}. Any other positive-support solution projects to the same velocity equation and therefore has the same velocity; its mean-zero pressure then agrees. The support estimate follows from the theorem, bilinear convolution, and support-preserving real operators. If $\Pi f=0$, the velocity is zero and the pressure carries the gradient forcing.
\end{proof}

For $f=tf_1$ and $u=\sum_{n\geq1}t^nu_n$, the recursion is
\begin{align}
 u_1&=(-\nu\Delta)^{-1}\Pi f_1,\\
 u_n&=-(-\nu\Delta)^{-1}\Pi
       \sum_{a=1}^{n-1}(u_a\cdot\nabla)u_{n-a}\quad(n\geq2).
\end{align}
There are finitely many summands at each order. The theorem is stationary, periodic, and restricted to infinitesimal data around the zero solution. It makes no assertion about unrestricted time-dependent three-dimensional Navier--Stokes regularity.

\subsection{What changes when the leading operator degenerates}
If the viscosity itself is $t^\eta\nu_0$ with $\eta>0$, the leading linear operator in \eqref{vtf:3d:eq:NS} is no longer the real isomorphism used above. Rescaling may restore a useful balance, but the theorem cannot be applied unchanged. Similarly, linearization about a background with a kernel requires compatibility conditions or a finite-dimensional reduction. Noninvertible leading tensors are not a minor technical inconvenience: they change which scale controls the equation.

\section{The scale connection in three dimensions}\label{vtf:3d:sec:scaleconnection}

Let $\mathcal D_i=\partial_i+a_i\mathscr D_\chi$ be a scale connection with a real
smooth one-form $a$ (\cref{vtf:3d:prop:scale}), and define
$\grad_{\mathcal D}$, $\Div_{\mathcal D}$ and $\curl_{\mathcal D}$ by replacing
each $\partial_i$ in \eqref{vtf:3d:eq:operators} with $\mathcal D_i$.

\begin{corollary}[Modified vector identities; source 03]\label{vtf:3d:cor:scalevector}
With $b=\curl a$,
\begin{align}
 \curl_{\mathcal D}\grad_{\mathcal D}F&=b\,\mathscr D_\chi F,\label{vtf:3d:eq:scalecurl}\\
 \Div_{\mathcal D}\curl_{\mathcal D}V&=b\cdot\mathscr D_\chi V.\label{vtf:3d:eq:scalediv}
\end{align}
\end{corollary}
\begin{proof}
The first component of \eqref{vtf:3d:eq:scalecurl} is
$[\mathcal D_2,\mathcal D_3]F=b_{23}\mathscr D_\chi F$ by \cref{vtf:3d:prop:scale},
and $b_{23}$ is the first component of $\curl a$; cyclic permutations give the
rest. Expanding the divergence of the curl gives
$[\mathcal D_1,\mathcal D_2]V^3+[\mathcal D_2,\mathcal D_3]V^1+[\mathcal
D_3,\mathcal D_1]V^2$, which proves \eqref{vtf:3d:eq:scalediv}.
\end{proof}

Thus $\curl\grad=0$ and $\Div\curl=0$ hold for the scale-covariant operators
exactly when the scale connection is flat. The failure of these identities has a
precise source, the integrability obstruction $da\ne0$; it is not a property of
surreal numbers as such.

\section{Continuum applications}\label{vtf:3d:sec:physics}

\subsection{Maxwell fields and exact conservation}
Let $\tau$ be an ordinary time coordinate, distinct from the infinitesimal scale $t$. Work on a fixed real spacetime domain $I\times M$, with common support for all coefficients. In a flat spatial chart, take spatially and temporally constant material scalars $\kappa_E,\kappa_B\in K_{>0}$ and write
\begin{align}
 \partial_\tau B&=-\curl E, &\Div B&=0,\\
 \kappa_E\partial_\tau E&=\kappa_B^{-1}\curl B-J,
 &\kappa_E\Div E&=\rho.\label{vtf:3d:eq:maxwell}
\end{align}
All operations are meaningful for admissible fields. Divergence of the evolution equations shows that the magnetic constraint is propagated and that electric constraint propagation is equivalent to
\[
 \partial_\tau\rho+\Div J=0.
\]
These are exact identities over $K$.

\begin{proposition}[Poynting balance; source 03]
Every admissible solution of \eqref{vtf:3d:eq:maxwell} satisfies
\[
 \partial_\tau\left(\frac{\kappa_E}2\norm E^2
                   +\frac1{2\kappa_B}\norm B^2\right)
 +\Div\left(\frac1{\kappa_B} E\times B\right)=-J\cdot E.
\]
\end{proposition}
\begin{proof}
Dot the electric evolution equation with $E$ and the magnetic one with $B/\kappa_B$, add, and use \eqref{vtf:3d:eq:crossdiv}. Commutativity of the scalar coefficients and constancy of the material parameters are the only additional algebraic requirements.
\end{proof}
Integration over a compact region gives the exact energy balance with boundary flux. If the material parameters are real, any specified real linear Maxwell solution operator on a common time interval extends coefficientwise. If the parameters themselves have singular scale orders, the identity still holds for admissible solutions, but existence in this same function space must be re-established.

With $\kappa_E=\kappa_B=1$ the balance law is the time component of the stress
balance of \cref{vtf:mink:prop:stressdiv}, which source 02 proves in covariant
form; the two statements are kept because their hypotheses differ. The
fast-time obstruction that source 03 proves at this point is printed with source
02's general version in \cref{vtf:sec:boundaries}.

\subsection{Cauchy stress and momentum balance}
A Cauchy stress is conveniently a symmetric contravariant two-tensor $\sigma^{ij}$, with traction on a unit conormal $n_j$ given by
\[
 T^i=\sigma^{ij}n_j.
\]
Its divergence is $(\Div_g\sigma)^i=\nabla_j\sigma^{ij}$. In Cartesian coordinates, $\nabla=\partial$. Momentum balance for velocity $u$, density $\rho$, and force density $f$ is
\[
 \rho\bigl(\partial_\tau u+\nabla_u u\bigr)=\Div_g\sigma+f,
 \qquad
 \partial_\tau\rho+\Div_g(\rho u)=0
\]
when the metric is time independent. Each equation is a tensor identity in the admissible field algebra. The divergence theorem converts local linear momentum balance into the integral balance over a fixed region, with the usual transport terms when appropriate.

In Euclidean coordinates, the torque density from stress has an extra term proportional to the antisymmetric part of $\sigma$. Indeed,
\[
 \partial_j\bigl(\epsilon_{ki\ell}x_i\sigma^{\ell j}\bigr)
 =\epsilon_{kj\ell}\sigma^{\ell j}
   +\epsilon_{ki\ell}x_i\partial_j\sigma^{\ell j}.
\]
For symmetric stress the first term vanishes. Thus the usual link between stress symmetry and local angular-momentum balance is unchanged at infinitesimal or infinite coefficient scales.

\subsection{Isotropic linear elasticity as a fourth-order tensor}
For a displacement $u$, the infinitesimal strain is
\[
 e_{ij}=\tfrac12(\nabla_i u_j+\nabla_j u_i).
\]
A general elasticity tensor $C^{ijkl}$ has the minor symmetries
$C^{ijkl}=C^{jikl}=C^{ijlk}$ and, for an energy law, the major symmetry
$C^{ijkl}=C^{klij}$. It acts by $\sigma^{ij}=C^{ijkl}e_{kl}$.
In an orthonormal frame the isotropic law is
\[
 C^{ijkl}=\lambda_{\mathrm L}\delta^{ij}\delta^{kl}
 +\mu_{\mathrm L}(\delta^{ik}\delta^{jl}+\delta^{il}\delta^{jk}),
 \qquad \sigma=\lambda_{\mathrm L}(\tr e)I+2\mu_{\mathrm L} e.
\]
Here $\lambda_{\mathrm L},\mu_{\mathrm L}$ may be surreal constants or admissible coefficient fields. Write $e=e_0+\frac13(\tr e)I$, where $\tr e_0=0$. The energy density is
\begin{equation}\label{vtf:3d:eq:elastic}
 W(e)=\frac{\lambda_{\mathrm L}}2(\tr e)^2+\mu_{\mathrm L}\tr(e^2)
 =\mu_{\mathrm L}\norm{e_0}^2+\frac{3\lambda_{\mathrm L}+2\mu_{\mathrm L}}{6}(\tr e)^2.
\end{equation}
Therefore $W$ is positive definite on symmetric strains exactly when
$\mu_{\mathrm L}>0$ and $3\lambda_{\mathrm L}+2\mu_{\mathrm L}>0$. This proof uses only ordered-field positivity. A modulus can be positive infinitesimal without being zero, so very soft modes remain genuine elastic modes. Analytic coercivity estimates on a chosen real function space, however, cannot be inferred merely from this algebraic positivity.

\subsection{A Beltrami field at arbitrary amplitude scale}
On $\mathbb T^3$, let
\[
 u=a(t)(\sin z,\cos z,0),\qquad a(t)\in K.
\]
Then
\[
 \Div u=0,\qquad \curl u=u,\qquad \Delta u=-u,
 \qquad (u\cdot\nabla)u=0.
\]
Its kinetic density is $a^2/2$ and helicity density is $a^2$. For constant real viscosity it is an exact stationary solution with force $f=\nu u$ and pressure zero. Choosing $a=t^{-1}+t^2$ creates interacting energy scales
$a^2=t^{-2}+2t+t^4$, even though the velocity itself has only two scales. This simple example illustrates why nonlinear observables use convolution rather than independent coefficientwise multiplication.

\subsection{Singular fields and smoothing are different constructions}
Replacing a small real radius by a nonzero infinitesimal permits evaluating $1/r$ and $1/r^2$ there, but does not define either at $r=0$. It also does not prove finite energy or remove a singular source. This agrees with the repository's distinction between scalar enrichment and physical regularization \cite{RepoPhysics}.

For comparison, a \emph{changed} field law
\[
 B_a(x)=q\frac{x}{(r^2+a^2)^{3/2}},\qquad a>0,
\]
is defined even at zero and satisfies
\[
 \Div B_a=\frac{3qa^2}{(r^2+a^2)^{5/2}},\qquad
 \int_{S_R}B_a\cdot n\,\dd S
 =4\pi q\frac{R^3}{(R^2+a^2)^{3/2}}.
\]
It has replaced a point source by an extended density; that is extra modeling content. If $a=t$, its value and derivatives near zero live at different orders from their expansions on $r>0$. A single common-domain expansion cannot be presumed to cross this boundary layer. Scaled core charts $x=ay$ and exterior charts serve different purposes.

Surreal coefficients thus provide a rigorous language for exact scale hierarchies. They do not, by themselves, establish a new experimentally validated field theory or a resolution of gravitational or hydrodynamic singularities.

\subsection{A worked three-scale field}\label{vtf:3d:subsec:worked}

Let $t=\omega^{-1}$, $\phi=x^2y$, and
\[
 A_0=(0,0,xyz),\qquad h=(1,2,3).
\]
Define
\[
 V=t^{-1}\grad\phi+t\curl A_0+t^2h.
\]
Since
\[
 \grad\phi=(2xy,x^2,0),\qquad
 \curl A_0=(xz,-yz,0),
\]
we obtain
\[
 V=(2t^{-1}xy+txz+t^2,\quad
     t^{-1}x^2-tyz+2t^2,\quad3t^2).
\]
Direct calculation gives
\[
 \Div V=2t^{-1}y,\qquad
 \curl V=t(y,x,0),\qquad
 \Div\curl V=0.
\]
Thus the infinite irrotational component controls divergence, the infinitesimal solenoidal component controls curl, and the spatial constant component is invisible to both. This is an explicit example of why divergence and curl alone do not reconstruct a field without global or boundary information.

The dot product of the first two ordinary component fields is
\[
 (2xy,x^2,0)\cdot(xz,-yz,0)=x^2yz.
\]
Consequently the squared norm contains a scale-zero mixed term $2x^2yz$ arising from $t^{-1}\cdot t$. Taking standard part of the individual vector components is not defined here because of their negative support. It would be wrong to discard their infinite parts first and then compute the energy. One must normalize or otherwise select an appropriate finite observable before using standard part.

\part{Minkowski space}\label{vtf:part:mink}

\section{Scope and conventions of Part III}\label{vtf:mink:sec:scope}

This part is source 02's Lorentzian development. Its dimension-free material has
been moved to Part~I and is cited from there. As stated in
\cref{vtf:subsec:merge}, it is mathematics of fields on Minkowski space: physical
vocabulary names objects in the conventions fixed below, and nothing in it is a
claim about the physical world.

There are two different readings of ``Minkowski space with surreal components.'' In the first, the points themselves have surreal coordinates: an event is an element of $K^4$, and the quadratic form is defined over $K$. In the second, events remain ordinary real points, while vectors, tensors, currents, amplitudes, and coupling parameters take values in $K$. These readings share the same finite tensor algebra but not the same analytic foundations.

We develop both. The first supplies unrestricted algebraic Lorentz geometry. The second supplies a rigorous, useful field theory with differential equations, ordinary causal domains, and coefficientwise integration. A third, intermediate construction evaluates coherent fields on infinitesimal neighborhoods of real events. It permits bounded surreal coordinate changes without pretending that every real smooth function has a globally defined value at every surreal argument.

\begin{principle}
\textbf{Three carriers, not one implicit substitution.}
\[
\underbrace{\M_\R=\R^4}_{\text{real events}}
\qquad
\underbrace{\OO_K^4\supset\R^4}_{\text{finite surreal events}}
\qquad
\underbrace{\M_K=K^4}_{\text{all events in the workspace}}.
\]
The full Lorentz group acts algebraically on $K^4$. Real Lorentz transformations act on ordinary spacetime and coherent fields. Bounded surreal transformations act on finite halos by Taylor prolongation. The assertions are different and will not be conflated.
\end{principle}

\subsection{Conventions}
Throughout, $K$ is a specified set-sized real-closed ordered subfield of $\No$; our main choice is a Hahn field described in \cref{vtf:sec:scalars}. We use units $c=1$, rationalized electromagnetic units, and the signature
\begin{equation}\label{vtf:mink:eq:metric}
 \eta_{\mu\nu}=\diag(1,-1,-1,-1),\qquad
 x=(x^0,x^1,x^2,x^3)=(\tau,\mathbf x),\qquad
 q_\eta(x)=\eta_{\mu\nu}x^\mu x^\nu.
\end{equation}
Greek indices range from $0$ to $3$; Latin spatial indices range from $1$ to $3$. Repeated tensor indices are summed once up and once down. The orientation is
\begin{equation}
 \vol=\dd x^0\wedge\dd x^1\wedge\dd x^2\wedge\dd x^3,
 \qquad \epsilon_{0123}=1,\qquad\epsilon^{0123}=-1.
\end{equation}
The symbol $\epsilon_{\mu\nu\rho\sigma}$ denotes the alternating tensor; the monomial $t=t^1=\omega^{-1}$ (source 02's $\eps$) is the basic positive
infinitesimal, and $1\in\Gamma$ is assumed throughout Part~III. Ordinary time is $\tau$, never the Hahn monomial.

The wave operator and electromagnetic conventions are
\begin{align}
 \Box&=\partial_\mu\partial^\mu=\partial_\tau^2-\Delta,\label{vtf:mink:eq:box}\\
 F&=\dd A,\qquad F_{0i}=E_i,\qquad
 F_{ij}=-\epsilon_{ijk}B_k,\qquad
 \partial_\mu F^{\mu\nu}=J^\nu.\label{vtf:mink:eq:maxconventions}
\end{align}
All spacetime derivatives annihilate constant elements of $K$. They are \emph{not} the Berarducci--Mantova scalar derivation.

\subsection{Provenance and the main results}
Conway normal forms, real closedness, and the strong-summation calculus are classical inputs; the relation to derivations is discussed in \cite{Conway,Gonshor,MM,BM}. Ordered-field relativity also has an established literature; \cite{MSS} is a relevant precedent. Our ordinary wave-equation input is the real smooth Cauchy and Green theory summarized in \cite{Bar,BGP}. The conventional electromagnetic formulas can be compared with \cite{TongEM,TongRad}; we derive the signs used here.

At the pin, the repository's documentation already emphasized set-sized workspaces, distinct derivative notions, strong versus coefficientwise operations, and the fact that scalar extension does not regularize a singular spacetime \cite{Repo,RepoPhysics}. Source 02 specializes these foundations to Minkowski fields. Its substantive
outputs include the following.

\begin{center}\small
\begin{tabularx}{\textwidth}{@{}P{0.22\textwidth}>{\raggedright\arraybackslash}XP{0.2\textwidth}@{}}
\toprule
Construction or result & Precise content & Location\\
\midrule
Lorentz geometry & Cones, boosts, frames, tensors over $K$; no analytic transfer assumed & \cref{vtf:mink:sec:lorentz,vtf:mink:sec:tensors}\\
Bounded reduction & A split reduction map to the real Lorentz group; its kernel is the exponential of infinitesimal Lorentz generators & \cref{vtf:mink:sec:bounded}\\
Coherent calculus & The coefficient algebra of Part~I on real spacetime, bounded Poincar\'e action on halos & \cref{vtf:sec:prolong,vtf:mink:sec:fields}\\
Causal dynamics & Coefficientwise Cauchy solutions and retarded inverses, including positive-valuation operator perturbations & \cref{vtf:mink:sec:wave}\\
Maxwell scale detection & Positive energy, dominant energy, null stress, and exact doubling of the leading field valuation & \cref{vtf:mink:sec:maxwell,vtf:mink:sec:energy}\\
Nonlinear lifting & A unique positive-order deformation of a given real cubic-wave solution on its fixed time slab & \cref{vtf:mink:sec:nonlinear}\\
Obstructions & Infinite boosts, infinite frequencies, scalar-derivative confusion, and singular-domain limits & \cref{vtf:sec:boundaries,vtf:mink:sec:limits}\\
\bottomrule
\end{tabularx}
\end{center}
These are theorems of the explicitly defined framework, not assertions that all are new in the literature. The classical algebraic identities are not claimed as new merely because their coefficient field has changed.

\section{Algebraic Minkowski geometry over a surreal field}\label{vtf:mink:sec:lorentz}

\subsection{Vectors, cones, and the reverse triangle inequality}
Put $V=K^4$. The metric $\eta$ identifies $V$ with its algebraic dual $V^*=\Hom_K(V,K)$ by $u^\flat(v)=\inner{u}{v}$. A nonzero vector is timelike, null, or spacelike according as $q_\eta(u)>0$, $q_\eta(u)=0$, or $q_\eta(u)<0$. Future causal means
\begin{equation}
 u^0>0,\qquad q_\eta(u)\ge0.
\end{equation}
The zero vector is included when discussing the closed future cone $C^+$.

Spatial norms are algebraic:
$|\mathbf u|=(\sum_i(u^i)^2)^{1/2}$, using real closedness. The identity
\begin{equation}
 |\mathbf a|^2|\mathbf b|^2-(\mathbf a\cdot\mathbf b)^2
 =\sum_{i<j}(a_i b_j-a_j b_i)^2
\end{equation}
proves Cauchy--Schwarz over $K$ without compactness or limits. Hence future causality is equivalent to $u^0\ge|\mathbf u|$ with $u^0\ge0$, and $C^+$ is closed under addition and nonnegative scalar multiplication.

\begin{proposition}[Reverse Cauchy--Schwarz]\label{vtf:mink:thm:reverse}
For future timelike vectors $u,w$,
\begin{equation}
 \inner{u}{w}\ge\sqrt{q_\eta(u)q_\eta(w)}.
\end{equation}
Equality holds precisely when $u,w$ are positive scalar multiples. Moreover,
\begin{equation}
 \sqrt{q_\eta(u+w)}\ge\sqrt{q_\eta(u)}+\sqrt{q_\eta(w)}.
\end{equation}
\end{proposition}
\begin{proof}
Divide $u$ by $\sqrt{q_\eta(u)}$ and write the resulting unit vector as $(a,\mathbf b)$, with $a>0$ and $a^2-|\mathbf b|^2=1$. The explicit Lorentz transformation in \cref{vtf:mink:eq:restboost} below takes this vector to $e_0$ after inversion and preserves future causality. In that frame the desired inequality is $w^0\ge\sqrt{(w^0)^2-|\mathbf w|^2}$, with equality exactly when $\mathbf w=0$. Squaring the second assertion and expanding $q_\eta(u+w)$ gives the first assertion. All square roots are of nonnegative elements.
\end{proof}
This argument is algebraic; ``frame'' here means a basis change, not an appeal to connectedness of a non-Archimedean topological group.

\subsection{Lorentz transformations and explicit rest frames}
Define
\begin{equation}
 O(1,3;K)=\{\Lambda\in\GL_4(K):\Lambda^T\eta\Lambda=\eta\},
\end{equation}
and
\begin{equation}\label{vtf:mink:eq:G}
 \mathcal G=\SO^+(1,3;K)=
 \{\Lambda\in O(1,3;K):\det\Lambda=1,\ \Lambda^0{}_0>0\}.
\end{equation}
The superscript $+$ is an algebraic time-orientation condition, not the assertion that $\mathcal G$ is a topological identity component.

For a unit future timelike $u=(a,\mathbf b)$ define
\begin{equation}\label{vtf:mink:eq:restboost}
 B_u=\begin{pmatrix}
 a&\mathbf b^T\\
 \mathbf b&I_3+\dfrac{\mathbf b\mathbf b^T}{a+1}
 \end{pmatrix}.
\end{equation}
Since $a\ge1$, the denominator is nonzero. Multiplication, using $|\mathbf b|^2=a^2-1$, proves $B_u^T\eta B_u=\eta$, $B_ue_0=u$, and
$B_u^{-1}=B_{(a,-\mathbf b)}$. Its determinant is $1$: on the spatial orthogonal complement of $\mathbf b$ it is the identity, and the remaining two-dimensional block has determinant $a^2-|\mathbf b|^2=1$; $\mathbf b=0$ is immediate.

To see that it preserves the future cone, its time component on $w\in C^+$ satisfies
$a w^0+\mathbf b\cdot\mathbf w\ge(a-|\mathbf b|)w^0\ge0$, strictly for nonzero $w$. Spatial rotations preserve that cone too. More generally, if $\Lambda\in\mathcal G$, its first column is future unit timelike; $B_{\Lambda e_0}^{-1}\Lambda$ fixes $e_0$ and has an $\SO_3(K)$ spatial block. Thus every element of $\mathcal G$ preserves $C^+$ and decomposes into a boost and a rotation.

\subsection{Boosts without a globally available rapidity}
A useful parametrization uses any $\lambda\in K_{>0}$:
\begin{equation}\label{vtf:mink:eq:boostlambda}
 B(\lambda)=
 \begin{pmatrix}
 (\lambda+\lambda^{-1})/2&(\lambda-\lambda^{-1})/2&0&0\\
 (\lambda-\lambda^{-1})/2&(\lambda+\lambda^{-1})/2&0&0\\
 0&0&1&0\\0&0&0&1
 \end{pmatrix}.
\end{equation}
This is an active boost. The passive coordinate boost uses its inverse. It satisfies
\begin{equation}
 B(\lambda)B(\mu)=B(\lambda\mu),\qquad
 B(\lambda)k_+=\lambda k_+,\quad
 B(\lambda)k_-=\lambda^{-1}k_-,
 \quad k_\pm=(1,\pm1,0,0).
\end{equation}
No logarithm of $\lambda$ is required inside the fixed Hahn workspace. If a suitable exponential extension has been selected, one may set $\lambda=e^r$ and recover rapidity.

For $\lambda=t^{-1}$,
\begin{equation}\label{vtf:mink:eq:infiniteboost}
 \gamma=\frac{t^{-1}+t}{2},\qquad
 \beta=\frac{1-t^2}{1+t^2},\qquad
 1-\beta^2=\frac{4t^2}{(1+t^2)^2}>0.
\end{equation}
The velocity remains strictly subluminal, yet the Lorentz factor is unlimited. This is an exact algebraic fact, not a physical claim about attainable observers.

\subsection{Four-momentum and the mass shell}
For a positive rest mass $m\in K$ and future unit velocity $U$, define the four-momentum $p=mU$. Then
\begin{equation}
 q_\eta(p)=m^2,\qquad p^0=\sqrt{m^2+|\mathbf p|^2},
 \qquad \mathbf v=\frac{\mathbf p}{p^0},\qquad |\mathbf v|<1.
\end{equation}
A future null momentum has $|\mathbf v|=1$. The energy measured by a unit future observer $u$ is $E_u=\eta(p,u)$; reverse Cauchy--Schwarz gives $E_u\geq m$, with equality exactly in the rest frame. These are algebraic kinematic relations, valid for finite, infinitesimal, and unlimited positive $m$.

Collinear velocity addition also remains exact. For $|v|,|w|<1$, put
\begin{equation}
 v\oplus w=\frac{v+w}{1+vw},\qquad
 1-(v\oplus w)^2=\frac{(1-v^2)(1-w^2)}{(1+vw)^2}>0.
\end{equation}
The denominator is positive because $|vw|<1$. No limiting process is needed to prove that a composition remains subluminal.

In reference units, take $m=t$ and boost its rest momentum by $B(t^{-1})$. The resulting momentum is
\begin{equation}
 p=\left(\frac{1+t^2}{2},\frac{1-t^2}{2},0,0\right),
 \qquad q_\eta(p)=t^2.
\end{equation}
It is timelike and finite but has the nonzero null shadow $(1/2,1/2,0,0)$. This illustrates how a real massless approximation can coexist with an exact positive infinitesimal mass, without identifying the two exact causal types.

\subsection{Affine events and causal order}
The Poincar\'e group $K^4\rtimes\mathcal G$ acts on $\M_K$ by $x\mapsto\Lambda x+a$. Define $x\preceq y$ if $y-x\in C^+$. Cone closure gives transitivity, and $C^+\cap(-C^+)=\{0\}$ gives antisymmetry. Thus this is a partial order invariant under the time-oriented Poincar\'e group.

If $y-x$ is timelike, the proper length of the affine segment is $\sqrt{q_\eta(y-x)}$. That is an algebraic invariant. It does not by itself define integrals along arbitrary $K$-parameterized curves or prove any causal compactness theorem for the order topology of $K^4$.

\section{Tensor algebra, exterior algebra, and Hodge duality}\label{vtf:mink:sec:tensors}

\subsection{Tensor types and their transformations}
A tensor of type $(r,s)$ is an element of
\begin{equation}
 T^r{}_s(V)=V^{\otimes_K r}\otimes_K(V^*)^{\otimes_K s},
 \qquad \dim_K T^r{}_s(V)=4^{r+s}.
\end{equation}
Tensor products, permutation symmetries, contractions, traces and metric raising
and lowering are finite operations and require no summability. Under
$x'=\Lambda x$, upper indices transform with $\Lambda$ and lower indices with
$\Lambda^{-1}$; this is \eqref{vtf:eq:transform} with $P=\Lambda^{-1}$.

The dual map is inverse transpose in column-vector notation. Confusing a covector with a vector merely because both have four components is especially dangerous when some entries are infinite.

A rank-two covariant tensor has the direct decomposition
\begin{equation}
 T_{\mu\nu}=T_{(\mu\nu)}^{\mathrm{tf}}
 +\frac14\eta_{\mu\nu}\eta^{\alpha\beta}T_{\alpha\beta}
 +T_{[\mu\nu]},
\end{equation}
with dimensions $9+1+6=16$. The formula remains valid over $K$.

\subsection{Indefinite self-adjointness is not Euclidean spectral theory}
The adjoint of an endomorphism is $A^\dagger=\eta^{-1}A^T\eta$. A Lorentz-self-adjoint operator need not be diagonalizable. For a nonzero null vector $k$, let
\begin{equation}\label{vtf:mink:eq:nullnilpotent}
 N(v)=k\inner{k}{v}.
\end{equation}
Then $N^\dagger=N$, $N\ne0$, and $N^2=0$. Therefore no theorem saying that every symmetric tensor has an orthonormal eigenbasis can be imported from positive-definite geometry. This counterexample already exists over $\R$ and persists over the surreals.

For a unit timelike observer $u$, however, the orthogonal space $u^\perp$ is three-dimensional and $-\eta$ is positive definite on it. All ordinary finite-dimensional Euclidean algebra works there over the real-closed field $K$. The projector
\begin{equation}
 h^\mu{}_\nu=\delta^\mu{}_\nu-u^\mu u_\nu
\end{equation}
gives the observer's time--space decomposition
$v=(u\cdot v)u+h(v)$.

\subsection{Forms and the Lorentzian Hodge operator}
The exterior powers have dimensions $1,4,6,4,1$. The metric gives the pairing
\begin{equation}
 \langle\alpha,\beta\rangle
 =\frac1{p!}\alpha_{\mu_1\cdots\mu_p}
                  \beta^{\mu_1\cdots\mu_p}
 \quad\text{on }\Lambda^pV^*.
\end{equation}
Define $*$ uniquely by
\begin{equation}\label{vtf:mink:eq:hodge}
 \alpha\wedge*\beta=\langle\alpha,\beta\rangle\vol.
\end{equation}
For a two-form,
\begin{equation}
 (* F)_{\mu\nu}
 =\frac12\epsilon_{\mu\nu\rho\sigma}F^{\rho\sigma}.
\end{equation}
On a basis wedge indexed by $I$, $*$ sends it to the complementary wedge, multiplied by the product of metric signs on $I$ and the permutation sign of $(I,I^c)$. Applying this rule twice proves
\begin{equation}\label{vtf:mink:eq:starsquare}
 *^2=(-1)^{p(4-p)+3}\id.
\end{equation}
Thus $*^2=-1$ on two-forms, and $*^2=1$ on one- and three-forms; this is
\eqref{vtf:eq:star2} with $n=4$ and $q=3$ (\cref{vtf:rem:specializations}). The construction commutes with orientation-preserving Lorentz transformations.

There is no direct four-dimensional replacement for the three-dimensional cross product of \emph{two} vectors with vector output: $u\wedge v$ has six components. The natural output is a bivector. The Hodge dual of a wedge of \emph{three} covectors does have one-form degree; after raising its index it is orthogonal to all three original vectors. An observer $u$ and an orientation of $u^\perp$ recover an ordinary spatial cross product within that three-space.

\subsection{Surcomplex self-dual fields}
After extending coefficients to $K[i]$, set
\begin{equation}
 F_+=\frac12(F-i* F),\qquad
 F_-=\frac12(F+i* F).
\end{equation}
Then $* F_+=iF_+$, $* F_-=-iF_-$, and each eigenspace has dimension three over $K[i]$. For a real $K$-valued field, conjugation exchanges $F_+$ and $F_-$. This is a representation-theoretic decomposition, not an invocation of surcomplex contour integration.

\section{Bounded Lorentz transformations and real shadows}\label{vtf:mink:sec:bounded}

\subsection{The reduction group}
Write $\mathcal G=\SO^+(1,3;K)$ for the group of \cref{vtf:mink:sec:lorentz}, and define
\begin{equation}
 \mathcal G_{\OO_K}=\mathcal G\cap\operatorname{Mat}_4(\OO_K),\qquad
 \mathcal G_1=\mathcal G\cap\bigl(I+\operatorname{Mat}_4(\mm_K)\bigr).
\end{equation}
The word \emph{bounded} means entrywise finite; it does not prescribe one common real bound for every matrix in the group. The inverse formula $\Lambda^{-1}=\eta^{-1}\Lambda^{\mathsf T}\eta$ shows that $\mathcal G_{\OO_K}$ is a group.

\begin{theorem}[Bounded Lorentz reduction]\label{vtf:mink:thm:reduction}
Entrywise standard part gives a split exact sequence
\begin{equation}\label{vtf:mink:eq:split}
 1\longrightarrow \mathcal G_1\longrightarrow \mathcal G_{\OO_K}
 \mathop{\longrightarrow}^{\st}\SO^+(1,3;\R)
 \longrightarrow1.
\end{equation}
Let
\begin{equation}
 \mathfrak g_1=
 \{X\in\operatorname{Mat}_4(\mm_K):X^{\mathsf T}\eta+\eta X=0\}.
\end{equation}
The strongly summable matrix series $\exp$ and $\log$ define mutually inverse bijections
\begin{equation}
 \exp:\mathfrak g_1\longrightarrow \mathcal G_1,
 \qquad \log:\mathcal G_1\longrightarrow\mathfrak g_1.
\end{equation}
Consequently every bounded Lorentz matrix has a unique expression
\begin{equation}\label{vtf:mink:eq:lorentzfactor}
 \Lambda=\exp(X)L_0,\qquad
 L_0=\st\Lambda\in\SO^+(1,3;\R),\quad X\in\mathfrak g_1.
\end{equation}
\end{theorem}
\begin{proof}
Reduction preserves the defining polynomial equations, including determinant one. The first column of a Lorentz matrix satisfies
$(\Lambda^0{}_0)^2=1+\sum_i(\Lambda^i{}_0)^2$; time orientation implies $\Lambda^0{}_0\geq1$. Hence its standard part is at least one and the reduced matrix is orthochronous. The real constant matrices give a section, and the kernel is exactly $\mathcal G_1$.

For an infinitesimal matrix, the finite union of the supports of its entries lies in $\Gamma_{>0}$. The Neumann lemma makes every coefficient of each matrix series a finite sum, even when the value group has arbitrary rank. Formal one-variable exponential and logarithm identities therefore apply.

Put $A^\dagger=\eta^{-1}A^{\mathsf T}\eta$. If $X^\dagger=-X$, then
$(\exp X)^\dagger\exp X=\exp(-X)\exp X=I$. Thus $\det(\exp X)^2=1$. Its determinant has standard part one, so it equals one, rather than minus one. Its $00$ entry also has standard part one, so it is orthochronous. Conversely, if $\Lambda=I+H\in \mathcal G_1$, then $\Lambda^\dagger=\Lambda^{-1}$, and the formal identities give
\[
 (\log\Lambda)^\dagger=\log(\Lambda^\dagger)
 =\log(\Lambda^{-1})=-\log\Lambda.
\]
The identities $\exp\log(I+H)=I+H$ and $\log\exp X=X$ prove the bijection. Finally $\Lambda L_0^{-1}\in \mathcal G_1$, which gives \cref{vtf:mink:eq:lorentzfactor} and its uniqueness.
\end{proof}

As groups, \cref{vtf:mink:eq:split} gives $\mathcal G_{\OO_K}\simeq \mathcal G_1\rtimes\SO^+(1,3;\R)$, with the real group acting by conjugation. The logarithmic coordinate on $\mathcal G_1$ is \emph{not} additive: multiplying exponentials uses the Baker--Campbell--Hausdorff series, again with positive support. The theorem uses only an infinitesimal exponential; it does not require a globally defined exponential on arbitrary Hahn matrices.

\subsection{Tensor valuations and bounded-frame invariance}
For a nonzero finite-component tensor, define its component valuation by
\begin{equation}\label{vtf:mink:eq:tensorval}
 v_V(T)=\min\{v(T^{\mu_1\ldots\mu_r}{}_{\nu_1\ldots\nu_s})\},
 \qquad v_V(0)=+\infty.
\end{equation}
This minimum refers to a specified frame, not to the scalar contractions of $T$.

\begin{proposition}[Frame invariance and reduction]\label{vtf:mink:prop:frame}
If $\Lambda\in \mathcal G_{\OO_K}$, then $v_V(\Lambda\cdot T)=v_V(T)$ for every tensor type. For tensors with finite components,
\begin{equation}
 \st(\Lambda\cdot T)=(\st\Lambda)\cdot\st(T).
\end{equation}
\end{proposition}
\begin{proof}
Every coefficient in the tensor transformation matrix is a finite product of entries of $\Lambda$ and $\Lambda^{-1}$, and therefore has nonnegative valuation. The finite transformation sums cannot lower the minimum component valuation. Applying the same argument to the inverse gives the opposite inequality. The reduction identity follows from multiplicativity and additivity of standard part on $\OO_K$.
\end{proof}

This is the invariance recorded in \cref{vtf:sec:algebra}, for the bounded
Lorentz group.

\begin{example}[An infinite boost reveals an infinitesimal vector]\label{vtf:mink:ex:boostshadow}
Let $k_+=(1,1,0,0)$ and let $B(t^{-1})$ be the null-eigenvalue boost of \cref{vtf:mink:eq:boostlambda}. Then
\begin{equation}
 V=t k_+,\qquad \st V=0,\qquad
 B(t^{-1})V=k_+,\qquad \st(B(t^{-1})V)=k_+.
\end{equation}
Both the input and output vectors have finite components. No real linear transformation can send the zero shadow to this nonzero shadow. Thus there is no extension of the standard-part equivariance statement to the full Lorentz group. The same boost sends the unit rest vector to a vector with unlimited components.
\end{example}

\subsection{Near-null geometry and a quantitative stability criterion}
The causal type of a finite vector need not be visible in its real shadow. The three vectors
\begin{equation}
 (1,1-t,0,0),\qquad(1,1,0,0),\qquad(1,1+t,0,0)
\end{equation}
are respectively timelike, null, and spacelike, but all have the same standard part. Their squared lengths are $2t-t^2$, $0$, and $-2t-t^2$.

For $V\ne0$, put $\alpha=v_V(V)$. If $q_\eta(V)\ne0$, define its \emph{null depth}
\begin{equation}
 d(V)=v(q_\eta(V))-2\alpha\in\Gamma_{\geq0};
\end{equation}
put $d(V)=+\infty$ for a null vector. A positive depth records cancellation of the leading quadratic form, not smallness of every component.

\begin{theorem}[Causal stability beyond the null depth]\label{vtf:mink:thm:causalstability}
For $q_\eta(V)\ne0$, let $d=d(V)$. If
\begin{equation}\label{vtf:mink:eq:stabilitybound}
 v_V(W)>\alpha+d,
\end{equation}
then $q_\eta(V+W)$ has the same leading exponent, leading coefficient, and sign as $q_\eta(V)$. Both $d(V)$ and condition \cref{vtf:mink:eq:stabilitybound} are invariant under bounded Lorentz changes of frame.
\end{theorem}
\begin{proof}
Let $\beta=v_V(W)$. The bilinear expression obeys
$v(\eta(V,W))\geq\alpha+\beta$ and $v(q_\eta(W))\geq2\beta$. Since $d\geq0$ and $\beta>\alpha+d$, both bounds are strictly greater than $2\alpha+d=v(q_\eta(V))$. In
$q_\eta(V+W)=q_\eta(V)+2\eta(V,W)+q_\eta(W)$, the leading term therefore survives. Frame invariance follows from \cref{vtf:mink:prop:frame} and the exact invariance of $q$.
\end{proof}
The criterion is sufficient, not necessary. A much larger perturbation tangent to an appropriate level set can preserve the sign. Its significance is that it works without an Archimedean comparison of scales. For an exactly null vector there is no finite depth protecting the zero value against arbitrary transverse perturbations.

\section{Fields on real Minkowski spacetime}\label{vtf:mink:sec:fields}

On an open $U\subset\M_\R$ the coefficient algebra is $\A_\Gamma(U)$ of
\cref{vtf:sec:sheaf}; a tensor field of type $(r,s)$ is an element of
$\A_\Gamma(U)^{4^{r+s}}$ with the transition laws of the tensor representation,
forms are $\Omega^p_\Gamma(U)=\Omega^p(U;\R)((t^\Gamma))$, and $\partial_\mu$ is
the coefficientwise derivative. These are modules over a coefficient algebra,
not finite-dimensional vector spaces over $K$, and they contain much more than
$K\otimes_\R C^\infty(U)$ (\cref{vtf:3d:rem:nottensor}). Everything in
\cref{vtf:sec:sheaf,vtf:sec:prolong,vtf:sec:fields,vtf:sec:integration,vtf:sec:cohomology}
applies with $n=4$. A vector with nonreal entries is not literally an ordinary
real tangent vector to $U$. In Minkowski Cartesian coordinates the flat
connection is $\nabla_\mu=\partial_\mu$, and curvature and torsion vanish
identically, over any coefficient workspace (source 02).

\subsection{Moving frames and invertibility}
An admissible moving frame is a matrix $E\in\GL_4(\A_\Gamma(U))$. Pointwise invertibility alone does not certify that its inverse has one coherent common-domain expansion. A useful sufficient condition is
\begin{equation}
 E=E_0+H,\qquad E_0\in C^\infty(U,\GL_4(\R)),\qquad
 H\in\operatorname{Mat}_4(\A_{>0}(U)).
\end{equation}
Then
\begin{equation}
 E^{-1}=\sum_{n\geq0}(-E_0^{-1}H)^n E_0^{-1}
\end{equation}
is coherent by the positive-support lemma. In the moving frame the connection one-form is $\omega=E^{-1}\dd E$. Direct differentiation of $E^{-1}E=I$ proves
\begin{equation}
 \dd\omega+\omega\wedge\omega=0.
\end{equation}
Thus a varying frame need not have zero connection coefficients, although the underlying connection is flat. If $E^{\mathsf T}\eta E=\eta$, differentiating this identity proves $\omega^{\mathsf T}\eta+\eta\omega=0$. These formulas allow accelerated tetrads and nonconstant observer fields while retaining the metric and support conditions explicitly.

\subsection{Bounded Poincar\'e transformations of coherent fields}
Let
\begin{equation}
 g(z)=\Lambda z+b,\qquad \Lambda\in \mathcal G_{\OO_K},\quad b\in\OO_K^4,
\end{equation}
and let $g_0(x)=(\st\Lambda)x+\st b$. For real $x\in g_0^{-1}U$, define
\begin{equation}\label{vtf:mink:eq:boundedpullback}
 (g^*f)(x)=\ev f(g(x)).
\end{equation}
The positive-support displacement
$(\Lambda-\st\Lambda)x+(b-\st b)$ has coefficient functions polynomial in $x$ and an exponent support independent of $x$. The proof of \cref{vtf:thm:prolongation} therefore shows that \cref{vtf:mink:eq:boundedpullback} is a coherent smooth series on $g_0^{-1}U$.

The formal chain rule proves
\begin{equation}
 \partial_\mu(g^*f)=\Lambda^\nu{}_\mu\,g^*(\partial_\nu f),
 \qquad \Box(g^*f)=g^*(\Box f).
\end{equation}
Composition obeys $(g\circ h)^*=h^*g^*$. The derivative statement extends to covariant and contravariant tensors with their respective transformation matrices; thus the flat tensor calculus and the field equations below are covariant under this bounded halo action. For finite fields, reduction commutes with the action:
\begin{equation}
 \st(g^*f)=g_0^*(\st f).
\end{equation}

This is the pullback \eqref{vtf:eq:nonrealpull} for the affine map $g$: its real
part $g_0$ is a real Poincar\'e transformation and its correction has positive
support, so \cref{vtf:thm:prolongation} and \cref{vtf:prop:diffeo} apply.

\subsection{What remains valid on all of \texorpdfstring{$K^4$}{K4}?}
Polynomial functions $K[x^0,x^1,x^2,x^3]$ have a globally defined full-Poincar\'e action. Their formal partial derivatives obey the ordinary chain rule at arbitrary $K$-valued points. Rational functions admit the same operations on their explicitly stated nonvanishing denominator domains. One can develop additional analytic classes on valuation balls, but their summability domains must be specified separately.

Consequently, the finite algebraic tensor theory has full $K$-Lorentz covariance, while the coherent smooth theory just constructed has real-spacetime and bounded-halo covariance. An infinite Lorentz transformation need not preserve the chosen coefficient function algebra. \Cref{vtf:mink:sec:limits} proves an explicit obstruction rather than hiding this boundary in notation.

\subsection{Codifferentiation and the de Rham wave operator}
For a $p$-form with $p\geq1$, define the codifferential in Cartesian components by
\begin{equation}
 (\delta\alpha)_{\mu_2\cdots\mu_p}
 =-\partial^\rho\alpha_{\rho\mu_2\cdots\mu_p},
\end{equation}
and put $\delta=0$ on scalars. With our four-dimensional Lorentzian Hodge convention, this is $\delta=*\dd*$. The identity follows on each basis wedge from the metric signs and the complementary permutation. In particular $\delta^2=0$. Expanding the alternating derivative gives, in every degree,
\begin{equation}
 (\dd\delta+\delta\dd)\alpha=-\Box\alpha,
\end{equation}
where the right side acts componentwise. The mixed second derivatives cancel and only $-\partial^\rho\partial_\rho$ on each component remains. The resulting de Rham operator is hyperbolic, rather than the positive elliptic operator associated with Euclidean Hodge theory. This is another place where the metric signature matters.

\section{Electromagnetic fields over surreal coefficients}\label{vtf:mink:sec:maxwell}

\subsection{The two-form, sources, and the three-plus-one decomposition}
A potential is a coherent one-form $A=A_\mu\dd x^\mu$. With $A_\mu=(\phi,-\mathbf A)$, its field strength $F=\dd A$ is
\begin{equation}\label{vtf:mink:eq:Fmatrix}
 (F_{\mu\nu})=
 \begin{pmatrix}
 0&E_1&E_2&E_3\\
 -E_1&0&-B_3&B_2\\
 -E_2&B_3&0&-B_1\\
 -E_3&-B_2&B_1&0
 \end{pmatrix},
 \qquad \mathbf E=-\grad\phi-\partial_\tau\mathbf A,
 \quad \mathbf B=\curl\mathbf A.
\end{equation}
In particular $F^{0i}=-E_i$, $F^{i0}=E_i$, and $F^{ij}=F_{ij}$. With $J^\mu=(\rho,\mathbf j)$, the equations
\begin{equation}\label{vtf:mink:eq:maxwellcov}
 \partial_{[\lambda}F_{\mu\nu]}=0,
 \qquad \partial_\mu F^{\mu\nu}=J^\nu
\end{equation}
are exactly
\begin{align}
 \Div\mathbf E&=\rho,&
 \curl\mathbf B-\partial_\tau\mathbf E&=\mathbf j,\label{vtf:mink:eq:maxwell3a}\\
 \Div\mathbf B&=0,&
 \partial_\tau\mathbf B+\curl\mathbf E&=0.\label{vtf:mink:eq:maxwell3b}
\end{align}
Every operation here is a finite polynomial differential operation in the coherent algebra. The equations can therefore be interpreted either as equations of coherent fields on real spacetime or, for polynomial fields, as equations on all of $K^4$.

The differential-form version, with the Hodge convention \cref{vtf:mink:eq:hodge}, is
\begin{equation}\label{vtf:mink:eq:maxwellforms}
 \dd F=0,\qquad \dd* F=-* J^\flat,
 \qquad J^\flat=J_\mu\dd x^\mu.
\end{equation}
The minus sign is convention dependent and is intentional. For example the spatial $123$ component of $\dd* F$ is $-\Div\mathbf E$, whereas the corresponding component of $* J^\flat$ is $\rho$. Equivalently, the excitation form $-* F$ has exterior derivative $* J^\flat$.

\subsection{Gauge symmetry and constraints}
For any coherent scalar $\chi$,
\begin{equation}
 A\longmapsto A+\dd\chi
\end{equation}
leaves $F$ unchanged. Conversely, on a star-shaped real domain, potentials for the same field differ by such a gauge term, by the coefficientwise Poincar\'e lemma. Antisymmetry and commuting derivatives imply current conservation:
\begin{equation}
 \partial_\nu J^\nu
 =\partial_\nu\partial_\mu F^{\mu\nu}=0.
\end{equation}
The Lorenz condition, written with a \emph{z}, is
\begin{equation}
 \partial_\mu A^\mu=0.
\end{equation}
In this gauge the field equation reduces to $\Box A^\nu=J^\nu$. Reaching the gauge requires solving $\Box\chi=-\partial_\mu A^\mu$ with suitable initial or support conditions; it is not a purely algebraic consequence of having a field extension.

The constraints in \cref{vtf:mink:eq:maxwell3a,vtf:mink:eq:maxwell3b} propagate under the evolution equations. Indeed,
\begin{align}
 \partial_\tau(\Div\mathbf E-\rho)
 &=-\Div\mathbf j-\partial_\tau\rho=0,\\
 \partial_\tau\Div\mathbf B&=-\Div\curl\mathbf E=0.
\end{align}
Thus compatible coherent initial data and a conserved coherent source give coefficientwise constraint propagation. Existence and uniqueness are obtained by lifting the ordinary Maxwell Cauchy problem, or by an appropriately constrained wave equation as in \cref{vtf:mink:sec:wave}.

\subsection{Invariants, duality, and observer dependence}
The Hodge dual acts on the three-vectors by
\begin{equation}
 *:(\mathbf E,\mathbf B)\longmapsto(-\mathbf B,\mathbf E).
\end{equation}
Direct contraction gives the two Lorentz invariants
\begin{equation}\label{vtf:mink:eq:EMinvariants}
 I_1=F_{\mu\nu}F^{\mu\nu}=2(\mathbf B^2-\mathbf E^2),
 \qquad
 I_2=F_{\mu\nu}(* F)^{\mu\nu}=4\mathbf E\cdot\mathbf B.
\end{equation}
The sign of $I_2$ changes under a different orientation convention. It is a pseudoscalar under orientation-reversing transformations and a scalar under the proper group used here.

For a future unit observer $u$, the electric and magnetic parts can be defined covariantly as
\begin{equation}
 E_u^\mu=F^{\mu\nu}u_\nu,
 \qquad B_u^\mu=(* F)^{\mu\nu}u_\nu.
\end{equation}
Both are orthogonal to $u$. In the rest frame of $u$, these are $(0,\mathbf E)$ and $(0,-\mathbf B)$ with our particular Hodge convention. One can instead put a minus sign in the definition of $B_u$ to make its spatial entries equal to the displayed $\mathbf B$; no physical assertion depends on this naming choice. The field itself, rather than either observer decomposition, is the Lorentz tensor.

In vacuum, $F\mapsto aF+b* F$ preserves the equations for constant $a,b\in K$. It preserves the stress tensor when $a^2+b^2=1$. This circle can be parameterized rationally over $K$ without defining trigonometric functions of unlimited angles:
\begin{equation}
 a=\frac{1-s^2}{1+s^2},\qquad b=\frac{2s}{1+s^2},\qquad s\in K,
\end{equation}
with the remaining point $(-1,0)$ added separately. The unit circle and its
angle group over the surreals are treated in \path{docs/surcomplex/trigonometry}. Infinitesimal duality angles can alternatively be treated by strongly summable power series.

\subsection{An exact null wave at any amplitude scale}\label{vtf:mink:sub:planewave}
Let $a\in K$, let $f\in C^\infty(\R,\R)$, and put $s=\tau-x^3$. Then
\begin{equation}\label{vtf:mink:eq:nullwave}
 \mathbf E=a f(s)(1,0,0),\qquad
 \mathbf B=a f(s)(0,1,0)
\end{equation}
is a source-free coherent Maxwell field. For example,
$\curl\mathbf B=(a f'(s),0,0)=\partial_\tau\mathbf E$, while
$\curl\mathbf E=(0,-a f'(s),0)=-\partial_\tau\mathbf B$; the divergences vanish. Both invariants in \cref{vtf:mink:eq:EMinvariants} are zero for every value of $a$, including infinitesimal and unlimited values.

The energy propagates in the $+x^3$ direction. With $k=(1,0,0,1)$, the stress tensor derived below is
\begin{equation}
 T^{\mu\nu}=a^2f(s)^2k^\mu k^\nu.
\end{equation}
A plane wave independent of $x^1,x^2$ generally has infinite total energy on a spatial slice even for a real amplitude. It is an exact local or per-unit-area example, not an example satisfying the compact-support energy hypotheses below.

For a finite wave covector $k=k_0+\kappa$, with real $k_0$ and infinitesimal components of $\kappa$, the surcomplex phase
\begin{equation}
 e^{i k\cdot x}:=e^{i k_0\cdot x}
                  \sum_{n\geq0}\frac{(i\kappa\cdot x)^n}{n!}
\end{equation}
belongs to the coherent complex coefficient algebra on real spacetime. With $k^2=0$ and transverse polarization, it gives the usual complex potential description of a vacuum wave. An unlimited covector requires a different function class; it is not covered by this definition.

\section{Stress--energy, positivity, and leading scales}\label{vtf:mink:sec:energy}

\subsection{The electromagnetic stress tensor}
Define
\begin{equation}\label{vtf:mink:eq:stress}
 T^{\mu\nu}
 =-F^{\mu\alpha}F^{\nu}{}_{\alpha}
   +\frac14\eta^{\mu\nu}F_{\alpha\beta}F^{\alpha\beta}.
\end{equation}
It is symmetric and traceless. Its Cartesian components are
\begin{align}
 T^{00}&=\mathcal E=\frac12(\mathbf E^2+\mathbf B^2),\label{vtf:mink:eq:energydensity}\\
 T^{0i}&=S_i=(\mathbf E\times\mathbf B)_i,\\
 T^{ij}&=\mathcal E\delta_{ij}-E_iE_j-B_iB_j.\label{vtf:mink:eq:stressspatial}
\end{align}
These follow by multiplying the finite matrices in \cref{vtf:mink:eq:Fmatrix}. In particular positivity of $T^{00}$ is positivity in the ordered coefficient field; no absolute value of an indefinite quadratic form has been substituted for energy.

\begin{proposition}[Local energy--momentum balance]\label{vtf:mink:prop:stressdiv}
If $F$ satisfies Maxwell's equations, then
\begin{equation}\label{vtf:mink:eq:stressdiv}
 \partial_\mu T^{\mu\nu}=-F^{\nu\lambda}J_\lambda.
\end{equation}
The time component is Poynting's identity,
\begin{equation}
 \partial_\tau\mathcal E+\Div\mathbf S=-\mathbf j\cdot\mathbf E.
\end{equation}
\end{proposition}
\begin{proof}
Differentiate \cref{vtf:mink:eq:stress}. The term differentiating the first factor of $F^{\mu\alpha}F^{\nu}{}_{\alpha}$ is $-J^\alpha F^{\nu}{}_{\alpha}$. Contracting the Bianchi identity with $F^{\mu\alpha}$ gives
\[
 F^{\mu\alpha}\partial_\mu F^{\nu}{}_{\alpha}
 =\frac14\partial^\nu(F_{\alpha\beta}F^{\alpha\beta}).
\]
This cancels the differentiated scalar term. The displayed balance law follows, and its $\nu=0$ component has the stated sign because $F^{0i}=-E_i$ and $J_i=-j_i$.
\end{proof}

\subsection{Dominant energy over a real-closed field}
The Euclidean identity $|\mathbf E\times\mathbf B|^2=\mathbf E^2\mathbf B^2-(\mathbf E\cdot\mathbf B)^2$ gives
\begin{equation}\label{vtf:mink:eq:DECidentity}
 \mathcal E^2-|\mathbf S|^2
 =\frac14(\mathbf E^2-\mathbf B^2)^2+(\mathbf E\cdot\mathbf B)^2
 =:\chi^2\geq0.
\end{equation}
All terms on the right are squares in $K$.

\begin{theorem}[Pointwise dominant energy]\label{vtf:mink:thm:DEC}
For any future unit timelike $u$, the vector
$j_u^\mu=T^{\mu\nu}u_\nu$ is future causal or zero. If $F\ne0$, then $T^{\mu\nu}u_\mu u_\nu>0$. These conclusions hold for arbitrary, not merely finite, components in $K$.
\end{theorem}
\begin{proof}
Use an algebraic rest boost for $u$. In that frame the current is $(\mathcal E,\mathbf S)$. Equation \cref{vtf:mink:eq:DECidentity} gives its causal character and the sum-of-squares expression gives its nonnegative time component. If $F\ne0$, at least one entry of $\mathbf E$ or $\mathbf B$ is nonzero, so $\mathcal E>0$. Transforming back preserves the future cone and the scalar energy measured by $u$.
\end{proof}

\subsection{The Rainich identity and the null case}
As an endomorphism, let $\mathcal T^\mu{}_{\nu}=T^{\mu\alpha}\eta_{\alpha\nu}$. Then
\begin{equation}\label{vtf:mink:eq:Rainich}
 \mathcal T^2=\chi^2 I,
 \qquad \chi^2=\frac{I_1^2+I_2^2}{16}.
\end{equation}
Here is a finite algebraic proof. Write $C=\mathcal E I_3-\mathbf E\mathbf E^{\mathsf T}-\mathbf B\mathbf B^{\mathsf T}$. The mixed tensor has blocks
\[
 \mathcal T=\begin{pmatrix}\mathcal E&-\mathbf S^{\mathsf T}\\
 \mathbf S&-C\end{pmatrix}.
\]
Orthogonality of $\mathbf S$ to both $\mathbf E$ and $\mathbf B$ gives $C\mathbf S=\mathcal E\mathbf S$. The identity
$C^2-\mathbf S\mathbf S^{\mathsf T}=(\mathcal E^2-|\mathbf S|^2)I_3$
follows by expanding the two rank-one products and using the cross-product identity. The block square is therefore \cref{vtf:mink:eq:Rainich}.

If $\chi>0$, the minimal polynomial divides $(X-\chi)(X+\chi)$, so $\mathcal T$ is diagonalizable over $K$. Tracelessness makes the multiplicities of $+\chi$ and $-\chi$ both two. This diagonalizability uses the special electromagnetic identity; it is not a general spectral theorem for Lorentz-self-adjoint tensors.

If $F\ne0$ and $\chi=0$, ordered-field positivity forces $\mathbf E^2=\mathbf B^2$ and $\mathbf E\cdot\mathbf B=0$. Thus $|\mathbf S|=\mathcal E>0$. Set $\mathbf n=\mathbf S/\mathcal E$ and $k=(1,\mathbf n)$. A spatial orthonormal basis adapted to $\mathbf E,\mathbf B,\mathbf n$ gives
\begin{equation}\label{vtf:mink:eq:nullstress}
 T^{\mu\nu}=\mathcal E k^\mu k^\nu,
 \qquad q_\eta(k)=0,
 \qquad \mathcal T^2=0,\quad \mathcal T\ne0.
\end{equation}
The stress is rank one, whereas its mixed endomorphism is nonzero nilpotent. Its scalar polynomial invariants can all vanish without its energy vanishing.

\subsection{Energy detects the leading electromagnetic scale}
For a two-form at an event, let $v_V(F)$ be the least valuation of its six independent entries. For coherent fields use the analogous least exponent of a nonzero component function.

\begin{theorem}[Exact doubling of the energy scale]\label{vtf:mink:thm:energyscale}
For every nonzero electromagnetic field at a $K$-valued event,
\begin{equation}\label{vtf:mink:eq:energyscale}
 v_V(T)=v(T^{00})=2v_V(F).
\end{equation}
For a nonzero coherent field on a real domain,
\begin{equation}
 \vH(T)=\vH(T^{00})=2\vH(F).
\end{equation}
These statements remain true for null fields with $I_1=I_2=0$.
\end{theorem}
\begin{proof}
Put $\alpha=v_V(F)$ and write each electric and magnetic entry as $t^\alpha$ times a finite Hahn number. At least one normalized entry has nonzero real residue. The coefficient at $2\alpha$ in $T^{00}$ is half the sum of the six real squared residues, and is strictly positive. Every component of $T$ is quadratic in the entries, so no component has valuation less than $2\alpha$. This proves the pointwise statement.

For coherent fields the leading entries are real functions. Their squared sum is a nonnegative smooth function and is not identically zero, because at least one leading entry is not identically zero. The same conclusion follows for global exponent order, despite the presence of zero divisors in the coefficient algebra.
\end{proof}
Together with \cref{vtf:mink:prop:frame}, this gives a bounded-frame-invariant scale detector. It avoids the false criterion that a small or zero scalar invariant must imply a small tensor.

\subsection{Integrated energy and conservation at every scale}
Assume each electric and magnetic coefficient on a spatial slice is compactly supported or Schwartz, and that the corresponding spacetime coefficients have the regularity and decay needed for integration by parts. Define
\begin{equation}
 \mathscr E(\tau)=\int_{\R^3}T^{00}(\tau,\mathbf x)\,\dd^3x
\end{equation}
coefficientwise. A finite convolution of such coefficient functions is integrable. If $\alpha$ is the leading exponent of the nonzero initial electromagnetic data, the coefficient at $2\alpha$ is
\begin{equation}
 \frac12\int_{\R^3}
 \bigl(|\mathbf E_\alpha|^2+|\mathbf B_\alpha|^2\bigr)\,\dd^3x>0.
\end{equation}
Consequently $v(\mathscr E)=2\alpha$. In vacuum, Poynting's identity and the absence of a boundary flux imply that $\mathscr E$ is constant, coefficient by coefficient. The leading field scale cannot disappear during this evolution. With sources, the corresponding work term is retained at every exponent.

This positivity argument is stronger than just reducing to the real shadow: an infinitesimal wave has zero real field and zero real energy shadow but still has a strictly positive exact Hahn energy at its first nonzero scale.

\section{Causal linear equations and retarded solution operators}\label{vtf:mink:sec:wave}

\subsection{The imported real theorem and the lifted data class}
The ordinary input is the smooth Cauchy theorem for a real normally hyperbolic operator on a globally hyperbolic spacetime: compatible smooth compactly supported initial data and source have a unique smooth solution, with finite propagation speed. We use its standard Minkowski specialization and the associated Green operators \cite{Bar,BGP}. This is an imported real theorem, not a consequence of real closedness.

Let $P_0$ be such an operator on a finite-rank real vector bundle, trivialized over the real domain under consideration. Basic examples are $\Box$ and $\Box+m_0^2$ with $m_0\in\R$. On the Cauchy slice $\Sigma=\{\tau=0\}$, prescribe coherent data
\begin{equation}
 f=\sum t^\gamma f_\gamma,\quad
 g=\sum t^\gamma g_\gamma,\quad
 h=\sum t^\gamma h_\gamma,
\end{equation}
where $h$ is the source. Suppose the union $S$ of their exponent supports is well ordered and each real coefficient belongs to the ordinary Cauchy theorem's data class. A common compact spacetime support is useful for uniform geometric statements, but is not needed merely to form the coefficientwise solution.

\begin{theorem}[Support-preserving Cauchy lifting]\label{vtf:mink:thm:Cauchy}
There is a unique coherent solution of
\begin{equation}\label{vtf:mink:eq:Cauchy}
 P_0u=h,\qquad u|_\Sigma=f,\qquad
 \nabla_nu|_\Sigma=g.
\end{equation}
Its exponent support is contained in $S$. Each real coefficient satisfies the ordinary finite-propagation estimate. If all coefficient data and sources have support in a common closed set $C$ for which the classical causal support theorem applies, then every coefficient solution is supported in $J(C)$; when $J(C)$ is closed, the closure of their union is also contained there.
\end{theorem}
\begin{proof}
For each $\gamma\in S$, solve the real Cauchy problem with data $(f_\gamma,g_\gamma,h_\gamma)$, setting missing coefficients to zero. Call the solution $u_\gamma$. The series $u=\sum_{\gamma\in S}t^\gamma u_\gamma$ is coherent because its support is a subset of $S$. The equations follow by coefficient extraction. If two coherent solutions have the same data, every coefficient of their difference satisfies a real homogeneous zero-data Cauchy problem and is zero. The real finite-propagation theorem applies at each coefficient, proving the support assertions.
\end{proof}

For a nonzero solution of a homogeneous equation, the least exponent of its initial data equals the least exponent of the global solution, by real uniqueness. On an individual time slice, a field component may vanish while its normal derivative does not, so the invariant data are the pair of Cauchy traces, not a single chosen component.

\begin{remark}[What kind of well-posedness has been proved?]
The theorem gives existence, uniqueness, support preservation, and real continuous dependence at each coefficient. It does not assert a uniform Sobolev bound across all exponents, convergence after setting $t$ equal to a positive real number, or well-posedness in the fine topology of $\No$. Such assertions need additional norms, growth conditions, and estimates.
\end{remark}

\subsection{The retarded massless kernel}
For the sign convention $\Box=\partial_\tau^2-\Delta$ in three spatial dimensions, the real retarded inverse is
\begin{equation}\label{vtf:mink:eq:retarded}
 (G_{0,\mathrm{ret}}h)(\tau,\mathbf x)
 =\frac1{4\pi}\int_{\R^3}
 \frac{h(\tau-|\mathbf x-\mathbf y|,\mathbf y)}{|\mathbf x-\mathbf y|}
 \,\dd^3y,
\end{equation}
for smooth compactly supported sources. Its distribution kernel is
$\delta(\tau-r)/(4\pi r)$ with future support, interpreted distributionally at the cone vertex. Applying \cref{vtf:mink:eq:retarded} to every coefficient gives the retarded Hahn inverse with unchanged exponent support. The sign is opposite to conventions that define the wave operator as $\Delta-\partial_\tau^2$; compare the radiation notes \cite{TongRad}.

For a conserved current with admissible retarded support, $A^\nu=G_{0,\mathrm{ret}}J^\nu$ obeys the Lorenz condition. Indeed, real differentiation commutes with this translation-invariant Green operator on its support class, so
$\partial_\nu A^\nu=G_{0,\mathrm{ret}}(\partial_\nu J^\nu)=0$.
Then $F=\dd A$ solves the sourced Maxwell equations. This does not prescribe arbitrary independent initial radiation; a homogeneous solution must be added when such data are intended.

\subsection{Positive-valuation perturbations of a hyperbolic operator}
Consider
\begin{equation}\label{vtf:mink:eq:perturbedP}
 P=P_0+Q,\qquad
 Q=\sum_{\sigma\in S_Q}t^\sigma Q_\sigma,
 \qquad S_Q\subset\Gamma_{>0}\text{ well ordered},
\end{equation}
where every $Q_\sigma$ is a real smooth local differential operator of a fixed finite order. Work in a retarded class in which $G_{0,\mathrm{ret}}$ is defined for the intermediate real sources. One convenient class consists of fields supported in $J^+(C)$ for a common compact set $C$: it is past compact in Minkowski space, the standard extension of the real retarded Green operator applies, and local differential operators preserve that support. Alternatively, use a fixed real time slab and a zero-initial-data solution operator throughout. The retarded extension to sources supported in $J^+(C)$ can be constructed locally: for a compact target region $D$, the intersection $J^-(D)\cap J^+(C)$ lies in a compact causal diamond. Cut off the source outside a neighborhood of that diamond and apply the compact-source Green operator. Finite propagation makes the result on $D$ independent of the cutoff, and uniqueness glues these local results.

\begin{theorem}[Exact causal inverse for a positive-order perturbation]\label{vtf:mink:thm:retardedperturb}
For coherent sources in the specified retarded class, the operator
\begin{equation}\label{vtf:mink:eq:NeumannGreen}
 G_{P,\mathrm{ret}}
 =\sum_{n\geq0}(-G_{0,\mathrm{ret}}Q)^nG_{0,\mathrm{ret}}
\end{equation}
is well defined by strong summation and gives the unique retarded coherent solution of $Pu=h$. If $h$ has exponent support $S$, its solution has exponent support contained in
\begin{equation}
 S+\langle S_Q\rangle.
\end{equation}
Every coefficient remains supported in the ordinary causal future prescribed by the real background operator.
\end{theorem}
\begin{proof}
In a term of length $n$, choose $n$ exponents $\sigma_1,\ldots,\sigma_n$ from $S_Q$. A contribution to
exponent $\gamma$ has $\gamma=s+\sigma_1+\cdots+\sigma_n$ with $s\in S$. The positive-support lemma makes the union of all such supports well ordered and leaves only finitely many contributing words, over all lengths, at any fixed $\gamma$. Each word is a finite composition of real differential and Green operators. Thus every coefficient is a finite sum of smooth functions in the admitted support class.

The series solves $(I+G_{0,\mathrm{ret}}Q)u=G_{0,\mathrm{ret}}h$ by formal telescoping. Applying $P_0$ gives $(P_0+Q)u=h$. For uniqueness, a nonzero retarded solution of the homogeneous equation would have a least exponent $\alpha$. At that exponent $Q$ contributes nothing from the difference, because all its exponents are positive. The leading coefficient would solve $P_0u_\alpha=0$ with zero retarded data, contradicting the real uniqueness theorem. Locality of $Q_\sigma$ and causality of $G_{0,\mathrm{ret}}$ prove the support claim for each word.
\end{proof}

For example a finite mass-squared parameter $m^2=m_0^2+\mu$, with $\mu$ of positive order, is handled by $Q=\mu$. A positive-order perturbation of the principal coefficients is also admissible as a \emph{formal coefficient problem}. The theorem's causal support is relative to the real body $P_0$. It is not a statement that the exact characteristic cone of every $K$-valued principal symbol coincides with the real Minkowski cone, or that the formal inverse converges for small positive real substitutions of the scale.

\subsection{Tensor wave systems}
A real normally hyperbolic operator may act on any fixed finite tensor representation. The proof of \cref{vtf:mink:thm:Cauchy} is unchanged because there are only finitely many components. Additional algebraic or differential constraints must be checked to propagate. For Maxwell theory this was done explicitly; for a proposed higher-rank gauge field it is a separate obligation. Scalar extension does not by itself remove gauge degeneracy or turn a constrained system into a normally hyperbolic one.

\section{Actions, scalar fields, fluids, and conserved charges}\label{vtf:mink:sec:actions}

\subsection{A variational calculus with real test variations}
Let $\psi^a$ be finitely many coherent fields, and let
$L(x,\psi,\partial\psi)$ be a polynomial in their entries and first derivatives with coherent coefficients. On a compact real spacetime region $D$, define
\begin{equation}
 S_D[\psi]=\int_D L(x,\psi,\partial\psi)\,\dd^4x
\end{equation}
coefficientwise. The first variation is the coefficient of $s$ in $S_D[\psi+s\zeta]$ for an ordinary real formal parameter $s$. For compactly supported variations in the interior, coefficientwise integration by parts gives
\begin{equation}\label{vtf:mink:eq:EL}
 \delta S_D[\psi;\zeta]
 =\int_D\sum_a
 \left(\frac{\partial L}{\partial\psi^a}
 -\partial_\mu\frac{\partial L}{\partial(\partial_\mu\psi^a)}\right)
 \zeta^a\,\dd^4x.
\end{equation}
It is enough to test all ordinary real smooth compactly supported $\zeta$: extracting any exponent and applying the real fundamental lemma of the calculus of variations proves the Euler--Lagrange equations at that exponent. Conversely, those equations make \cref{vtf:mink:eq:EL} zero for every coherent test variation for which the products and integral are defined.

This is a theory of stationarity, with the fundamental-lemma argument of
\cref{vtf:sec:variation}. It does not infer the existence of action minimizers from the order completeness of $K$; that completeness is not available. It also does not require a measure on all of $K^4$.

\subsection{Scalar fields and their stress tensor}
For a real $K$-valued scalar $\phi$, let
\begin{equation}
 L_0=\frac12\partial_\mu\phi\,\partial^\mu\phi-V(\phi),
 \qquad L=L_0+J\phi,
\end{equation}
where $V$ is a polynomial with constant coefficients in $K$. With the sign convention of \cref{vtf:rem:potentialsign} for $V$, the
equation of motion is
\begin{equation}\label{vtf:mink:eq:scalarfield}
 \Box\phi+V'(\phi)=J.
\end{equation}
The symmetric scalar stress tensor is
\begin{equation}\label{vtf:mink:eq:scalarstress}
 T^{\mu\nu}_\phi
 =\partial^\mu\phi\,\partial^\nu\phi-\eta^{\mu\nu}L_0.
\end{equation}
Differentiating and cancelling the mixed-derivative terms gives
\begin{equation}
 \partial_\mu T^{\mu\nu}_\phi
 =(\Box\phi+V'(\phi))\partial^\nu\phi
 =J\partial^\nu\phi.
\end{equation}
For $V(\phi)=\tfrac12m^2\phi^2+\tfrac\lambda4\phi^4$, with $m^2,\lambda\geq0$, the rest-frame energy density
\begin{equation}
 T^{00}_\phi=\frac12(\partial_\tau\phi)^2
             +\frac12|\grad\phi|^2+V(\phi)
\end{equation}
is nonnegative in $K$. Negative potentials or other matter Lagrangians require their own energy analysis; positivity is not automatic for every tensor field.

For electromagnetism the Lagrangian
\begin{equation}
 L_{\mathrm{EM}}=-\frac14F_{\mu\nu}F^{\mu\nu}-J^\mu A_\mu
\end{equation}
produces $\partial_\mu F^{\mu\nu}=J^\nu$ with the conventions already fixed. Its symmetric gauge-invariant field stress is \cref{vtf:mink:eq:stress}. The sign of the source term differs from the scalar source convention because the fields and their defining kinetic terms differ; each equation follows from its displayed Lagrangian.

\subsection{Translations, Lorentz transformations, and charges}
For a translation-invariant Lagrangian the canonical current is
\begin{equation}
 \Theta^\mu{}_{\nu}
 =\sum_a\frac{\partial L}{\partial(\partial_\mu\psi^a)}
          \partial_\nu\psi^a-\delta^\mu{}_{\nu}L.
\end{equation}
A direct application of the Euler--Lagrange equations proves its conservation when there is no explicit coordinate dependence. For a symmetric conserved stress tensor,
\begin{equation}\label{vtf:mink:eq:angularcurrent}
 \mathcal M^{\lambda\mu\nu}
 =x^\mu T^{\lambda\nu}-x^\nu T^{\lambda\mu},
 \qquad \partial_\lambda \mathcal M^{\lambda\mu\nu}=0.
\end{equation}
Indeed the differentiated coordinates produce $T^{\mu\nu}-T^{\nu\mu}=0$, and the remaining terms contain the stress divergence.

On a real Cauchy slice, with coefficientwise integrability and vanishing boundary flux, the quantities
\begin{equation}
 P^\nu=\int T^{0\nu}\,\dd^3x,
 \qquad J^{\mu\nu}=\int \mathcal M^{0\mu\nu}\,\dd^3x
\end{equation}
are conserved Hahn-valued charges. More invariantly, one integrates $T^{\mu\nu}\dd\Sigma_\mu$ over a real spacelike slice. Stokes' theorem gives independence of the slice under the stated boundary conditions. Ordinary real Lorentz covariance follows coefficientwise. Integration over a genuinely surreal hypersurface would be an additional construction; the real-chain definition does not silently supply it.

For finite fields satisfying polynomial differential equations with finite constant coefficients, \cref{vtf:mink:eq:stdifferential} gives the corresponding real equations for their shadows. With appropriate integrability the same is true of charges. Thus shadow reduction is compatible with the local dynamics in this regime, although it discards positive-order energy and momentum.

\subsection{A perfect-fluid tensor}
Let $u$ be a future unit coherent vector field and let $\rho,p$ be coherent scalar fields. Algebraically,
\begin{equation}\label{vtf:mink:eq:fluid}
 T^{\mu\nu}=(\rho+p)u^\mu u^\nu-p\eta^{\mu\nu}
\end{equation}
is a well-defined perfect-fluid stress. Put $\mathrm D_u=u^\mu\partial_\mu$, $\vartheta=\partial_\mu u^\mu$, and
$a^\nu=\mathrm D_uu^\nu$. Contracting $\partial_\mu T^{\mu\nu}=0$ first with $u_\nu$ and then with the spatial projector gives
\begin{equation}
 \mathrm D_u\rho+(\rho+p)\vartheta=0,
 \qquad (\rho+p)a^\nu=h^{\nu\mu}\partial_\mu p,
 \quad h^{\nu\mu}=\eta^{\nu\mu}-u^\nu u^\mu.
\end{equation}
In an instantaneous rest frame the second equation has the usual negative spatial pressure gradient. The identity $u\cdot a=0$ follows by differentiating $u\cdot u=1$.

If $\rho\geq|p|$, the tensor satisfies the dominant energy condition. In its rest frame the mixed tensor is $\diag(\rho,-p,-p,-p)$, so the energy current of a future causal vector $(v^0,\mathbf v)$ is $(\rho v^0,-p\mathbf v)$, whose squared length is at least $\rho^2((v^0)^2-|\mathbf v|^2)\geq0$ and whose time component is nonnegative. This is an ordered-field proof. An equation of state, regularity, shock theory, and hyperbolic well-posedness are separate dynamical inputs, not consequences of \cref{vtf:mink:eq:fluid}.

\section{A nonlinear wave theorem by transfinite coefficient recursion}\label{vtf:mink:sec:nonlinear}

\subsection{Fixing the real background}
Let $I\subset\R$ be an open time interval containing zero and let $U=I\times\R^3$. Fix $m_0^2,\lambda\in\R$ and a smooth real solution $\phi_0$ on $U$ of
\begin{equation}
 (\Box+m_0^2)\phi_0+\lambda\phi_0^3=f_0.
\end{equation}
No claim is made that every real datum has such a global background. Choose positive-order coherent perturbations of the source and initial data. Assume their nonzero exponents have well-ordered union $S_r\subset\Gamma_{>0}$ and that the real coefficient data are smooth, with compact support or other hypotheses sufficient for the real linear Cauchy problems used below. Let $\mathfrak M=\langle S_r\rangle$.

The linearized real operator
\begin{equation}\label{vtf:mink:eq:linearized}
 L_0=\Box+m_0^2+3\lambda\phi_0^2
\end{equation}
is normally hyperbolic. It therefore has a smooth real linear Cauchy theory on the fixed background slab. When a recursive source is not compact in time, the needed smooth Cauchy solution is obtained by causal localization on compact diamonds and uniqueness; compact initial supports and finite propagation are sufficient for the recursive products, and an arbitrary smooth source may likewise be localized. Smooth coefficients can grow toward the ends of $I$; the construction asserts smoothness on $U$, not a uniform bound at an endpoint outside $I$.

\begin{theorem}[Positive-order cubic-wave deformation]\label{vtf:mink:thm:cubic}
There is a unique coherent solution $\phi=\phi_0+H$, with $H$ of strictly positive exponent support, having the prescribed perturbed initial data and satisfying
\begin{equation}\label{vtf:mink:eq:cubic}
 (\Box+m_0^2)\phi+\lambda\phi^3=f_0+r.
\end{equation}
The support of $H$ is contained in $\mathfrak M\setminus\{0\}$. Each coefficient is obtained by solving one real linear Cauchy problem for $L_0$ with a source formed from finitely many lower-exponent coefficients.
\end{theorem}
\begin{proof}
Subtract the background equation and expand the cubic:
\begin{equation}\label{vtf:mink:eq:cubicH}
 L_0H=r-\lambda(3\phi_0H^2+H^3).
\end{equation}
The positive-support lemma makes $\mathfrak M$ well ordered and gives only finitely many decompositions into words from $S_r$ at a fixed exponent. Suppose the coefficients $H_\beta$ have been constructed at all positive $\beta<\gamma$ in $\mathfrak M$. The equation for the coefficient at $\gamma$ is
\begin{equation}\label{vtf:mink:eq:cubicrecursion}
 \begin{split}
 L_0H_\gamma={}&r_\gamma
 -3\lambda\phi_0\sum_{\substack{\beta+\zeta=\gamma\\\beta,\zeta>0}}
 H_\beta H_\zeta\\
 &-\lambda\sum_{\substack{\beta+\zeta+\kappa=\gamma\\\beta,\zeta,\kappa>0}}
 H_\beta H_\zeta H_\kappa.
 \end{split}
\end{equation}
Every exponent on the right is smaller than $\gamma$, and each sum is finite. The prescribed initial coefficients at $\gamma$ and the real Cauchy theorem determine $H_\gamma$ uniquely. Transfinite recursion along $\mathfrak M\setminus\{0\}$ constructs the family; at limit positions there is no additional topological limit, only the previously constructed lower coefficients.

The family has well-ordered support by construction, and coefficientwise substitution into \cref{vtf:mink:eq:cubicH} verifies the equation. To prove uniqueness among all coherent positive-order deformations, not just those preassigned to $\mathfrak M$, let a nonzero difference have least exponent $\gamma$. All terms in the nonlinear difference have that difference multiplied by at least one positive-order perturbation, so they cannot contribute at $\gamma$. The coefficient difference solves a homogeneous zero-data Cauchy problem for $L_0$ and must be zero, a contradiction.
\end{proof}

The result is an exact formal existence and uniqueness theorem over a Hahn coefficient algebra. It does not prove convergence of the deformation at a small real parameter, nor prevent a different real solution from developing a singularity outside the background domain. It is stronger than writing an unevaluated perturbation ansatz: every coefficient is determined, all convolution sums are certified finite, and the support class is closed.

The theorem is the Lorentzian instance of the recursion of
\cref{vtf:thm:lifting}, with a real Cauchy problem in place of a bijective real
linearization.

\subsection{A worked spatially homogeneous example}
On a real time interval, consider
\begin{equation}
 y''+y^3=0,\qquad y(0)=t,\quad y'(0)=0.
\end{equation}
This is the spatially homogeneous, massless, unit-coupling special case. It is used as an ODE example rather than as compactly supported spatial data. The recursion yields
\begin{equation}\label{vtf:mink:eq:cubicoscillator}
 y(\tau)=t-\frac{\tau^2}{2}t^3
 +\frac{\tau^4}{8}t^5
 -\frac{3\tau^6}{80}t^7
 +\frac{7\tau^8}{640}t^9+\cdots.
\end{equation}
For example $y_1=1$, $y_3''=-1$, and $y_5''=-3y_3=3\tau^2/2$, with zero higher-order initial data. At exponent seven the source is
$-(3y_5+3y_3^2)=-9\tau^4/8$, giving $y_7=-3\tau^6/80$. At exponent nine it is
$-(3y_7+6y_3y_5+y_3^3)=49\tau^6/80$, giving $y_9=7\tau^8/640$.

The exact conserved energy is
\begin{equation}
 \frac12(y')^2+\frac14y^4=\frac14t^4.
\end{equation}
Its derivative is $y'(y''+y^3)$, so it vanishes coefficientwise and the initial value determines it. The accompanying symbolic program verifies the displayed truncation through exponent nine; the theorem, not the finite test, establishes the all-order solution.

\subsection{Extension to other systems}
The proof uses three ingredients: a specified real background, a well-posed real linearized Cauchy problem on a common domain, and nonlinear terms that raise positive exponent order. Thus the same architecture applies to finite polynomial systems when those ingredients are separately verified. A nonpolynomial nonlinearity can be treated by a support-admissible formal Taylor expansion around the body, but its domain and coefficient regularity must be stated.

For nonabelian gauge fields, the finite formulas $F_A=\dd A+A\wedge A$ and $\dd_AF_A=0$ extend over coherent matrix coefficients, the latter by associativity and the graded Leibniz rule. A full evolution theorem additionally requires a gauge choice, constraint propagation, and a real linearized hyperbolic theory. Those extra steps are not replaced by the scalar recursion above.

\section{Worldlines, local scales, and derivative choices}\label{vtf:mink:sec:worldlines}

\subsection{Parameterized surreal-valued trajectories}
A coherent trajectory on an ordinary real interval is $X\in\A_\Gamma(I)^4$. Its velocity and acceleration are defined coefficientwise. If its parameter $s$ satisfies $q_\eta(\dot X)=1$ and $\dot X^0>0$, it is proper-time normalized within the coefficient theory. Differentiating gives
\begin{equation}
 \eta(\dot X,\ddot X)=0.
\end{equation}
Such a trajectory is a real-parameter family of $K$-valued events, not an ordinary path in $\R^4$ unless all coefficients outside exponent zero vanish. It also need not be continuous for the fine topology on the surreal target.

A proper-time lapse can be constructed locally under a nondegenerate leading-term hypothesis. Suppose
\begin{equation}
 q_\eta(\dot X)=t^{2\alpha}\bigl(r_0+R\bigr),
 \qquad r_0\in C^\infty(I,\R_{>0}),\quad R\in\A_{>0}(I).
\end{equation}
Then
\begin{equation}
 \sqrt{q_\eta(\dot X)}
 =t^\alpha\sqrt{r_0}
 \sum_{n\geq0}\binom{1/2}{n}(R/r_0)^n
\end{equation}
is a coherent positive lapse. If the leading real coefficient vanishes at some parameter value, a single such expansion need not extend through it. One must reanalyze the scale or restrict the domain. The existence of pointwise square roots in $K$ alone does not give one coherent smooth square-root field across every zero.

\subsection{The Lorentz force and admissible composition}
For nonzero constant mass $m\in K$, charge $e\in K$, and proper velocity $U$, the Lorentz force equation is
\begin{equation}\label{vtf:mink:eq:Lorentzforce}
 m\dot U^\mu=eF^{\mu\nu}(X)U_\nu,
 \qquad \dot X=U.
\end{equation}
The expression $F(X)$ must be admitted by the chosen function class: for finite trajectories with real body in the domain one uses the halo construction; for polynomial fields one may evaluate at arbitrary $K$-valued trajectories. Antisymmetry gives
\begin{equation}
 \frac{\dd}{\dd s}q_\eta(U)=\frac{2e}{m}U_\mu F^{\mu\nu}(X)U_\nu=0.
\end{equation}
Thus any solution with unit initial velocity remains unit. Existence of solutions, especially for unlimited $eF/m$, is not implied by this conserved algebraic identity.

For constant coefficients of finite body plus positive-order perturbation, the corresponding linear system has a coefficientwise perturbative fundamental matrix built around its ordinary real fundamental matrix. Conversely, an unlimited oscillatory generator can fall outside the coherent class. Unlimited coefficients are not universally forbidden: a nilpotent constant matrix has a finite exponential polynomial, so its evolution can remain coherent even at unlimited scale. It is the support of the actual operation, not a blanket finiteness slogan, that decides admissibility.

\subsection{Inner coordinates and scale-dependent differential operators}\label{vtf:mink:subsec:inner}
Fix a center $a\in K^4$ and a scale $t^b>0$. In an inner chart put
\begin{equation}
 x=a+t^b X.
\end{equation}
For an admissible scalar expression, the chain rule yields
\begin{equation}
 \partial_{x^\mu}=t^{-b}\partial_{X^\mu},
 \qquad \Box_x=t^{-2b}\Box_X.
\end{equation}
The metric written in the $X$ coordinates is $t^{2b}\eta$, and the volume form is scaled by $t^{4b}$. The inverse metric supplies the inverse factor in $\Box$. Thus coordinate derivative scales cannot be inferred just from the amplitude of the field.

This inner chart is an algebraic or halo parameterization of a region of $K^4$. When $b\ne0$, it is not an ordinary diffeomorphism between two open subsets of real spacetime. Inner and outer coherent coefficient expansions require an explicit overlap and matching rule. Such scale bookkeeping is useful for boundary layers, concentrated sources, and asymptotic regimes without declaring a divergent quantity regular.

\section{Covariance, singular sources and interpretation}\label{vtf:mink:sec:limits}

\subsection{An infinite boost leaves the coherent class}\label{vtf:mink:subsec:covariance}

\Cref{vtf:mink:thm:highfrequency} explains a genuine covariance boundary. A real null plane wave transformed by a suitable infinite boost acquires an unlimited frequency in an ordinary real coordinate direction. Its would-be transformed phase would satisfy an equation of the form \cref{vtf:mink:eq:highfreq}; it therefore cannot remain a nonzero element of the same coherent function algebra. Full $K$-Lorentz covariance and this particular globally coherent smooth coefficient class cannot both be asserted without enlarging the class or changing the base.

\subsection{Singular sources and punctured domains}
The electrostatic Coulomb field
\begin{equation}
 \mathbf E(\mathbf x)=\frac{e}{4\pi}\frac{\mathbf x}{r^3},
 \qquad r=\sqrt{(x^1)^2+(x^2)^2+(x^3)^2},
 \qquad r\ne0,
\end{equation}
has an algebraic interpretation at every nonzero $K$-valued spatial point. At an infinitesimal radius its components may be unlimited. The denominator is still zero at the origin, so scalar extension has not defined a field there. A distributional source at the origin belongs to the coefficientwise linear distribution theory, but products needed for nonlinear self-energy require additional regularization and cannot be inferred from that linear statement.

This is the flat-space analogue of the repository's broader warning that scalar extension can preserve a singular expression's exact asymptotic content without regularizing its domain \cite{RepoPhysics}. No claim about black-hole interiors or quantum gravity follows from the Minkowski construction.

\subsection{What standard part does and does not preserve}
For coherent fields with nonnegative support, standard part preserves polynomial differential expressions. It preserves rational expressions only when the denominators have nonvanishing real body on the chosen domain, so that their inverses have coherent finite expansions. For example the finite nonzero scalar $t$ has an unlimited inverse, and reduction cannot turn the identity $\eps t^{-1}=1$ into multiplication of two real shadows.

Even when every displayed field and tensor is finite, shadow reduction need not preserve a strict causal classification: the near-null example of \cref{vtf:mink:sec:bounded} proves otherwise. It also fails to commute with unbounded frame changes, as \cref{vtf:mink:ex:boostshadow} shows. Accordingly, a real readout map defined by standard part has a natural bounded-frame covariance domain, not a universal full-Lorentz equivariance property.

\subsection{Mathematics versus a physical interpretation}
The construction provides exact relativistic tensor identities, coherent multiscale fields, and controlled formal evolution. It does not by itself select the physically realized exponent group, attach operational units to an arbitrary surreal scale, prescribe detector response, or produce experimentally different predictions. A physical model would need those additional choices together with a correspondence to measured real quantities.

In particular, an unlimited field value is an exact mathematical value in an ordered extension, not a proof that an ordinary apparatus can measure an infinite real energy. An infinitesimal correction has no nonzero standard-part readout without a specified amplification or scaling procedure. The infinite-boost example is an algebraic illustration of amplification; it is not a claim that an observer of unlimited Lorentz factor can be physically prepared.

These interpretive limits are source 02's own. Together with the firewall of
\cref{vtf:subsec:merge} they govern the whole of Part~III.

\part{Synthesis}\label{vtf:part:synthesis}

\section{What the theory establishes, and what it leaves as choices}\label{vtf:sec:conclusion}

The resulting theory has three genuine strengths. It accommodates every
finite tensor type in every finite dimension with exact infinite and
infinitesimal scalar values. It provides a coordinate-compatible calculus
with certified infinite operations, including actual intrinsic derivatives
on a large class of surreal-coordinate domains. And it converts specified
real linear solvability into exact nonlinear multiscale solutions without
requiring ordinary convergence of perturbation series.

Its boundaries are equally concrete. A bare map into $\No^n$ is not a
smooth field. Pointwise nondegeneracy is not a certificate that an inverse
belongs to a particular function algebra. The intrinsic topology of a fixed
workspace is not the topology inherited from the full proper class.
Coefficientwise integration is not an unrestricted measure theory on all
surreal subsets. Standard part preserves regular geometry but can destroy
rank in a multiscale metric. A field derivation on scalars is not a spatial
partial derivative. None of these distinctions disappears when $n$ is
specialized to three or four; Parts~II and~III show them at work there.

Source 03's three-dimensional development has a definite core. Vectors and
finite tensors over surreal scalars obey the expected algebraic identities;
compact-chain integration obeys exact Stokes identities; cohomology keeps the
ordinary topological degrees of freedom while periods become surreal; Hodge
decomposition survives both scalar extension and positive-order metric
perturbation; and invertible real operators and quadratic nonlinearities admit
exact support-controlled lifts. Its boundary is equally definite: a
common-support field is more restrictive than a collection of pointwise surreal
values, pointwise invertibility does not give a field inverse, singular scale
operators can force solutions out of the calculus, and standard part can lose
rank, flux, energy and curvature and can fail to commute with raising indices for
a degenerate leading metric.

Source 02's Lorentzian package is a compatible collection of algebraic, coherent
and halo theories with explicit overlaps, not a single unrestricted analytic
theory on all of $\No^4$. The finite algebra has full $K$-Lorentz covariance; the
coherent calculus has real-spacetime and bounded-halo covariance; Maxwell theory
has exact gauge symmetry, local conservation, dominant energy, null radiation and
a leading-scale energy detector; and real hyperbolic solution operators lift
without enlarging support, while positive-order perturbations enlarge it only by
explicitly generated monoids. The infinite-frequency theorem and the
infinite-time example show that omitting these boundaries would produce false
results.

\subsection{Open directions and their status}\label{vtf:subsec:open}

Each source ends with proposed directions. They are merged below; where another
source answers part of a direction, the answered part and its source are
recorded and the rest stays open. None of these is asserted to be unsolved in
every existing formalism.
\begin{enumerate}
\item \emph{Larger intrinsic function classes} admitting coordinate gluing and
integration (source 01), in particular an amplitude algebra enlarged by
explicitly controlled oscillatory phases and boundary-layer charts with precise
transition rules (sources 02 and 03), so that a larger class of infinite boosts
can act without violating \cref{vtf:mink:thm:highfrequency}. Open.
\item \emph{Hodge and elliptic results for genuinely surreal leading metrics}
(source 01), and a Hodge theory for anisotropic metrics whose real reduction is
singular, tracking scale-dependent cohomological and coercivity data (source 03).
\emph{Partly answered:} for a regular Riemannian Hahn metric, a nondegenerate
positive-definite real leading part plus a positive-support perturbation, source
03 proves the orthogonal decomposition, harmonic representatives and Betti
dimension (\cref{vtf:3d:thm:regularhodge}). \emph{Open:} regular metrics of
indefinite signature, and every metric whose standard part is degenerate, such as
$\diag(1,t^2,t^4)$.
\item \emph{Scale charts around rank-changing real limits} (source 01). Open;
\cref{vtf:subsec:warped} and the softened monopole of \cref{vtf:3d:sec:physics}
are worked examples, not a classification.
\item \emph{Support certificates for specific gauge-fixed nonlinear systems}
(source 01), and a variable Lorentzian metric with nondegenerate real body
(source 02). \emph{Partly answered:} the inverse, Levi--Civita connection,
curvature, Bianchi identities, first-order response and standard-part reduction
of a regular metric of any signature are proved in
\cref{vtf:sec:metric,vtf:sec:curvature,vtf:sec:reduction} (source 01); specific
nonlinear instances are proved for stationary incompressible flow on the
three-torus (\cref{vtf:3d:cor:ns}, source 03) and for a cubic wave equation
around a real background (\cref{vtf:mink:thm:cubic}, source 02). \emph{Open:}
gauge fixing, causal propagation and the nonlinear gravitational equations, and
any global Einstein, Yang--Mills or Navier--Stokes evolution statement.
\item \emph{A locally coherent sheaf theory with precise global support criteria
for integration and evolution} (source 02). \emph{Partly answered:} the sheaf
$\A^{\mathrm{loc}}_\Gamma$, its local calculus, and finite localization over
compact sets and compact chains for integration (\cref{vtf:sec:sheaf},
source 01). \emph{Open:} global support criteria for evolution on noncompact
domains.
\item \emph{Computation and formalization}: a computational treatment of support
certificates that reliably detects which truncations are finite (source 03), and
formalization of the prolongation and unit-admissibility interfaces (source 01),
following the layer orders of \cref{vtf:sec:implementation}. Open.
\end{enumerate}
The constructions already proved give a working foundation for these projects,
while retaining a sharp distinction between exact scale representation and
analytic or physical regularization.

\appendix

\section{Notation and hypothesis ledger}\label{vtf:app:ledger}

\begin{longtable}{P{0.24\textwidth}P{0.68\textwidth}}
\toprule
Symbol or phrase & Meaning and scope\\
\midrule\endhead
$\No$ & The proper class of Conway surreal numbers; not a set-sized scalar type without a universe qualification.\\
$\Gamma$ & A nonzero set-sized divisible ordered additive subgroup of $\No$.\\
$t^\gamma$ & Conway monomial $\omega^{-\gamma}$, not an implicit use of surreal exponentiation.\\
$K$ & The real closed Hahn workspace $\R((t^\Gamma))$, with its intrinsic topology when derivatives are taken.\\
$\A_\Gamma(U)$ & Smooth real coefficient functions on one common real domain, with one well-ordered outer Hahn support.\\
$\A^{\mathrm{loc}}_\Gamma$ & The sheaf of local such presentations; globally uniform support is not assumed on noncompact bases.\\
$v$, $\vH$ & Minimum exponent of a scalar; minimum exponent of a nonzero coefficient function in a local field presentation.\\
$\halo U$ & $\st^{-1}(U)\subset K^n$; a nearstandard open set, not all of $K^n$.\\
$D_i$ & Coefficientwise spatial differentiation, fixing every constant in $K$.\\
$\partial_{y^i}$ & Intrinsic ordered-field derivative of a realized function; agrees with $\ev D_i$ only where the bridge applies.\\
$\BM$ & The optional Berarducci--Mantova scalar derivation, not a spatial derivative.\\
Admissible metric & Symmetric metric whose inverse lies in the chosen coefficient algebra; its density is also required when used.\\
Regular metric & $g=g_0+h$ with real nondegenerate $g_0$ and positive-support $h$.\\
$\Delta_g$, $\Delta_H$ & Divergence--gradient scalar Laplacian; $d\delta+\delta d$. On functions $\Delta_H=-\Delta_g$.\\
$\int^H$, $\int^{\mathrm{alg}}$ & Coefficientwise integration over real chains; polynomial moment integration on affine $K$-simplices.\\
Strong sum & Well-ordered support union and finitely many nonzero contributions at every exponent; no ordinary convergence assertion.\\

$\eps$ & A chosen positive infinitesimal $t^\eta$; source 02's $\eps$ is $t$ here.\\
$t$, $\tau$ & $t=t^1=\omega^{-1}$ when $1\in\Gamma$ (Parts II and III); $\tau$ is ordinary time.\\
$\partial_i$, $\partial_\mu$ & In Parts II and III, the coefficientwise derivatives $D_i$.\\
$\mathscr D_\chi$, $\mathcal D_i$ & Character scale derivation and scale connection of \cref{vtf:sec:BM}; not spatial derivatives.\\
$v_V$ & Least component valuation of a finite tensor in a stated frame (source 02's $\nu$).\\
Coherent field & An element of $\A_\Gamma(U)$: one well-ordered support on all of $U$.\\
Regular Riemannian metric & Regular metric with $g_0$ positive definite (source 03's regular Hahn metric).\\
$\epsilon_{ijk}$, $\epsilon_{\mu\nu\rho\sigma}$ & Alternating symbols.\\
$\eta_{\mu\nu}$, $q_\eta$ & Minkowski metric $\diag(1,-1,-1,-1)$ and its quadratic form (Part III only).\\
$\mathcal G$, $\mathcal G_{\OO_K}$, $\mathcal G_1$ & $\SO^+(1,3;K)$, its entrywise finite subgroup, and the kernel of reduction.\\
$\Box$ & $\partial_\tau^2-\Delta$.\\
$S_Q$, $S_r$ & Exponent sets of a positive-order operator perturbation and of cubic-wave data (source 02's $\Delta$).\\
$\kappa_E$, $\kappa_B$; $\lambda_{\mathrm L}$, $\mu_{\mathrm L}$ & Permittivity and permeability; Lam\'e constants (source 03's $\epsilon_{\mathrm{el}},\mu$ and $\lambda,\mu$).\\
$\operatorname{in}(F)$ & Leading normalized observable $t^{\vH(F)}F_{\vH(F)}$.\\
\bottomrule
\end{longtable}

\section{Lorentzian sign conventions}\label{vtf:app:lorentzconv}

The conventions of Part~III (source 02).

The following table collects signs that must be changed together when importing a different convention.
\begin{center}\small
\begin{tabularx}{\textwidth}{@{}p{0.36\textwidth}X@{}}
\toprule
Object & Convention in Part~III\\
\midrule
Metric and wave operator & $\eta=\diag(1,-1,-1,-1)$; $\Box=\partial_\tau^2-\Delta$\\
Alternating tensor & $\epsilon_{0123}=+1$, $\epsilon^{0123}=-1$\\
Potential components & $A_\mu=(\phi,-\mathbf A)$; $F=\dd A$\\
Electric and magnetic entries & $F_{0i}=E_i$, $F_{ij}=-\epsilon_{ijk}B_k$\\
Hodge action on fields & $(\mathbf E,\mathbf B)\mapsto(-\mathbf B,\mathbf E)$\\
Hodge square by degree $0,1,2,3,4$ & $-1,+1,-1,+1,-1$\\
Inhomogeneous Maxwell equation & $\partial_\mu F^{\mu\nu}=J^\nu$; $\dd* F=-* J^\flat$\\
Stress divergence & $\partial_\mu T^{\mu\nu}=-F^{\nu\lambda}J_\lambda$\\
Wave kernel & $G_{\mathrm{ret}}=\delta(\tau-r)/(4\pi r)$ as a distribution\\
Small Hahn monomial & $t^\gamma\mapsto\omega^{-\gamma}$, $\gamma>0$\\
Spacetime differentiation & $\partial_\mu c=0$ for all constant $c\in K$\\
\bottomrule
\end{tabularx}
\end{center}

Three small diagnostics prevent large errors. A pure electric field in the $x^1$ direction has $F\wedge* F=-E_1^2\vol$; a null plane wave has a nonzero positive $T^{00}$ even though $I_1=I_2=0$; and a source-free field has $\partial_\tau\mathcal E+\Div\mathbf S=0$. These check, respectively, the signature/Hodge pairing, the distinction between invariants and energy, and the source sign in the stress balance.

\section{Sources, provenance and repository audit}\label{vtf:sec:audit}

\subsection{The three manuscripts}

\begin{center}\small
\begin{tabular}{P{0.07\textwidth}P{0.28\textwidth}P{0.16\textwidth}P{0.33\textwidth}}
\toprule
Source & Delivered as & Checks & Placed in this report\\
\midrule
01 & $n$-dimensional article, 31 pages; README, Makefile, \path{verify.py},
recorded output & 7 groups, $5{,}533$ exact assertions & Base text of Part~I,
the synthesis, the ledger and this audit's first part\\
02 & Minkowski article, 37 pages; README, source audit, Makefile, requirements,
\path{check_identities.py}, log and JSON record & 42 named exact checks &
Part~III; \cref{vtf:sec:boundaries}, \cref{vtf:3d:thm:cohomology} (with 03),
remarks in Part~I, and \cref{vtf:app:lorentzconv}\\
03 & Three-dimensional article, 33 pages; README, \path{verify_examples.py},
recorded output & 11 check groups & Part~II;
\cref{vtf:subsec:hahnspaces}, \cref{vtf:3d:prop:darboux,vtf:3d:thm:cohomology,vtf:3d:thm:regularhodge,vtf:3d:thm:quadratic,vtf:3d:prop:scale}
and remarks in Part~I\\
\bottomrule
\end{tabular}
\end{center}
All three state the repository pin \texttt{\repoSHA} and the date September 22,
2026. The manuscripts themselves are not shipped; their code and recorded data
are, byte-identical, under \path{code/} and \path{data/}, and source 02's audit
is \path{02-minkowski-SOURCE_AUDIT.md}. The page counts are those stated in the
delivered READMEs. The report was merged in the collection after placement
commit \texttt{5fe7f8d}.

\subsection{Statements printed once}\label{vtf:subsec:printedonce}

The following results are proved by more than one source and printed once; the
sources are named at the place of printing.
\begin{itemize}
\item The workspace $K=\R((t^\Gamma))$, its real closedness, the localization of
set-sized data, valuation, $\OO_K$, $\mm_K$ and standard part; strong
summability, the Neumann support lemma, and the examples showing that strong
sums are not limits: all three sources, \cref{vtf:sec:scalars}.
\item Finite tensor algebra over $K$ and the invariance of component valuations
under bounded frames: all three sources, \cref{vtf:sec:algebra}; source 02's
bounded-Lorentz version is kept as \cref{vtf:mink:prop:frame} because it adds the
equivariance of standard part and the infinite-boost counterexample.
\item The common-domain differential algebra, the failure of the sheaf property
for global supports, and the pointwise-invertible field $x^2+\eps^2$: all three
sources or sources 01 and 03, \cref{vtf:sec:sheaf}.
\item Taylor--Hahn prolongation and the flat-function remark: all three sources,
\cref{vtf:thm:prolongation}.
\item Lie brackets, forms, Cartan calculus, Lie derivatives and covariant
derivatives of tensors: all three sources, \cref{vtf:sec:fields,vtf:sec:metric}.
\item Inverse, density, signature, Levi--Civita connection, curvature and
Bianchi identities of regular metrics: sources 01 and 03, \cref{vtf:sec:metric,vtf:sec:curvature}.
\item The Hodge star, its square and the codifferential: all three sources,
\cref{vtf:sec:hodge}; source 02's four-dimensional formulas are kept in
\cref{vtf:mink:sec:tensors} as the case $n=4$, $q=3$.
\item Coefficientwise integration and Stokes, the radial Poincar\'e lemma and
polynomial-simplex integration: all three sources or sources 01 and 03,
\cref{vtf:sec:integration}.
\item The cohomology of the Hahn-extended complex: sources 02 and 03,
\cref{vtf:3d:thm:cohomology}; compact base change with products is source 01's
\cref{vtf:thm:derham}. The Hodge decomposition for a real metric: sources 01
and 03, \cref{vtf:cor:hodge}.
\item Standard-part reduction: all three sources, \cref{vtf:thm:reduction} and
the remark after it.
\item The infinite-frequency obstruction: sources 02 and 03,
\cref{vtf:mink:thm:highfrequency}; distributional coefficients: sources 02 and
03, \cref{vtf:subsec:distributions}.
\item The distinction between spatial and scalar derivations, the Euler scale
derivation and the curvature of a scale connection: all three sources or sources
02 and 03, \cref{vtf:sec:BM}.
\item The radial flux computation is \cref{vtf:subsec:radial} (source 01); its
case $n=3$ is kept in \cref{vtf:3d:subsec:monopole} for the vector-potential
obstruction.
\end{itemize}
Kept twice, because the hypotheses or proofs differ: nonlinear lifting
(\cref{vtf:thm:lifting}, source 01) and quadratic lifting by trees
(\cref{vtf:3d:thm:quadratic}, source 03); the Poynting balance with material
constants (source 03, \cref{vtf:3d:sec:physics}) and the covariant stress balance
in vacuum units (source 02, \cref{vtf:mink:prop:stressdiv}); the general
variational calculus with a scalar action (source 01) and with a polynomial
first-order Lagrangian (source 02).

\subsection{Repository audits of the sources}

\paragraph{Source 01.}

The repository was inspected read-only at commit
\begin{center}\small\texttt{\repoSHA}.\end{center}
The targeted review used the root \path{README.md}, \path{docs/README.md},
the opening scope of
\path{docs/foundations-and-computation/foundations/article.tex}, and the
normal-form, workspace, function-class, derivative and support conventions in
\path{docs/surcomplex/analysis/article.tex}. The documentation catalogue was
used to identify the nearby contours, differential-equations, spectral and
physics reports. This was not a line-by-line audit of every manuscript or a
local rebuild of the Lean project. The catalogue's report count is not used
as a claim that the entire repository, including newly added material, was
exhaustively surveyed.

The points carried forward are the use of increasing Hahn exponents,
set-sized divisible workspaces, common domains, strong summability, the
separation of scalar and coordinate derivations, and explicit statements of
formalization scope. Source 01's tensor and PDE developments are
proved independently above. The repository itself describes its manuscripts
as research drafts with theorem-specific, rather than blanket, Lean coverage.

The external literature check included the original Berarducci--Mantova
paper, the general-base-field calculus of Bertram--Gl\"ockner--Neeb, the
van den Dries--Ehrlich paper, Poonen's discussion of Hahn units, and classical
smooth and de Rham geometry sources. Two sources relevant to the cutoff date
are Poonen's January 2026 revision on units and Mantova's January 2026
Taylor-approximation preprint. Their roles are identified in the bibliography;
neither is represented as a theorem about all tensor fields on all of
$\No^n$.

\paragraph{Source 02.} The repository was inspected at the same pin through a
connected GitHub reader: the recursive tree and a large excerpt of the root
README (both responses truncated), \path{docs/README.md} lines 1--160, the
physics report's README lines 1--170, \path{docs/NORMAL_FORM_BRIDGE.md} lines
1--95, directory listings, and a repository search for ``Minkowski tensor'',
which returned no matches and is not evidence that every topic is absent. The
full physics article, the other manuscripts and the Lean sources were not read
line by line, and no Lean build was performed. The external sources checked were
the Mantova--Matusinski survey, Berarducci--Mantova, B\"ar and
B\"ar--Ginoux--Pf\"affle, Tong's notes (whose PDF pages were viewed to fix
conventions and the sign of the retarded kernel), and Madar\'asz--Stannett--Sz\'ekely;
Neumann's paper was identified bibliographically, and Conway and Gonshor are cited
as foundational monographs cross-checked against the survey. The delivered audit
is shipped as \path{02-minkowski-SOURCE_AUDIT.md}.

\paragraph{Source 03.} The repository was inspected at the same pin: the root
README, the documentation catalogue and the contours-and-stokes README. The
spectral, physics and surquaternion reports are cited as described in the
catalogue. External sources are Gonshor, Berarducci--Mantova (both papers),
Costin--Ehrlich, Lee and Warner.

\paragraph{Checked for this merge.} Every repository statement of the three
sources was checked against the tree at the pin and against the current tree.
The cited paths exist at both:
\path{docs/foundations-and-computation/foundations/article.tex},
\path{docs/surcomplex/analysis/article.tex},
\path{docs/surcomplex/contours-and-stokes},
\path{docs/surcomplex/spectral-theory},
\path{docs/physics/surreal-scalars-and-spacetime},
\path{docs/surquaternions/surquaternions} and \path{docs/NORMAL_FORM_BRIDGE.md}.
The foundations report proves that small subsets and small-index nets are
discrete (\path{found:thm:discrete}). The contours-and-stokes report contains
the coefficientwise Stokes theorem, the support-preserving Poincar\'e lemma,
polynomial-chain Stokes and the finite-partition obstruction, as source 03
states. The physics README distinguishes scalar extension from regularization
and records bounded-frame valuation invariance, as source 02 states. The
opening assessment of \path{docs/NORMAL_FORM_BRIDGE.md}, which records the
infinite normal-form ordered-field bridge, is the same at the pin and now. The
root README's statement that the Hahn layer proves the full Neumann support
lemma is present at the pin, supporting source 03's remark.

\subsection{Stale or incomplete statements corrected}\label{vtf:subsec:stale}

\begin{enumerate}
\item Source 01 described ``the repository's current documentation''. The
statement is now pinned: it held at \texttt{d22a5b3} and still holds.
\item Source 01 developed Hahn Stokes, the Poincar\'e lemma, de Rham base change
and polynomial Stokes without naming the report
\path{docs/surcomplex/contours-and-stokes}, which at the pin already contained
their complex-coefficient versions. The relation is now stated
(\cref{vtf:subsec:neighbours}), and those results are not claimed as new.
\item Source 01's remark that a Hodge decomposition for any metric other than a
real one ``requires a separate argument'' is updated: source 03 supplies the
argument for regular Riemannian Hahn metrics (\cref{vtf:3d:thm:regularhodge});
the rest stays open (\cref{vtf:subsec:open}).
\item Source 02's future direction on a variable Lorentzian metric with
nondegenerate real body is partly answered inside this report by source 01, and
source 02's direction on a locally coherent sheaf theory partly by source 01's
sheaf $\A^{\mathrm{loc}}_\Gamma$; both are re-scoped in \cref{vtf:subsec:open}.
\item Source 02's audit recorded that a repository search for ``Minkowski
tensor'' returned no matches. That held at the pin; the collection now contains
this report. The audit file is shipped as delivered.
\item Source 03's remark that parametrizing $\SO(3,K)$ by large angles needs a
separate trigonometric convention now points to the report
\path{docs/surreal/euclidean-three-space}, placed with this one, and to the spin
description in the surquaternion report.
\item The three sources name their programs by their delivered paths
(\path{code/verify.py}, \path{verification/check_identities.py},
\path{verify_examples.py}) and their records correspondingly. The prefixed names
under which they are shipped are used throughout
(\cref{vtf:subsec:verification}). The delivered Makefiles of sources 01 and 02
are shipped under \path{code/} and still refer to the delivered layout.
\end{enumerate}

\subsection{Non-claims, source by source}\label{vtf:subsec:nonclaims}

Every limitation stated by a source is kept in the text; this ledger collects
them.

\paragraph{Source 01 (18).}
(1) No completed Lean formalization and no Lean build; no repository file
modified. (2) No priority claim for every result, no resolution of a named open
problem, no exhaustive novelty search. (3) The finite checks do not prove the
infinite-support and existence theorems, and typesetting has no bearing on
correctness. (4) The framework represents exact hierarchies of scales; it is not
an automatic regularization of singularities or a transfer of all real analytic
existence theorems. (5) A bare map into $\No^n$ is not a smooth field, and
pointwise nondegeneracy does not certify an admissible inverse. (6) The
$K$-valued norm carries no completeness theorem. (7) Localization of scalar data
localizes neither regularity claims nor closure under scalar derivations.
(8) The intrinsic topology of a workspace is not the topology inherited from the
full class, and prolongation is not an identity theorem for analytic functions
nor unique outside the chosen class. (9) An arbitrary element of $K[[Z]]$ is not
promised to evaluate on the whole halo; a formal inverse function theorem is not
a convergence theorem. (10) The algebraic integration is not a countably
additive Lebesgue theory on subsets of $K^n$. (11) The positivity statement
concerns smooth common-support densities, not arbitrary Hahn-valued measures.
(12) The Hodge corollary is not an unrestricted Hodge theorem for every surreal
metric, every intrinsic $K$-manifold or a completed normed space. (13) The Hahn
cohomology is not singular cohomology of the halo. (14) The parametric sphere is
not a fine-continuous map, and the singularity at $r=0$ is not assigned a value.
(15) Flows: neither completeness of arbitrary real fields nor all-real-time flows
for infinite leading coefficients. (16) Lifting does not establish global
Einstein evolution, Navier--Stokes regularity or singularity removal; gauge
kernels and compatibility conditions must be handled first. (17) Variational and
gauge equations come without existence theorems. (18) The repository review was
targeted, not a line-by-line audit, and the mixed scalar-derivation connection
is not forced to be a physical direction.

\paragraph{Source 02 (16).}
(1) A mathematical framework, not a new experimentally validated theory; it does
not select a physical exponent group, units or detector response. (2) No complete
formal verification, no Lean build, and no claim that its theorems have checked
Lean proofs. (3) No exhaustive novelty or priority claim; classical identities are
not new because the coefficient field changed. (4) The 42 checks do not
establish Neumann summability, real closedness, transfinite recursion, PDE
existence, halo functoriality or a Lean theorem. (5) Cauchy lifting gives no
uniform Sobolev bound across exponents, no convergence at a positive real scale,
and no fine-topology well-posedness. (6) Retarded perturbation: causal support is
relative to the real body; no statement about the characteristic cone of a
$K$-valued principal symbol, and no convergence at small real scales. (7) Cubic
wave: no convergence at a small real parameter, no control of singularities of
other real solutions outside the background domain, and no claim that every real
datum has a global background. (8) Infinite boosts are algebraic facts, not
statements about attainable observers. (9) An unlimited field value is not a
measurable infinite energy, and an infinitesimal correction has no real readout
without a specified amplification. (10) No claim about black-hole interiors or
quantum gravity; scalar extension does not define the Coulomb field at the
origin; products of distributions remain undefined. (11) The infinite-frequency
theorem is a limitation of the chosen algebra, not an impossibility of every
oscillatory extension. (12) Conservation of $q_\eta(U)$ under the Lorentz force
does not give existence of solutions. (13) Stationarity only: no minimizers from
order completeness, no measure on $K^4$. (14) The fluid tensor carries no
equation of state, regularity, shock theory or well-posedness. (15) The
halo prolongation is a specified extension, not a uniqueness theorem for
functions on $K^4$ nor smoothness in the full fine topology. (16) The source
audit was targeted, and its search for ``Minkowski tensor'' is not evidence of
absence.

\paragraph{Source 03 (17).}
(1) No unrestricted integration theory on $\No^3$ and no theorem about all
functions on $\No^3$; replacing full-class language by a workspace does not
preserve an unmentioned topology or function class. (2) No physical singularity
resolution and no new experimentally validated field theory. (3) No Lean source,
no Lean-checked proof, and a formalization section that is an architecture
proposal; the repository was not modified. (4) No priority claim; Hahn
summability, coefficientwise Stokes, the radial Poincar\'e operator and the
real-closed spectral theorem are not presented as new. (5) The finite checks do
not prove support lemmas, arbitrary-rank summability, the operator theorems, the
imported Hodge theorem or the infinite-dimensional consequences. (6) The
stationary-flow corollary is periodic, stationary and restricted to infinitesimal
data around zero; nothing is asserted about time-dependent three-dimensional
Navier--Stokes regularity. (7) Quadratic lifting is not a claim that a real
nonlinear problem is solvable for arbitrary real forcing. (8) The regular-metric
Hodge theorem does not cover metrics with singular standard part; pointwise
positive definiteness is no substitute. (9) Reduction is a consistency result,
not a converse existence theorem. (10) The spectral theorem is pointwise; no
globally smooth eigenvector fields across repeated eigenvalues. (11) Taylor
prolongation is not a canonical extension and does not solve the extension
problems of Costin and Ehrlich. (12) The rotation groups are algebraic; no
parametrization by large angles is provided. (13) The fast-time obstruction is a
limitation of the chosen calculus, not a contradiction with real oscillators or
formal time series. (14) Algebraic positivity of the elastic energy gives no
analytic coercivity. (15) Products of distributions remain undefined, and the jet
distribution is not an ordinary distribution at a nonreal point. (16) A Lie
derivative does not give a global flow. (17) Its open directions are proposals,
not statements that the problems are unsolved in every formalism.

\section{A proof-dependency summary}\label{vtf:app:dependency}

The support-sensitive claims have short, explicit dependency chains.
The field and size setup uses normal forms and the classical Hahn-field
closure theorem. Differential algebra uses the finite-pair support lemma.
Prolongation adds positive-support substitution and a valuation remainder
estimate. Metrics add only a certified geometric inverse and a certified
square root. Curvature identities are finite differential-algebra identities.
Stokes and Poincar\'e use fixed real coefficient operators. Compact de Rham
base change uses finite real cohomology and support-preserving linear
splittings; the Hodge refinement imports the real Hodge theorem. Nonlinear
lifting uses a specified real linear inverse plus positive-support recursion.
No circular reliance on the desired surreal existence theorem occurs in
these constructions.

Parts~II and~III add the following chains. The finite vector identities use sums
of squares in an ordered field; the vector differential identities are finite
differential-algebra identities extended by \cref{vtf:3d:thm:closure}. The
Helmholtz decomposition, the elliptic example and the stationary-flow corollary
use real Fourier multipliers on the torus, the Leray projection, and
\cref{vtf:3d:thm:inverse,vtf:3d:thm:quadratic}. The Lorentz algebra uses real
closedness and sums of squares; the bounded reduction uses the infinitesimal
exponential and logarithm \eqref{vtf:mink:eq:formalfunctions}; the energy scale
uses positivity of sums of real squares. Cauchy and retarded lifting import the
real Cauchy and Green theorems \cite{Bar,BGP}; the cubic-wave theorem adds
positive-support recursion. The regular-metric Hodge theorem imports the real
Hodge theorem and uses the positive-support inverse.

\begin{thebibliography}{99}
\small\raggedright
\bibitem{Repo}
Vladimir Reshetnikov, \emph{Surreal}, repository documentation, inspected by all
three sources at commit \texttt{\repoSHA}, September 22, 2026.
\url{https://github.com/VladimirReshetnikov/Surreal}.
The root README and \path{docs/README.md} give the implementation and manuscript
scope; these are project sources, not a blanket verification claim.

\bibitem{RepoFound}
\emph{Foundations for Surreal and Surcomplex Mathematics}, Surreal project,
\path{docs/foundations-and-computation/foundations/article.tex}, same pinned
revision as \cite{Repo}. Targeted use: sizes, workspaces and the fine-topology
warning.

\bibitem{RepoAnalysis}
\emph{Surcomplex Analysis}, Surreal project,
\path{docs/surcomplex/analysis/article.tex}, same pinned revision as
\cite{Repo}. Targeted use: monomials, coherent coefficient fields,
summability and the distinction between derivative operators.

\bibitem{RepoContours}
\emph{Contours and Stokes: Jordan separation, Cauchy theory, and Stokes calculus
on the surcomplex plane}, Surreal project,
\path{docs/surcomplex/contours-and-stokes}, same pinned revision as \cite{Repo}.

\bibitem{RepoSpectral}
\emph{Spectral theory}, Surreal project, \path{docs/surcomplex/spectral-theory},
same pinned revision as \cite{Repo}.

\bibitem{RepoPhysics}
\emph{Surreal Scalars in Theoretical Physics: Black-Hole Singularities,
Asymptotic Structure, and the Limits of Replacing the Scalar Field}, Surreal
project, \path{docs/physics/surreal-scalars-and-spacetime}, same pinned revision
as \cite{Repo}.

\bibitem{RepoQuaternion}
\emph{Surquaternions}, Surreal project, \path{docs/surquaternions/surquaternions},
same pinned revision as \cite{Repo}.

\bibitem{RepoBridge}
\emph{Next steps for the infinite normal-form bridge}, implementation note,
Surreal project, \path{docs/NORMAL_FORM_BRIDGE.md}, assessment dated September
22, 2026; same pinned revision as \cite{Repo}.

\bibitem{Conway}
J. H. Conway, \emph{On Numbers and Games}, Academic Press, 1976; second edition,
A K Peters, 2001.

\bibitem{Gonshor}
Harry Gonshor, \emph{An Introduction to the Theory of Surreal Numbers},
London Mathematical Society Lecture Note Series 110, Cambridge University
Press, 1986.
\url{https://www.cambridge.org/core/books/an-introduction-to-the-theory-of-surreal-numbers/312AE504A3E88E804054BFB390446374}.

\bibitem{vdDE}
Lou van den Dries and Philip Ehrlich, ``Fields of surreal numbers and
exponentiation,'' \emph{Fundamenta Mathematicae} 167 (2001), 173--188.
DOI: \href{https://doi.org/10.4064/fm167-2-3}{10.4064/fm167-2-3}.
See also their published erratum, \emph{Fundamenta Mathematicae} 168 (2001).
Used as background, not as an unrestricted transfer theorem for this calculus.

\bibitem{Poonen93}
Bjorn Poonen, ``Maximally complete fields,'' \emph{L'Enseignement
Math\'ematique} (2) 39 (1993), 87--106.
\url{https://math.mit.edu/~poonen/papers/amsval.pdf}.
Used for the algebraic-closure theorem for Hahn fields with algebraically
closed residue field and divisible value group.

\bibitem{Neumann}
B. H. Neumann, ``On ordered division rings,'' \emph{Transactions of the
American Mathematical Society} 66 (1949), 202--252.
DOI: \href{https://doi.org/10.1090/S0002-9947-1949-0032593-5}{10.1090/S0002-9947-1949-0032593-5}.
Classical source for the well-ordered support and positive-word finiteness
arguments.

\bibitem{Poonen26}
Bjorn Poonen, ``Units in Hahn--Mal'cev--Neumann rings,'' author manuscript,
January 14, 2026.
\url{https://math.mit.edu/~poonen/papers/malcev.pdf}.
Theorems 1.1 and 4.1 give the positive-support unit theorem and its explicit
geometric-series form.

\bibitem{MM}
V. Mantova and M. Matusinski, ``Surreal numbers with derivation, Hardy fields
and transseries: a survey,'' in \emph{Ordered Algebraic Structures and Related
Topics}, Contemporary Mathematics 697, American Mathematical Society, 2017,
265--290. Preprint: \url{https://arxiv.org/abs/1608.03413}, version 2 (2016).

\bibitem{BM}
Alessandro Berarducci and Vincenzo Mantova, ``Surreal numbers, derivations
and transseries,'' \emph{Journal of the European Mathematical Society} 20
(2018), 339--390. DOI: \href{https://doi.org/10.4171/JEMS/769}{10.4171/JEMS/769}.
\url{https://arxiv.org/abs/1503.00315}.

\bibitem{BMGerms}
Alessandro Berarducci and Vincenzo Mantova, ``Transseries as germs of surreal
functions,'' arXiv:1703.01995 (2017).
\url{https://arxiv.org/abs/1703.01995}.
Constructs composition on the omega-series and distinguishes it from the
situation on all surreals; not a claim about arbitrary surreal-coordinate tensor
fields.

\bibitem{Mantova26}
Vincenzo Mantova, ``Monotonicity and a Taylor approximation theorem for
transseries,'' preprint, January 12, 2026, arXiv:2601.07747.
\url{https://arxiv.org/abs/2601.07747}.
Recent context for omega-series composition and its Taylor theory; not used
in the proof of the smooth finite-halo prolongation theorem.

\bibitem{CE}
O. Costin and P. Ehrlich, ``Integration on the Surreals,'' arXiv:2208.14331
(2022). \url{https://arxiv.org/abs/2208.14331}.
Cited for its specified function-extension and integration program and its
limitations, not as a premise that all smooth functions have a global canonical
extension.

\bibitem{BGN}
Wolfgang Bertram, Helge Gl\"ockner and Karl-Hermann Neeb,
``Differential calculus, manifolds and Lie groups over arbitrary infinite
fields,'' preprint, 2003, arXiv:math/0303300.
\url{https://arxiv.org/abs/math/0303300}.

\bibitem{MSS}
J. X. Madar\'asz, M. Stannett and G. Sz\'ekely, ``Groups of worldview
transformations implied by Einstein's special principle of relativity over
arbitrary ordered fields,'' \emph{The Review of Symbolic Logic} 15 (2022),
334--361. DOI:
\href{https://doi.org/10.1017/S1755020321000149}{10.1017/S1755020321000149}.

\bibitem{Lee}
John M. Lee, \emph{Introduction to Smooth Manifolds}, second edition,
Graduate Texts in Mathematics 218, Springer, 2012 (copyright 2013).
DOI: \href{https://doi.org/10.1007/978-1-4419-9982-5}{10.1007/978-1-4419-9982-5}.

\bibitem{Petersen}
Peter Petersen, \emph{Manifolds, Transversality, and de Rham Cohomology},
author lecture notes, accessed September 22, 2026, especially Chapters 6--7.
\url{https://www.math.ucla.edu/~petersen/manifolds.pdf}.

\bibitem{Warner}
Frank W. Warner, \emph{Foundations of Differentiable Manifolds and Lie
Groups}, Graduate Texts in Mathematics 94, Springer, 1983, Chapter 6.
DOI: \href{https://doi.org/10.1007/978-1-4757-1799-0}{10.1007/978-1-4757-1799-0}.
The real compact-manifold Hodge theorem is the imported analytic input to
\cref{vtf:cor:hodge,vtf:3d:thm:regularhodge}.

\bibitem{Bar}
C. B\"ar, ``Linear wave equations on Lorentzian manifolds,'' 2010.
\url{https://arxiv.org/abs/1006.2354}.
The real Cauchy and finite-propagation theorems are used with their stated
hypotheses.

\bibitem{BGP}
C. B\"ar, N. Ginoux and F. Pf\"affle, \emph{Wave Equations on Lorentzian
Manifolds and Quantization}, ESI Lectures in Mathematics and Physics, European
Mathematical Society, 2007. DOI: \href{https://doi.org/10.4171/037}{10.4171/037}.
\url{https://arxiv.org/abs/0806.1036}.

\bibitem{TongEM}
D. Tong, \emph{Electromagnetism}, lecture notes, chapter ``Electromagnetism and
Relativity,'' accessed September 22, 2026.
\url{https://davidtong.org/pdfs/teaching/electromagnetism/electro4.pdf}.

\bibitem{TongRad}
D. Tong, \emph{Electromagnetism}, lecture notes, chapter ``Electromagnetic
Radiation,'' accessed September 22, 2026.
\url{https://davidtong.org/pdfs/teaching/electromagnetism/electro5.pdf}.
\end{thebibliography}
\end{document}
